% Effet photoélectrique et effet Compton — Physique Tle C — Cours signé OnBuch+
% Programme officiel MINESEC. Compiler avec : tectonic P18-effet-photoelectrique-effet-compton.tex
\newcommand{\DOCMATIERE}{Physique}
\newcommand{\DOCNIVEAU}{Tle C}
\newcommand{\DOCMODULE}{Module 4 : Onde, matière et transformations nucléaires}
\newcommand{\DOCLECON}{P18}
\newcommand{\DOCTITRE}{Effet photoélectrique et effet Compton}
% OnBuch+ — préambule « cours signés » ENRICHI (compilé avec tectonic / XeLaTeX)
% Système visuel « Pop » d'OnBuch+ : fond crème (jamais blanc), bordures encre
% épaisses, ombres « dures » décalées, titres Archivo Black en capitales.
% Version MATHÉMATIQUES : graphes (pgfplots/tikz), boîtes pédagogiques
% (définition, propriété, méthode, attention, le savais-tu, exercice résolu,
% exercice, TP, bilan), page de garde et signature OnBuch+ ancrée en bas.
%
% Macros attendues dans main.tex AVANT \input{preamble} :
%   \DOCMATIERE  \DOCNIVEAU  \DOCMODULE  \DOCLECON (numéro)  \DOCTITRE
\documentclass[11pt]{article}

\usepackage{fontspec}
\usepackage[french]{babel}
\usepackage[a4paper,margin=2cm,top=3.4cm,bottom=2.8cm,headheight=1.4cm,footskip=1.3cm]{geometry}
\usepackage{amsmath,amssymb,amsfonts}
\usepackage{mathtools} % \xmapsto, \coloneqq... souvent utilisés par les modèles
\usepackage{cancel} % simplifications barrées (bilans de forces)
% \abs{z} (module, valeur absolue) et \norm{u} (norme), utilisés par les modèles
\providecommand{\abs}{}\let\abs\relax\DeclarePairedDelimiter{\abs}{\lvert}{\rvert}
\providecommand{\norm}{}\let\norm\relax\DeclarePairedDelimiter{\norm}{\lVert}{\rVert}
\usepackage[table]{xcolor} % [table] : fournit aussi \rowcolors (alternance de lignes)
\usepackage{pagecolor}
\usepackage{tikz}
\usetikzlibrary{arrows.meta,positioning,calc,shapes.geometric,shapes.misc,shapes.arrows,decorations.pathmorphing,decorations.markings,patterns,shadows,mindmap,trees,fit,backgrounds,matrix,intersections,angles,quotes}
\usepackage{pgfplots}
\usepackage[version=4]{mhchem} % équations nucléaires \ce{^{235}_{92}U}
\usepackage[european,siunitx]{circuitikz} % schémas électriques
\pgfplotsset{compat=1.18}
\usepgfplotslibrary{fillbetween} % aires entre courbes (calcul intégral)
\usepackage{enumitem}
\usepackage{fancyhdr}
% [most] charge toutes les bibliothèques tcolorbox (listings, minted, raster,
% documentation...) et gonfle inutilement la mémoire TeX sur un document long
% (~300 boîtes) : on ne charge que ce qui est réellement utilisé.
\usepackage[skins,breakable]{tcolorbox}
\usepackage{graphicx}
\usepackage{colortbl}
\usepackage{booktabs}
\usepackage{tabularx}
\usepackage{diagbox} % case d'en-tête diagonale des tableaux à double entrée
% Colonnes extensibles centrée / gauche / droite (C, L, R), souvent utilisées par les modèles
\newcolumntype{C}{>{\centering\arraybackslash}X}
\newcolumntype{L}{>{\raggedright\arraybackslash}X}
\newcolumntype{R}{>{\raggedleft\arraybackslash}X}
\usepackage{array}
\usepackage{multicol}
\usepackage{siunitx}
% parse-numbers=false : accepte aussi \SI{2,5 \times 10^{-3}}{...}
\sisetup{locale=FR,output-decimal-marker={,},group-minimum-digits=4,parse-numbers=false}
\DeclareSIUnit\ohm{\ensuremath{\symup{\Omega}}} % U+2126 absent de la police
\DeclareSIUnit\division{div} % réglages d'oscilloscope (ms/div, V/div)
\usepackage{qrcode}
% \degree : commande fréquemment utilisée par les modèles pour un angle ou une
% température en dehors de \SI{}{} (non fournie par les paquets chargés).
\providecommand{\degree}{^{\circ}}
% \qed (fin de démonstration), utilisé par les modèles sans amsthm
\providecommand{\qed}{\hfill$\square$}
% \barre : double barre des valeurs interdites dans un tableau de variations (tkz-tab)
\providecommand{\barre}{\ensuremath{\|}}
% environnement proof (amsthm non chargé) utilisé par les modèles
\ifdefined\proof\else\newenvironment{proof}[1][Démonstration]{\par\noindent\textit{#1.}\ }{\qed\par}\fi
\usepackage{lastpage}
\usepackage{hyperref}
\hypersetup{hidelinks}

% ============================================================
% Palette « Pop » OnBuch+ — mêmes hex que la classe P (plus_ui.dart)
% ============================================================
\definecolor{poporange}{HTML}{F59321}
\definecolor{popdark}{HTML}{E07A0C}
\definecolor{popcream}{HTML}{FFF6E8}
\definecolor{popink}{HTML}{1C1714}
\definecolor{poppurple}{HTML}{7A5AE0}
\definecolor{popgreen}{HTML}{2E9E6A}
\definecolor{poppink}{HTML}{E0457B}
\definecolor{popblue}{HTML}{2F7DE1}
\definecolor{popgold}{HTML}{C98A12}
\definecolor{popmuted}{HTML}{8A7F76}
\definecolor{popline}{HTML}{EADFD2}
\definecolor{popgray}{HTML}{8A7F76}
\colorlet{popgrey}{popgray}
% Alias tolérants (le modèle nomme parfois les couleurs différemment)
\colorlet{popwhite}{white}
\colorlet{popblack}{popink}
\colorlet{popred}{poppink}
\colorlet{popyellow}{popgold}
% Teintes claires pour fonds de tableaux / zones
\colorlet{popblueL}{popblue!10!white}
\colorlet{poporangeL}{poporange!14!white}
\colorlet{popgreenL}{popgreen!12!white}
\colorlet{poppurpleL}{poppurple!12!white}
\colorlet{poppinkL}{poppink!12!white}
\colorlet{popgoldL}{popgold!14!white}

\pagecolor{popcream}

% ============================================================
% Typographie
% ============================================================
\defaultfontfeatures{Ligatures=TeX}
\setmainfont{PlusJakartaSans-Medium.ttf}[Path=../../../fonts/,BoldFont=PlusJakartaSans-Bold.ttf]
% Maths en sans-serif (Fira Math) avec chiffres et lettres droites de Plus
% Jakarta, pour que les formules se fondent dans le texte.
\usepackage[math-style=ISO,bold-style=ISO]{unicode-math}
\setmathfont{FiraMath-Regular.otf}[Path=../../../fonts/]
\setmathfont{PlusJakartaSans-Medium.ttf}[Path=../../../fonts/,range={up/num,up/{Latin,latin},\mathup}]
\usepackage{icomma} % virgule décimale sans espace en mode math
% Caractères mathématiques Unicode absents des polices de texte (Plus Jakarta,
% Archivo Black) : rendus par leur équivalent mathématique, en texte comme en
% formule — sinon ils s'affichent en carré vide (ex. « Arithmétique dans ℤ »).
\usepackage{newunicodechar}
\newunicodechar{ℝ}{\ensuremath{\mathbb{R}}}
\newunicodechar{ℤ}{\ensuremath{\mathbb{Z}}}
\newunicodechar{ℂ}{\ensuremath{\mathbb{C}}}
\newunicodechar{ℕ}{\ensuremath{\mathbb{N}}}
\newunicodechar{ℚ}{\ensuremath{\mathbb{Q}}}
\newunicodechar{𝔻}{\ensuremath{\mathbb{D}}}
\newunicodechar{α}{\ensuremath{\alpha}}
\newunicodechar{β}{\ensuremath{\beta}}
\newunicodechar{γ}{\ensuremath{\gamma}}
\newunicodechar{δ}{\ensuremath{\delta}}
\newunicodechar{ε}{\ensuremath{\varepsilon}}
\newunicodechar{θ}{\ensuremath{\theta}}
\newunicodechar{λ}{\ensuremath{\lambda}}
\newunicodechar{μ}{\ensuremath{\mu}}
\newunicodechar{σ}{\ensuremath{\sigma}}
\newunicodechar{τ}{\ensuremath{\tau}}
\newunicodechar{φ}{\ensuremath{\varphi}}
\newunicodechar{ψ}{\ensuremath{\psi}}
\newunicodechar{ω}{\ensuremath{\omega}}
\newunicodechar{Σ}{\ensuremath{\Sigma}}
% majuscules produites par \MakeUppercase dans les titres (α-aminés -> Α)
\newunicodechar{Α}{\ensuremath{\alpha}}
\newunicodechar{Β}{\ensuremath{\beta}}
\newunicodechar{Γ}{\ensuremath{\Gamma}}
\newunicodechar{Δ}{\ensuremath{\Delta}}
\newunicodechar{Ε}{\ensuremath{\varepsilon}}
\newunicodechar{Θ}{\ensuremath{\Theta}}
\newunicodechar{Λ}{\ensuremath{\Lambda}}
\newunicodechar{Μ}{\ensuremath{\mu}}
\newunicodechar{Τ}{\ensuremath{\tau}}
\newunicodechar{Φ}{\ensuremath{\Phi}}
\newunicodechar{Ψ}{\ensuremath{\Psi}}
\newunicodechar{Ω}{\ensuremath{\Omega}}
\newunicodechar{↦}{\ensuremath{\mapsto}}
\newunicodechar{∈}{\ensuremath{\in}}
\newunicodechar{∉}{\ensuremath{\notin}}
\newunicodechar{∧}{\ensuremath{\wedge}}
\newunicodechar{⇔}{\ensuremath{\Leftrightarrow}}
\newunicodechar{⟺}{\ensuremath{\Longleftrightarrow}}
\newunicodechar{⇒}{\ensuremath{\Rightarrow}}
\newunicodechar{⟹}{\ensuremath{\Longrightarrow}}
\newunicodechar{∘}{\ensuremath{\circ}}
\newunicodechar{∪}{\ensuremath{\cup}}
\newunicodechar{∩}{\ensuremath{\cap}}
\newunicodechar{≡}{\ensuremath{\equiv}}
\newunicodechar{⊂}{\ensuremath{\subset}}
\newunicodechar{⊥}{\ensuremath{\perp}}
\newunicodechar{∃}{\ensuremath{\exists}}
\newunicodechar{∀}{\ensuremath{\forall}}
\newunicodechar{∛}{\ensuremath{\sqrt[3]{\,}}}
\newunicodechar{∜}{\ensuremath{\sqrt[4]{\,}}}
\newunicodechar{⃗}{\ensuremath{{}^{\to}}}
\newunicodechar{ⁿ}{\ifmmode^{n}\else\textsuperscript{\ensuremath{n}}\fi}
\newunicodechar{ˣ}{\ifmmode^{x}\else\textsuperscript{\ensuremath{x}}\fi}
\newunicodechar{ᵇ}{\ifmmode^{b}\else\textsuperscript{\ensuremath{b}}\fi}
\newunicodechar{ʳ}{\ifmmode^{r}\else\textsuperscript{\ensuremath{r}}\fi}
\newunicodechar{ᵘ}{\ifmmode^{u}\else\textsuperscript{\ensuremath{u}}\fi}
\newunicodechar{ʸ}{\ifmmode^{y}\else\textsuperscript{\ensuremath{y}}\fi}
\newunicodechar{ᵃ}{\ifmmode^{a}\else\textsuperscript{\ensuremath{a}}\fi}
\newunicodechar{ᵉ}{\ifmmode^{e}\else\textsuperscript{\ensuremath{e}}\fi}
\newunicodechar{ᵏ}{\ifmmode^{k}\else\textsuperscript{\ensuremath{k}}\fi}
\newunicodechar{ᵖ}{\ifmmode^{p}\else\textsuperscript{\ensuremath{p}}\fi}
\newunicodechar{ᴺ}{\ifmmode^{N}\else\textsuperscript{\ensuremath{N}}\fi}
\newunicodechar{ᶜ}{\ifmmode^{c}\else\textsuperscript{\ensuremath{c}}\fi}
\newunicodechar{ʰ}{\ifmmode^{h}\else\textsuperscript{\ensuremath{h}}\fi}
\newunicodechar{ᵠ}{\ifmmode^{q}\else\textsuperscript{\ensuremath{q}}\fi}
\newunicodechar{ₙ}{\ifmmode_{n}\else\textsubscript{\ensuremath{n}}\fi}
\newunicodechar{ₐ}{\ifmmode_{a}\else\textsubscript{\ensuremath{a}}\fi}
\newunicodechar{ᵢ}{\ifmmode_{i}\else\textsubscript{\ensuremath{i}}\fi}
\newunicodechar{ᵣ}{\ifmmode_{r}\else\textsubscript{\ensuremath{r}}\fi}
\newunicodechar{ₖ}{\ifmmode_{k}\else\textsubscript{\ensuremath{k}}\fi}
\newunicodechar{ᵦ}{\ifmmode_{\beta}\else\textsubscript{\ensuremath{\beta}}\fi}
\newunicodechar{ₓ}{\ifmmode_{x}\else\textsubscript{\ensuremath{x}}\fi}
\newfontfamily\popdisplay{ArchivoBlack-Regular.ttf}[Path=../../../fonts/]
\newfontfamily\poplogo{SpaceGrotesk-Bold.ttf}[Path=../../../fonts/]
\newfontfamily\popsemi{PlusJakartaSans-SemiBold.ttf}[Path=../../../fonts/]
\newfontfamily\popxbold{PlusJakartaSans-ExtraBold.ttf}[Path=../../../fonts/]
\newfontfamily\popxboldls{PlusJakartaSans-ExtraBold.ttf}[Path=../../../fonts/,LetterSpace=3.0]

\setlist[enumerate]{leftmargin=1.7em,itemsep=5pt,topsep=4pt}
\setlist[itemize]{leftmargin=1.7em,itemsep=4pt,topsep=4pt,label={\color{poporange}\textbullet}}
\setlength{\parindent}{0pt}
\setlength{\parskip}{5pt}
\renewcommand{\arraystretch}{1.25}

% pgfplots : style Pop par défaut
\pgfplotsset{
  every axis/.append style={axis line style={popink,line width=1pt},tick style={popink},
    grid style={popline,line width=0.5pt},label style={font=\small},tick label style={font=\footnotesize}},
  popcurve/.style={poporange,line width=1.8pt,no markers,smooth},
}
% Même style disponible dans un \draw TikZ simple (hors pgfplots)
\tikzset{popcurve/.style={poporange,line width=1.8pt}}
\tikzset{popgrid/.style={popline,very thin}}
\colorlet{popcurve}{poporange} % le nom du style est parfois utilisé comme couleur

% ============================================================
% Pill / squiggle
% ============================================================
\newcommand{\poppill}[3]{%
  \tikz[baseline=-0.55ex]\node[draw=popink,line width=1.2pt,rounded corners=2.6pt,%
    fill=#2,inner xsep=6.5pt,inner ysep=3pt]{%
    \popxboldls\fontsize{7}{7}\selectfont\color{#3}\MakeUppercase{#1}};%
}
\newcommand{\popsquiggle}[1]{%
  \tikz[baseline=-0.4ex]{
    \draw[#1,line width=2.6pt,line cap=round]
      plot [smooth] coordinates {(0,0.09) (0.34,0.01) (0.68,0.21) (1.02,0.01) (1.36,0.10)};
  }%
}
% Mot-clé surligné (marqueur orange sous le texte)
\newcommand{\cle}[1]{{\popxbold\color{popdark}#1}}

% ============================================================
% En-tête / pied
% ============================================================
\pagestyle{fancy}
\fancyhf{}
\renewcommand{\headrulewidth}{0pt}
\renewcommand{\footrulewidth}{0pt}
\lhead{\raisebox{-0.15ex}{\poplogo\fontsize{13}{13}\selectfont\color{popink}OnBuch\color{poporange}+}}
\rhead{\poppill{\DOCMATIERE\ \textbullet\ \DOCNIVEAU\ \textbullet\ Leçon \DOCLECON}{white}{popink}}
\lfoot{\popxbold\fontsize{7.6}{7.6}\selectfont\color{popmuted}\MakeUppercase{Cours signé OnBuch+}\ \textbullet\ {\popsemi\DOCTITRE}}
\rfoot{\tikz[baseline=-0.6ex]\node[rounded corners=8.5pt,fill=popink,minimum height=17pt,minimum width=17pt,inner xsep=5pt,inner sep=0]{\popxbold\fontsize{8}{8}\selectfont\color{popcream}\thepage\,/\,\pageref*{LastPage}};}

% ============================================================
% Titres
% ============================================================
\newcommand{\chaptitre}[2]{%
  \begin{tcolorbox}[enhanced,colback=poporange,colframe=popink,boxrule=2.6pt,arc=11pt,
      drop shadow=popink,left=15pt,right=15pt,top=12pt,bottom=11pt,before skip=3pt,after skip=18pt]
    \poppill{#2}{popink}{popcream}\par\vspace{9pt}
    {\raggedright\hyphenpenalty=10000\exhyphenpenalty=10000\popdisplay\fontsize{22}{24}\selectfont\color{popcream}\MakeUppercase{#1}\par}
  \end{tcolorbox}
}
% Section numérotée : pastille orange avec le numéro + titre Archivo
\newcounter{popsec}
\newcommand{\coursec}[1]{%
  \stepcounter{popsec}\setcounter{popsub}{0}%
  \par\vspace{16pt}\noindent
  \begin{minipage}{\linewidth}
  \tikz[baseline=-0.7ex]\node[draw=popink,line width=1.6pt,fill=poporange,rounded corners=4pt,minimum size=22pt,inner sep=0]{\popdisplay\fontsize{12}{12}\selectfont\color{popcream}\thepopsec};%
  \hspace{8pt}{\popdisplay\fontsize{15}{17}\selectfont\color{popink}\MakeUppercase{#1}}\par
  \vspace{4pt}\noindent{\color{poporange}\rule{36pt}{3.6pt}}
  \end{minipage}\par\nopagebreak\vspace{8pt}
}
\newcounter{popsub}
\newcommand{\courssub}[1]{%
  \stepcounter{popsub}%
  \par\vspace{9pt}\noindent{\popxbold\fontsize{11.5}{13}\selectfont\color{popink}{\color{poporange}\thepopsec.\thepopsub}\hspace{6pt}#1}\par\nopagebreak\vspace{4pt}
}

% ============================================================
% Boîtes pédagogiques — même recette : fond blanc, bordure épaisse colorée,
% ombre dure encre, badge plein. Titre optionnel [#1].
% ============================================================
\tcbset{popbase/.style={breakable,enhanced,colback=white,boxrule=2.3pt,arc=9pt,
  shadow={3pt}{-3pt}{0pt}{popink},left=14pt,right=14pt,top=11pt,bottom=11pt,before skip=11pt,after skip=15pt,
  fontupper=\popsemi\fontsize{10.6}{14.5}\selectfont\color{popink},
  fontlower=\popsemi\fontsize{10.6}{14.5}\selectfont\color{popink}}}
% Coche dessinée (le glyphe ✓ n'existe pas dans les polices du document)
\newcommand{\popcheck}[1]{\tikz[baseline=-0.5ex]\draw[#1,line width=1.6pt,line cap=round,line join=round](-0.13,0.0)--(-0.04,-0.09)--(0.13,0.1);}
\providecommand{\coche}{\popcheck{popgreen}}
\providecommand{\resultat}[1]{\par\smallskip\noindent\fcolorbox{popgreen}{popgreenL}{\parbox{\dimexpr\linewidth-2\fboxsep-2\fboxrule\relax}{\textbf{Résultat.} #1}}\par\smallskip}
\newcommand{\popboxhead}[3]{\poppill{#1}{#2}{white}\ifx\relax#3\relax\else\hspace{6pt}{\popxbold\fontsize{10.6}{12}\selectfont\color{popink}#3}\fi\par\vspace{7pt}}

\newtcolorbox{aretenir}[1][]{popbase,colframe=popgreen,before upper={\popboxhead{À retenir}{popgreen}{#1}}}
\newtcolorbox{exemplebox}[1][]{popbase,colframe=poporange,before upper={\popboxhead{Exemple}{poporange}{#1}}}
\newtcolorbox{definition}[1][]{popbase,colframe=popblue,before upper={\popboxhead{Définition}{popblue}{#1}}}
\newtcolorbox{propriete}[1][]{popbase,colframe=poppurple,before upper={\popboxhead{Propriété}{poppurple}{#1}}}
\newtcolorbox{methode}[1][]{popbase,colframe=popgold,before upper={\popboxhead{Méthode}{popgold}{#1}}}
\newtcolorbox{attention}[1][]{popbase,colframe=poppink,before upper={\popboxhead{Attention}{poppink}{#1}}}
\newtcolorbox{experience}[1][]{popbase,colframe=popblue,colback=popblueL,before upper={\popboxhead{Expérience}{popblue}{#1}}}
% « Le savais-tu ? » : carte violette pleine (contexte, culture, Cameroun)
\newtcolorbox{savaistu}[1][]{popbase,colback=poppurple,colframe=popink,
  fontupper=\popsemi\fontsize{10.4}{14.2}\selectfont\color{white},
  before upper={\poppill{Le savais-tu\,?}{white}{poppurple}\ifx\relax#1\relax\else\hspace{6pt}{\popxbold\color{white}#1}\fi\par\vspace{7pt}}}
% Exercice résolu : énoncé en haut, \tcblower puis la solution sur fond crème
\newcounter{popexo}\newcounter{popexr}
\newtcolorbox{exoresolu}[1][]{popbase,colframe=poporange,colbacklower=popcream,
  segmentation style={popink,line width=1.2pt,dashed},
  before upper={\stepcounter{popexr}\popboxhead{Exercice résolu \thepopexr}{poporange}{#1}},
  before lower={\poppill{Solution}{popink}{popcream}\par\vspace{6pt}}}
% Exercice (non résolu en place) — la correction est dans la partie Corrigés
\newtcolorbox{exercice}[1][]{popbase,colframe=popink,
  before upper={\stepcounter{popexo}\popboxhead{Exercice \thepopexo}{popink}{#1}}}
% Corrigé
\newtcolorbox{corrige}[1][]{popbase,colframe=popgreen,colback=popgreenL,
  before upper={\popboxhead{Corrigé}{popgreen}{#1}}}
% Objectifs / prérequis / bilan
\newtcolorbox{objectifs}{popbase,colframe=popink,colback=poporangeL,
  before upper={\popboxhead{Objectifs de la leçon}{popink}{}}}
\newtcolorbox{prerequis}{popbase,colframe=popink,colback=popblueL,
  before upper={\popboxhead{Prérequis}{popblue}{}}}
\newtcolorbox{bilan}{popbase,colframe=popink,colback=white,boxrule=2.8pt,
  before upper={\popboxhead{Fiche bilan}{poporange}{L'essentiel à connaître pour l'examen}}}
% Figure encadrée avec légende
\newtcolorbox{popfigure}[1][]{enhanced,breakable=false,colback=white,colframe=popline,boxrule=1.4pt,arc=8pt,
  left=10pt,right=10pt,top=10pt,bottom=8pt,before skip=10pt,after skip=12pt,halign=center,
  lower separated=false,halign lower=center,
  fontlower=\popsemi\fontsize{9}{11.5}\selectfont\color{popmuted},
  #1}
\newcommand{\legende}[1]{\par\vspace{6pt}{\popsemi\fontsize{9}{11.5}\selectfont\color{popmuted}#1}}

% ============================================================
% Signature OnBuch+
% ============================================================
\newcommand{\poplogoinline}[1]{{\poplogo\fontsize{#1}{#1}\selectfont\color{popink}OnBuch\color{poporange}+}}
\newcommand{\signatureonbuch}{%
  \begin{tcolorbox}[enhanced,colback=popink,colframe=popink,boxrule=2.2pt,arc=10pt,
      drop shadow=poporange,left=14pt,right=14pt,top=10pt,bottom=10pt,before skip=0pt,after skip=0pt]
    \tikz[baseline=-0.7ex]\node[circle,fill=popgreen,draw=popcream,line width=1.3pt,minimum size=22pt,inner sep=0]{\popcheck{white}};%
    \hspace{9pt}\begin{minipage}[c]{0.62\linewidth}
      {\popxbold\fontsize{12}{13}\selectfont\color{popcream}Signé\ \textbullet\ {\poplogo OnBuch\color{poporange}+}}\par\vspace{2pt}
      {\raggedright\popsemi\fontsize{8}{10.5}\selectfont\color{popline}\MakeUppercase{Équipe pédagogique OnBuch+}\\\MakeUppercase{Conforme au programme officiel MINESEC}\par}
    \end{minipage}\hfill
    {\poplogo\fontsize{22}{22}\selectfont\color{popcream}OnBuch\color{poporange}+}
  \end{tcolorbox}%
}

% ============================================================
% Page de garde
% #1 titre · #2 matière · #3 niveau · #4 module · #5 n° leçon · #6 durée ·
% #7 description / objectifs (texte ou itemize)
% ============================================================
\newcommand{\pagegarde}[7]{%
  \thispagestyle{empty}
  \newgeometry{a4paper,margin=1.7cm,top=1.6cm,bottom=1.4cm}%
  \begin{tikzpicture}[remember picture,overlay]
    \fill[poporange!18] ([xshift=-3.2cm,yshift=-3.6cm]current page.north east) circle (4.6cm);
    \fill[poppurple!14] ([xshift=2.4cm,yshift=7.6cm]current page.south west) circle (3.8cm);
  \end{tikzpicture}%
  {\poplogo\fontsize{30}{30}\selectfont\color{popink}OnBuch\color{poporange}+}\hfill
  \raisebox{0.6ex}{\poppill{Cours signé \textbullet\ Premium}{popink}{popcream}}\par
  \vspace{1pt}\popsquiggle{poppurple}\par
  \vspace{8pt}
  \begin{tcolorbox}[enhanced,colback=poporange,colframe=popink,boxrule=2.8pt,arc=13pt,
      drop shadow=popink,left=18pt,right=18pt,top=12pt,bottom=13pt,after skip=10pt]
    \poppill{#2 \textbullet\ #3}{popink}{popcream}\hspace{5pt}\poppill{#4}{white}{popink}\par\vspace{8pt}
    {\popdisplay\fontsize{14}{14}\selectfont\color{popink}LEÇON #5}\par\vspace{4pt}
    {\raggedright\hyphenpenalty=10000\exhyphenpenalty=10000\popdisplay\fontsize{25}{27}\selectfont\color{popcream}\MakeUppercase{#1}\par}
    \vspace{8pt}
    {\popxbold\fontsize{9.5}{9.5}\selectfont\color{popink}\MakeUppercase{Durée conseillée : #6}}
  \end{tcolorbox}
  \begin{tcolorbox}[enhanced,colback=white,colframe=popink,boxrule=2.2pt,arc=10pt,
      drop shadow=popink,left=16pt,right=16pt,top=10pt,bottom=10pt,after skip=8pt]
    \poppill{Dans cette leçon}{poppurple}{white}\par\vspace{5pt}
    \popsemi\fontsize{9.3}{12.6}\selectfont\color{popink}#7
  \end{tcolorbox}
  \noindent\begin{minipage}[t]{0.487\linewidth}
    \begin{tcolorbox}[enhanced,colback=popgreen,colframe=popink,boxrule=2.2pt,arc=10pt,
        drop shadow=popink,left=11pt,right=11pt,top=10pt,bottom=10pt,height=3.1cm,valign=center]
      \tcbox[colback=white,colframe=popink,boxrule=1.3pt,arc=5pt,boxsep=0pt,nobeforeafter,
        left=3pt,right=3pt,top=3pt,bottom=3pt,valign=center]{\qrcode[height=1.9cm]{https://play.google.com/store/apps/details?id=com.luvvix.onbuch&hl=fr}}%
      \hspace{8pt}\begin{minipage}[c]{0.5\linewidth}
        {\popdisplay\fontsize{11}{12.5}\selectfont\color{white}\MakeUppercase{Télécharge OnBuch}}\par\vspace{4pt}
        {\raggedright\popsemi\fontsize{8.4}{11}\selectfont\color{white}Gratuit \textbullet\ Google Play\\Tuteur IA, annales, concours, résultats\par}
      \end{minipage}
    \end{tcolorbox}
  \end{minipage}\hfill
  \begin{minipage}[t]{0.487\linewidth}
    \begin{tcolorbox}[enhanced,colback=poppurple,colframe=popink,boxrule=2.2pt,arc=10pt,
        drop shadow=popink,left=11pt,right=11pt,top=10pt,bottom=10pt,height=3.1cm,valign=center]
      \tikz[baseline=-0.7ex]\node[circle,fill=white,draw=popink,line width=1.3pt,minimum size=1.6cm,inner sep=0]{\popdisplay\fontsize{24}{24}\selectfont\color{poppurple}+};%
      \hspace{8pt}\begin{minipage}[c]{0.6\linewidth}
        {\popdisplay\fontsize{11}{12.5}\selectfont\color{white}\MakeUppercase{Passe à OnBuch+}}\par\vspace{4pt}
        {\raggedright\popsemi\fontsize{8.4}{11}\selectfont\color{white}Tous les cours signés, Tuteur IA illimité, fiches PDF — onglet OnBuch+ de l'app\par}
      \end{minipage}
    \end{tcolorbox}
  \end{minipage}
  \vspace{8pt}
  \popcontenu
  \vfill
  \signatureonbuch
  \restoregeometry
  \newpage
}

% Rangée « Ce cours contient » (page de garde)
\newcommand{\poptile}[3]{% #1 couleur #2 gros texte #3 légende
  \begin{tcolorbox}[enhanced,colback=#1,colframe=popink,boxrule=2pt,arc=9pt,shadow={2.5pt}{-2.5pt}{0pt}{popink},
      left=6pt,right=6pt,top=9pt,bottom=9pt,halign=center,valign=center,height=1.75cm,width=\linewidth,nobeforeafter]
    {\popdisplay\fontsize{12.5}{13}\selectfont\color{white}\MakeUppercase{#2}}\par\vspace{3pt}
    {\centering\popsemi\fontsize{7.8}{9.5}\selectfont\color{white}#3\par}
  \end{tcolorbox}}
\newcommand{\popcontenu}{%
  \par\noindent{\popxboldls\fontsize{7.5}{7.5}\selectfont\color{popmuted}CE COURS SIGNÉ CONTIENT}\par\vspace{7pt}
  \noindent\begin{minipage}[t]{0.235\linewidth}\poptile{popblue}{Cours}{complet, illustré, pas à pas}\end{minipage}\hfill
  \begin{minipage}[t]{0.235\linewidth}\poptile{poporange}{Méthodes}{et exercices résolus}\end{minipage}\hfill
  \begin{minipage}[t]{0.235\linewidth}\poptile{poppink}{Exercices}{type examen + corrigés}\end{minipage}\hfill
  \begin{minipage}[t]{0.235\linewidth}\poptile{popgreen}{Fiche bilan}{l'essentiel à retenir}\end{minipage}\par}

% Sommaire
\newtcolorbox{sommaire}{enhanced,colback=white,colframe=popink,boxrule=2.2pt,arc=10pt,
  drop shadow=popink,left=18pt,right=18pt,top=13pt,bottom=13pt,before skip=6pt,after skip=20pt,
  before upper={\poppill{Sommaire}{poppurple}{white}\par\vspace{9pt}\popsemi\fontsize{10.6}{14.5}\selectfont\color{popink}}}

% Signature de fin de cours — poussée tout en bas de la dernière page
\newcommand{\findecours}{%
  \par\vfill\nopagebreak
  \begin{center}\popsquiggle{poporange}\end{center}\vspace{4pt}
  \signatureonbuch
}

%% FIGFIX-BEGIN v4 — correctifs globaux des figures (docs/onbuch-generation/tools/figfix)
\usepackage{adjustbox}\usepackage{etoolbox}
% toute figure TikZ/pgfplots plus large que la ligne est réduite à la largeur disponible
\BeforeBeginEnvironment{tikzpicture}{\begin{adjustbox}{max width=\linewidth}}
\AfterEndEnvironment{tikzpicture}{\end{adjustbox}}
% légendes pgfplots lisibles par-dessus les courbes
\pgfplotsset{every axis/.append style={legend style={fill=white,fill opacity=0.92,text opacity=1,draw=black!15,inner sep=3pt}}}
% étiquettes d'arêtes (syntaxe edge["..."]) avec fond blanc
\tikzset{every edge quotes/.append style={fill=white,inner sep=1.2pt}}
%% FIGFIX-END
\begin{document}

\pagegarde{Effet photoélectrique et effet Compton}{Physique}{Tle C}{Module 4 : Onde, matière et transformations nucléaires}{P18}{5 h}{Cette leçon étudie l'interaction lumière-matière à l'échelle quantique : l'effet photoélectrique (extraction d'électrons d'un métal par la lumière) et l'effet Compton (diffusion d'un photon par un électron). Elle établit les relations d'Einstein et de Compton, exploite la caractéristique d'une cellule photoélectrique et ouvre sur les applications photovoltaïques.
\vspace{4pt}
\begin{multicols}{2}\small
\begin{itemize}
\item À la fin de cette leçon, je sais décrire l'expérience de Hertz et expliquer pourquoi elle met en échec le modèle ondulatoire classique
\item À la fin de cette leçon, je sais définir le seuil photoélectrique et le travail d'extraction W0 = hν0
\item À la fin de cette leçon, je sais appliquer la relation d'Einstein Ec max = hν - W0 pour calculer l'énergie cinétique des photoélectrons
\item À la fin de cette leçon, je sais tracer et exploiter la caractéristique I = f(U) d'une cellule photoélectrique (courant de saturation, potentiel d'arrêt)
\item À la fin de cette leçon, je sais relier le potentiel d'arrêt à l'énergie cinétique maximale par eU0 = Ec max
\item À la fin de cette leçon, je sais définir et calculer le rendement quantique d'une cellule photoélectrique
\end{itemize}\end{multicols}}

\begin{sommaire}
\begin{multicols}{2}
\begin{enumerate}
\item Mise en évidence de l'effet photoélectrique
\item Seuil photoélectrique et relation d'Einstein
\item Caractéristique I = f(U) d'une cellule photoélectrique
\item Applications de l'effet photoélectrique
\item Quantité de mouvement du photon
\item Effet Compton
\item Activité d'intégration
\item Méthodes et exercices résolus
\item Exercices
\item Corrigés des exercices
\item Fiche bilan
\end{enumerate}
\end{multicols}
\end{sommaire}

\begin{objectifs}
\begin{itemize}
\item À la fin de cette leçon, je sais décrire l'expérience de Hertz et expliquer pourquoi elle met en échec le modèle ondulatoire classique
\item À la fin de cette leçon, je sais définir le seuil photoélectrique et le travail d'extraction W0 = hν0
\item À la fin de cette leçon, je sais appliquer la relation d'Einstein Ec max = hν - W0 pour calculer l'énergie cinétique des photoélectrons
\item À la fin de cette leçon, je sais tracer et exploiter la caractéristique I = f(U) d'une cellule photoélectrique (courant de saturation, potentiel d'arrêt)
\item À la fin de cette leçon, je sais relier le potentiel d'arrêt à l'énergie cinétique maximale par eU0 = Ec max
\item À la fin de cette leçon, je sais définir et calculer le rendement quantique d'une cellule photoélectrique
\item À la fin de cette leçon, je sais exprimer la quantité de mouvement d'un photon p = h/λ
\item À la fin de cette leçon, je sais expliquer l'effet Compton et exploiter la relation λ' - λ = (h/(mc))(1 - cos θ)
\item À la fin de cette leçon, je sais citer des applications de l'effet photoélectrique (panneaux photovoltaïques, cellules, automatismes)
\end{itemize}
\end{objectifs}

\begin{prerequis}
\begin{itemize}
\item Ondes lumineuses : relation c = λν, spectre électromagnétique (Première)
\item Énergie cinétique Ec = (1/2)mv² et travail d'une force électrique (Première)
\item Conservation de l'énergie et de la quantité de mouvement dans les chocs (Première)
\item Notion de photon et relation E = hν (début d'année de Terminale)
\item Champ électrique et énergie potentielle électrique, relation W = qU
\end{itemize}
\end{prerequis}

\newpage

\coursec{Mise en évidence de l'effet photoélectrique}

\courssub{Situation d'accroche et problématique}

À Douala, de nombreux quartiers s'équipent de lampadaires solaires qui s'allument automatiquement à la tombée de la nuit : comment la simple lumière du soleil peut-elle produire du courant électrique ? Au début du XX\textsuperscript{e} siècle, l'expérience de Hertz montrait déjà qu'une plaque de zinc éclairée par les ultraviolets perdait sa charge électrique, un phénomène que la physique classique des ondes était incapable d'expliquer. C'est Einstein qui, en 1905, proposa que la lumière est constituée de grains d'énergie, les \cle{photons}. Plus tard, Compton montra que ces photons rebondissent sur les électrons comme des billes de billard. Cette leçon explore ces deux effets qui ont révolutionné notre conception de la lumière et qui sont à la base des panneaux solaires qui transforment le Cameroun.

La problématique de cette première section est la suivante : \textbf{comment l'expérience met-elle en évidence l'interaction entre la lumière et la matière, et pourquoi cette interaction oblige-t-elle à abandonner la description purement ondulatoire de la lumière ?}

\courssub{L'expérience historique de Hertz (1887)}

\textbf{Observation fortuite.} En 1887, le physicien allemand Heinrich Hertz étudie les ondes électromagnétiques qu'il vient de découvrir. Il remarque un fait étrange : une étincelle naît plus facilement entre deux sphères métalliques chargées lorsqu'elles sont éclairées par la lumière ultraviolette d'un arc électrique. Intrigué, il simplifie le dispositif pour isoler le phénomène.

\begin{experience}[Expérience de Hertz : électroscope et plaque de zinc]
\textbf{Matériel :} un électroscope à feuilles d'or, une plaque de zinc soigneusement polie fixée sur le plateau de l'électroscope, une lampe à arc (riche en rayonnement ultraviolet), une lame de verre ordinaire, une baguette d'ébonite électrisée par frottement.

\textbf{Protocole :}
\begin{enumerate}
\item On charge négativement l'électroscope (par contact avec la baguette d'ébonite) : les feuilles d'or, portant des charges de même signe, se repoussent et s'écartent.
\item On éclaire la plaque de zinc avec la lampe à arc : \textbf{les feuilles retombent rapidement} : l'électroscope se décharge.
\item On interpose une lame de verre entre la lampe et le zinc : les feuilles \textbf{restent écartées} : plus aucune décharge.
\item On remplace la lampe à arc par une lampe ordinaire très puissante (lumière visible intense) : les feuilles \textbf{restent écartées}, même après une longue exposition.
\item On charge cette fois l'électroscope \emph{positivement} et on éclaire le zinc avec la lampe à arc : les feuilles \textbf{ne retombent pas}.
\end{enumerate}

\textbf{Interprétation :}
\begin{itemize}
\item L'électroscope chargé négativement possède un excès d'électrons. Sous l'éclairage UV, ces électrons sont \emph{arrachés} au métal : la charge négative s'évacue, les feuilles retombent.
\item Le verre absorbe les ultraviolets mais laisse passer la lumière visible : la disparition de l'effet avec la lame de verre prouve que ce sont les \textbf{UV}, et non la lumière visible, qui sont responsables.
\item La lumière visible, même très intense, ne produit aucun effet sur le zinc : \textbf{l'intensité ne suffit pas}.
\item Chargé positivement, l'électroscope ne se décharge pas : les électrons arrachés sont immédiatement rappelés par la plaque positive et y retombent. Cela confirme que ce sont bien des \emph{électrons} (charges négatives) qui sont émis.
\end{itemize}
\end{experience}

\begin{popfigure}
\begin{center}
\begin{tikzpicture}[scale=0.9, >=Stealth]
% Lampe UV
\draw[fill=popgoldL, draw=popgold] (-4.5,2.2) circle (0.45);
\node[popdark] at (-4.5,3.1) {\small Lampe à arc (UV)};
\foreach \a in {-20,0,20} \draw[->, poppurple, thick] (-4.5,2.2) ++(\a-90:0.5) -- ++(\a-90:1.1);
\draw[->, poppurple, thick, decorate, decoration={snake, amplitude=0.8mm}] (-4.0,1.9) -- (-1.3,0.9);
\draw[->, poppurple, thick, decorate, decoration={snake, amplitude=0.8mm}] (-4.0,2.4) -- (-1.3,1.6);
\node[poppurple] at (-2.6,2.35) {\small $\nu \geq \nu_0$};
% Plaque de zinc
\draw[fill=popmuted!40, draw=popdark] (-1.2,0.3) rectangle (-0.8,2.3);
\node[popdark, rotate=90] at (-1.0,1.3) {\small plaque de zinc};
% Electrons arrachés
\foreach \y in {0.8,1.3,1.8} {
  \draw[->, popblue, thick] (-0.8,\y) -- (0.4,\y+0.25);
  \node[popblue] at (0.65,\y+0.28) {\small $e^-$};
}
% Tige + plateau
\draw[popdark, thick] (-1.0,0.3) -- (-1.0,-0.6);
\draw[popdark, thick] (-1.0,-0.6) -- (0.6,-0.6);
% Cloche / tige de l'électroscope
\draw[popdark, thick] (0.6,-0.6) -- (0.6,-2.2);
% Cage de l'électroscope
\draw[popdark] (-0.3,-0.9) rectangle (1.5,-3.1);
% Feuilles d'or écartées
\draw[fill=popgoldL, draw=popgold, thick] (0.6,-2.2) -- (0.1,-3.0) -- (0.35,-3.0) -- cycle;
\draw[fill=popgoldL, draw=popgold, thick] (0.6,-2.2) -- (1.1,-3.0) -- (0.85,-3.0) -- cycle;
\node[popdark] at (0.6,-3.5) {\small feuilles d'or écartées (charge $-$)};
% Charges négatives
\node[popblue] at (-0.15,-1.3) {\small $-$};
\node[popblue] at (1.35,-1.7) {\small $-$};
\node[popblue] at (-0.15,-2.5) {\small $-$};
\end{tikzpicture}
\end{center}
\legende{Expérience de Hertz : le rayonnement ultraviolet arrache des électrons à la plaque de zinc ; l'électroscope chargé négativement se décharge et ses feuilles retombent.}
\end{popfigure}

\courssub{Définition de l'effet photoélectrique}

L'expérience précédente montre qu'un métal convenablement éclairé émet des électrons. Ces électrons arrachés par la lumière sont appelés \cle{photoélectrons}.

\begin{definition}[Effet photoélectrique]
L'\cle{effet photoélectrique} est l'émission d'électrons (appelés \cle{photoélectrons}) par un métal (ou plus généralement un matériau) lorsqu'il est éclairé par une radiation électromagnétique de \textbf{fréquence suffisante}.

Lorsque l'effet se produit à la surface libre d'un métal, on parle d'effet photoélectrique \emph{externe} : les électrons quittent le métal.
\end{definition}

\begin{attention}[Ce que l'effet photoélectrique n'est pas]
Ne confonds pas :
\begin{itemize}
\item l'effet photoélectrique (la \emph{lumière} arrache des électrons) ;
\item l'émission thermoélectronique (c'est la \emph{chaleur} qui arrache les électrons, comme dans les anciens tubes à vide) ;
\item l'effet thermique ordinaire (un métal éclairé s'échauffe sans nécessairement émettre d'électrons).
\end{itemize}
Dans l'effet photoélectrique, c'est la \textbf{fréquence} de la radiation qui décide, pas la température ni l'éclairement.
\end{attention}

\courssub{Les constats expérimentaux fondamentaux}

En reprenant systématiquement l'expérience avec différents métaux et différentes radiations, les physiciens (Hallwachs, Lenard, puis Millikan) ont établi trois constats qui déroutèrent la physique de l'époque.

\begin{propriete}[Constats expérimentaux de l'effet photoélectrique — admis]
\begin{enumerate}
\item \textbf{Existence d'un seuil :} pour un métal donné, l'émission de photoélectrons n'a lieu que si la fréquence $\nu$ de la radiation vérifie
\[ \nu \geq \nu_0 \qquad \text{ou, de façon équivalente,} \qquad \lambda \leq \lambda_0 = \frac{c}{\nu_0} \]
où $\nu_0$ est la \cle{fréquence seuil} et $\lambda_0$ la \cle{longueur d'onde seuil}, \textbf{caractéristiques du métal} utilisé.
\item \textbf{Quasi-instantanéité :} lorsque $\nu \geq \nu_0$, l'émission commence dès l'éclairage, avec un retard inférieur à $\SI{1e-9}{s}$, même pour une lumière extrêmement faible.
\item \textbf{Indépendance vis-à-vis de l'intensité :} si $\nu < \nu_0$, \emph{aucune} émission n'est observée, quelle que soit l'intensité lumineuse et quelle que soit la durée d'éclairage. Si $\nu \geq \nu_0$, l'intensité ne fixe que le \emph{nombre} de photoélectrons émis par seconde, pas leur existence.
\end{enumerate}
\end{propriete}

\begin{exemplebox}[Le seuil du zinc : pourquoi la lampe ordinaire échoue]
Le zinc possède une longueur d'onde seuil $\lambda_0 \approx \SI{290}{nm}$, située dans l'ultraviolet. La fréquence seuil correspondante vaut :
\[ \nu_0 = \frac{c}{\lambda_0} = \frac{\num{3,00e8}}{\num{290e-9}} \approx \num{1,03e15}{Hz}. \]
Or le spectre visible s'étend de $\SI{400}{nm}$ (violet) à $\SI{750}{nm}$ (rouge), soit des fréquences comprises entre $\num{4,0e14}{Hz}$ et $\num{7,5e14}{Hz}$, toutes \emph{inférieures} à $\nu_0$. Une lampe ordinaire, même très puissante, ne peut donc \textbf{jamais} déclencher l'effet photoélectrique sur le zinc : c'est exactement ce qu'observe Hertz. En revanche, la lampe à arc émet des UV de longueur d'onde inférieure à $\SI{290}{nm}$, donc de fréquence supérieure à $\nu_0$ : l'effet se produit.
\end{exemplebox}

\begin{methode}[Vérifier si une radiation déclenche l'effet photoélectrique]
\begin{enumerate}
\item Identifier le métal et sa fréquence seuil $\nu_0$ (ou sa longueur d'onde seuil $\lambda_0$).
\item Calculer la fréquence de la radiation : $\nu = \dfrac{c}{\lambda}$ avec $c = \num{3,00e8}{m.s^{-1}}$.
\item Comparer : si $\nu \geq \nu_0$ (c'est-à-dire $\lambda \leq \lambda_0$), l'effet a lieu ; sinon, il n'a pas lieu, quelle que soit l'intensité.
\end{enumerate}
\end{methode}

\courssub{L'échec du modèle ondulatoire classique}

À la fin du XIX\textsuperscript{e} siècle, la lumière est décrite comme une \textbf{onde électromagnétique} (théorie de Maxwell, confirmée par Hertz lui-même !). Or cette description échoue totalement à expliquer les trois constats précédents.

\textbf{Ce que prévoit le modèle ondulatoire.} Pour une onde, l'énergie transportée est proportionnelle à l'\emph{intensité} (au carré de l'amplitude) et se répartit de façon \emph{continue} sur le front d'onde. Un électron du métal, bombardé par une onde faible, devrait donc \emph{accumuler progressivement} de l'énergie jusqu'à pouvoir s'extraire.

\textbf{Trois contradictions flagrantes :}
\begin{enumerate}
\item \textbf{Le seuil en fréquence est inexplicable.} Pour une onde, seule l'énergie reçue compte : une lumière rouge très intense devrait finir par arracher des électrons. Or l'expérience montre que $\nu < \nu_0$ ne produit \emph{jamais} d'effet, même à très forte intensité.
\item \textbf{L'instantanéité est impossible.} Un calcul classique montre qu'avec une lumière faible, un électron devrait attendre plusieurs \emph{minutes}, voire des heures, pour accumuler l'énergie d'extraction. Or l'émission débute en moins de $\SI{1e-9}{s}$.
\item \textbf{L'énergie des photoélectrons ne dépend pas de l'intensité} (nous le verrons en section 2), alors qu'une onde plus intense devrait communiquer plus d'énergie à chaque électron.
\end{enumerate}

\begin{attention}[L'erreur la plus fréquente au Bac]
Beaucoup d'élèves écrivent : « si on augmente l'intensité lumineuse, on finit par déclencher l'effet photoélectrique ». C'est \textbf{faux}. Sous le seuil ($\nu < \nu_0$), aucune intensité, aucune durée d'exposition ne déclenche l'effet. \textbf{Seule la fréquence de la radiation décide de l'existence de l'effet} ; l'intensité ne règle que le \emph{débit} de photoélectrons quand l'effet existe déjà.
\end{attention}

\courssub{L'hypothèse révolutionnaire d'Einstein (1905)}

En 1900, Max Planck avait introduit, pour expliquer le rayonnement des corps chauds, l'idée que les échanges d'énergie entre matière et lumière se font par \emph{paquets} d'énergie $h\nu$, où $h = \num{6,63e-34}{J.s}$ est la \cle{constante de Planck}. En 1905, Albert Einstein va plus loin : ce n'est pas seulement l'\emph{échange} qui est quantifié, c'est la \textbf{lumière elle-même}.

\begin{definition}[Le photon]
La lumière de fréquence $\nu$ est constituée d'un flux de particules d'énergie appelées \cle{photons}. Chaque photon transporte une énergie indivisible :
\[ E = h\nu = \frac{hc}{\lambda} \]
où $h = \num{6,63e-34}{J.s}$ (constante de Planck), $\nu$ la fréquence et $\lambda$ la longueur d'onde dans le vide. Le photon se déplace à la célérité $c$ et possède une masse nulle.
\end{definition}

\textbf{Vérification dimensionnelle :} $[h\nu] = \text{J.s} \times \text{s}^{-1} = \text{J}$ : c'est bien une énergie.

\textbf{Le mécanisme selon Einstein.} Lorsqu'un photon rencontre un électron du métal, il est absorbé \emph{intégralement} (tout ou rien) : il cède la totalité de son énergie $h\nu$ à un seul électron, en un choc quasi instantané. Trois cas se présentent :
\begin{itemize}
\item si $h\nu$ est inférieur à l'énergie minimale d'extraction, l'électron reste prisonnier du métal : \textbf{pas d'effet}, quel que soit le nombre de photons (d'où l'existence du seuil) ;
\item si $h\nu$ suffit, l'électron est éjecté \emph{immédiatement} (d'où l'instantanéité, le transfert se faisant en un seul choc) ;
\item augmenter l'intensité revient à augmenter le \emph{nombre de photons} par seconde, donc le nombre d'électrons émis, mais pas l'énergie de chacun.
\end{itemize}

Toutes les contradictions du modèle ondulatoire disparaissent d'un coup. Nous quantifierons ce bilan énergétique (relation d'Einstein) dans la section suivante.

\begin{savaistu}[Un prix Nobel pour l'effet photoélectrique]
C'est pour son explication de l'effet photoélectrique — et \emph{non} pour la relativité, jugée trop controversée à l'époque — qu'Albert Einstein reçut le prix Nobel de physique en 1921. Ironie de l'histoire : Hertz, dont l'expérience ouvrit la voie, était précisément en train de \emph{confirmer} la théorie ondulatoire de Maxwell lorsqu'il découvrit le phénomène qui allait la compléter. La lumière possède un \emph{double aspect} : ondulatoire (interférences, diffraction) et corpusculaire (effet photoélectrique, effet Compton). C'est la \cle{dualité onde-corpuscule}, pilier de la physique quantique.
\end{savaistu}

\courssub{La cellule photoélectrique : du phénomène à la mesure}

Pour étudier quantitativement l'effet photoélectrique (et l'exploiter technologiquement), on utilise une \cle{cellule photoélectrique}.

\begin{definition}[Cellule photoélectrique]
Une \cle{cellule photoélectrique} (ou cellule photoémissive) est un tube à vide contenant deux électrodes :
\begin{itemize}
\item la \cle{photocathode} $C$ : une couche de métal photoémissif (alcalin, par exemple le césium) qui émet des photoélectrons lorsqu'elle est éclairée à travers une fenêtre ;
\item l'\cle{anode} $A$ : une électrode qui collecte les photoélectrons.
\end{itemize}
Le vide poussé (pression résiduelle très faible) permet aux électrons de circuler sans chocs avec les molécules de gaz. Placée dans un circuit comportant un générateur de tension réglable $U_{AC}$, un ampèremètre et un voltmètre, la cellule est traversée par un courant $I$ dès que la photocathode est convenablement éclairée : la lumière s'y transforme directement en courant électrique.
\end{definition}

\begin{popfigure}
\begin{center}
\begin{circuitikz}[european, scale=0.95, transform shape]
% Tube à vide
\draw[popdark, thick] (2,0.2) circle (1.9);
\node[popmuted] at (2,-1.4) {\small vide};
% Photocathode (arc)
\draw[popblue, very thick] (1.1,-0.9) arc (250:290:2.6);
\node[popblue] at (0.6,-1.15) {\small $C$};
% Anode
\draw[poporange, very thick] (2.9,-0.6) -- (2.9,0.6);
\node[poporange] at (3.25,0.75) {\small $A$};
% Fenêtre et lumière
\draw[poppurple, thick, ->, decorate, decoration={snake, amplitude=0.8mm}] (-1.3,1.2) -- (0.5,0.4);
\draw[poppurple, thick, ->, decorate, decoration={snake, amplitude=0.8mm}] (-1.3,0.6) -- (0.55,0.05);
\node[poppurple] at (-1.3,1.6) {\small lumière $\nu \geq \nu_0$};
% Photoélectrons
\draw[->, popblue, dashed, thick] (1.35,0.35) -- (2.75,0.3);
\draw[->, popblue, dashed, thick] (1.3,-0.15) -- (2.75,-0.15);
\node[popblue] at (2.05,0.62) {\small $e^-$};
% Connexions vers le circuit
\draw (1.1,-0.9) -- (1.1,-2.6);
\draw (2.9,0.6) -- (2.9,-2.6);
% Circuit extérieur : source de tension variable + ampèremètre
\draw (1.1,-2.6) -- (1.1,-3.4) to[ammeter, l_={$I$}] (3.0,-3.4) to[battery1, l_={$U_{AC}$}, invert] (5.2,-3.4) -- (5.2,-2.6) -- (2.9,-2.6);
% Voltmètre en parallèle sur la source (donc sur le tube)
\draw (3.0,-3.4) -- (3.0,-4.3) to[voltmeter, l_={$U$}] (5.2,-4.3) -- (5.2,-3.4);
\end{circuitikz}
\end{center}
\legende{Cellule photoélectrique en coupe et son circuit : la lumière frappe la photocathode $C$, les photoélectrons sont collectés par l'anode $A$ ; l'ampèremètre mesure le photocourant $I$ et le voltmètre la tension $U_{AC}$ aux bornes du tube.}
\end{popfigure}

\textbf{Principe de fonctionnement.} Sans lumière, le circuit est ouvert (le vide isole les deux électrodes) : $I = 0$. Lorsque la photocathode reçoit une radiation de fréquence $\nu \geq \nu_0$, les photoélectrons émis sont attirés par l'anode (si celle-ci est portée à un potentiel positif par rapport à la cathode) et le courant $I$ circule. La cellule réalise ainsi la conversion directe \emph{lumière} $\rightarrow$ \emph{courant électrique} : c'est l'ancêtre des capteurs qui équipent les lampadaires solaires automatiques de Douala. L'étude détaillée de la caractéristique $I = f(U)$ fera l'objet de la section 3.

\begin{exoresolu}[La radiation peut-elle déclencher l'effet ?]
Énoncé. La fréquence seuil du césium vaut $\nu_0 = \num{4,60e14}{Hz}$. Une cellule photoélectrique à photocathode de césium est éclairée successivement par : (a) une lumière rouge de longueur d'onde $\lambda_1 = \SI{700}{nm}$ ; (b) une lumière verte de longueur d'onde $\lambda_2 = \SI{550}{nm}$ ; (c) un rayonnement ultraviolet de longueur d'onde $\lambda_3 = \SI{250}{nm}$. Dans quels cas observe-t-on un photocourant ? On donne $c = \num{3,00e8}{m.s^{-1}}$.
\tcblower
Solution. Calculons d'abord la longueur d'onde seuil du césium :
\[ \lambda_0 = \frac{c}{\nu_0} = \frac{\num{3,00e8}}{\num{4,60e14}} = \num{6,52e-7}{m} = \SI{652}{nm}. \]
L'effet photoélectrique se produit si et seulement si $\lambda \leq \lambda_0$.
\begin{itemize}
\item (a) $\lambda_1 = \SI{700}{nm} > \lambda_0$ : \textbf{pas d'effet}, même en augmentant l'intensité de la lumière rouge.
\item (b) $\lambda_2 = \SI{550}{nm} < \lambda_0$ : \textbf{effet photoélectrique}, un photocourant apparaît.
\item (c) $\lambda_3 = \SI{250}{nm} < \lambda_0$ : \textbf{effet photoélectrique}, avec des photoélectrons encore plus énergétiques (nous le quantifierons avec la relation d'Einstein à la section 2).
\end{itemize}
Vérification par les fréquences : $\nu_1 = \dfrac{c}{\lambda_1} = \num{4,29e14}{Hz} < \nu_0$, $\nu_2 = \num{5,45e14}{Hz} > \nu_0$, $\nu_3 = \num{1,20e15}{Hz} > \nu_0$ : conclusions identiques.
\end{exoresolu}

\begin{aretenir}[L'essentiel de la section]
\begin{itemize}
\item L'\cle{effet photoélectrique} est l'émission d'électrons par un métal éclairé par une radiation de fréquence suffisante (expérience de Hertz, 1887).
\item L'émission n'a lieu que si $\nu \geq \nu_0$ (ou $\lambda \leq \lambda_0 = c/\nu_0$), seuil \textbf{caractéristique du métal} ; elle est quasi instantanée ($< \SI{1e-9}{s}$) et indépendante de l'intensité sous le seuil.
\item Le modèle ondulatoire classique échoue : Einstein (1905) propose que la lumière est un flux de \cle{photons} d'énergie $E = h\nu$, absorbés intégralement, un par un, par les électrons.
\item La \cle{cellule photoélectrique} (photocathode + anode sous vide) convertit la lumière en courant électrique mesurable.
\end{itemize}
\end{aretenir}

\coursec{Seuil photoélectrique et relation d'Einstein}

Dans la section précédente, nous avons mis en évidence l'effet photoélectrique : une plaque de zinc chargée négativement, éclairée par la lumière ultraviolette d'une lampe à vapeur de mercure, se décharge, tandis que la lumière visible, même très intense, reste sans effet. Cette observation expérimentale est un véritable scandale pour la physique classique : si la lumière était une simple onde, une lumière intense devrait toujours finir par arracher des électrons, quelle que soit sa couleur. Or ce n'est pas ce qu'on observe. Il faut donc changer de regard : c'est Einstein qui, en 1905, a proposé l'explication décisive en supposant que la lumière est transportée par des grains d'énergie, les \cle{photons}. Cette section construit rigoureusement ce modèle et aboutit à la célèbre \cle{relation d'Einstein}, l'une des formules les plus rentables du programme de Terminale C.

\courssub{Le travail d'extraction $W_0$ : le « péage énergétique » du métal}

\textbf{Intuition.} Un électron dans un métal n'est pas libre comme un oiseau dans le ciel : il est retenu par l'attraction électrique des ions positifs du réseau cristallin. Pour le faire sortir du métal, il faut vaincre cette attraction, c'est-à-dire lui fournir une énergie minimale. Imagine un enfant au fond d'un puits : pour sortir, il lui faut au moins l'énergie correspondant à la profondeur du puits. Cette « profondeur énergétique » est le travail d'extraction.

\begin{definition}[Travail d'extraction]
Le \cle{travail d'extraction} (ou travail de sortie) d'un métal, noté $W_0$, est l'\textbf{énergie minimale} qu'il faut fournir à un électron pour l'extraire du métal (avec une vitesse quasi nulle à la sortie). Il se mesure en joules (J) ou, plus commodément, en électron-volts (eV). $W_0$ est une \textbf{caractéristique du métal} : il ne dépend ni de la lumière incidente, ni de l'intensité de l'éclairage.
\end{definition}

\textbf{Lien avec la fréquence seuil.} L'expérience de Hertz reprise en TP montre que, pour chaque métal, il existe une fréquence minimale $\nu_0$ (la \cle{fréquence seuil}) en dessous de laquelle aucun électron n'est émis, quelle que soit l'intensité lumineuse. Puisqu'un photon de fréquence $\nu_0$ apporte exactement l'énergie minimale d'extraction, on écrit :

\begin{aretenir}[Travail d'extraction et seuil photoélectrique]
\[
W_0 = h\,\nu_0 = \dfrac{h\,c}{\lambda_0}
\]
où $h = \SI{6,63e-34}{J.s}$ est la constante de Planck, $c = \SI{3,00e8}{m/s}$ la célérité de la lumière, $\nu_0$ la fréquence seuil et $\lambda_0$ la \cle{longueur d'onde seuil} du métal.
\end{aretenir}

\textbf{Vérification dimensionnelle.} $[h\nu_0] = \mathrm{J.s \times s^{-1}} = J$ : c'est bien une énergie. De même $[hc/\lambda_0] = \mathrm{J.s \times m.s^{-1} / m} = J$. La relation est homogène.

\begin{exemplebox}[Longueur d'onde seuil du zinc]
Pour le zinc, $W_0 = \SI{4,3}{eV} = 4{,}3 \times 1{,}6 \times 10^{-19} = \SI{6,88e-19}{J}$. Alors :
\[
\lambda_0 = \dfrac{hc}{W_0} = \dfrac{6{,}63 \times 10^{-34} \times 3{,}00 \times 10^{8}}{6{,}88 \times 10^{-19}} \approx \SI{2,89e-7}{m} = \SI{289}{nm}.
\]
Cette longueur d'onde est dans l'\textbf{ultraviolet} : voilà pourquoi la lumière visible ($\lambda > \SI{400}{nm}$) ne décharge jamais l'électroscope à plaque de zinc, alors que la lampe à vapeur de mercure, riche en UV, y parvient instantanément. Le mystère de la section précédente est résolu !
\end{exemplebox}

\courssub{L'électron-volt : l'unité des physiciens de l'atome}

Les énergies mises en jeu à l'échelle atomique sont de l'ordre de $10^{-19}$ J : le joule est une unité mal adaptée, comme le kilomètre pour mesurer une fourmi. On utilise donc l'\cle{électron-volt}.

\begin{definition}[Électron-volt]
Un électron-volt ($\SI{1}{eV}$) est l'énergie cinétique acquise par un électron accéléré sous une tension de $\SI{1}{V}$ :
\[
\SI{1}{eV} = e \times \SI{1}{V} = 1{,}6 \times 10^{-19}\ \mathrm{C} \times 1\ \mathrm{V} = \SI{1,6e-19}{J}.
\]
Pour convertir : $E\ (\mathrm{J}) = E\ (\mathrm{eV}) \times 1{,}6 \times 10^{-19}$ et inversement $E\ (\mathrm{eV}) = \dfrac{E\ (\mathrm{J})}{1{,}6 \times 10^{-19}}$.
\end{definition}

\begin{attention}[Conversions et constantes à connaître par cœur]
Dans les exercices, les données sont souvent mélangées : $W_0$ en eV, $\lambda$ en nm, $h$ en J.s. Avant tout calcul, \textbf{convertis tout en unités SI} (joules, mètres, hertz), ou utilise la formule pratique en eV et nm : $h\nu\ (\mathrm{eV}) = \dfrac{1240}{\lambda\ (\mathrm{nm})}$ (car $hc/e \approx \SI{1240}{eV.nm}$). Vérifie systématiquement l'homogénéité du résultat.
\end{attention}

\courssub{Bilan énergétique et relation d'Einstein}

\textbf{Le modèle du photon.} Einstein postule que la lumière de fréquence $\nu$ est constituée de photons, chacun transportant l'énergie $E = h\nu$. Dans l'effet photoélectrique, un photon est absorbé \textbf{intégralement} par \textbf{un seul} électron du métal : c'est un choc à un contre un, pas un chauffage collectif. Cette hypothèse explique d'un coup les deux faits expérimentaux déroutants : l'émission est \textbf{instantanée} (un seul photon suffit) et elle ne dépend que de la \textbf{fréquence} (l'énergie d'un photon est $h\nu$, indépendante du nombre de photons).

\begin{propriete}[Relation d'Einstein pour l'effet photoélectrique]
Lorsqu'un métal de travail d'extraction $W_0$ est éclairé par une radiation de fréquence $\nu \geq \nu_0$, l'énergie cinétique maximale des photoélectrons émis est :
\[
\boxed{\;E_{c\,\max} = h\nu - W_0 = h(\nu - \nu_0) = hc\left(\dfrac{1}{\lambda} - \dfrac{1}{\lambda_0}\right)\;}
\]
La démonstration (bilan énergétique) est \textbf{exigible au Baccalauréat}.
\end{propriete}

\begin{experience}[Démonstration guidée de la relation d'Einstein]
Raisonnons par \textbf{conservation de l'énergie} lors de l'absorption d'un photon par un électron du métal.
\begin{enumerate}
\item \textbf{Énergie reçue} : l'électron absorbe un photon et reçoit l'énergie $h\nu$.
\item \textbf{Énergie dépensée} : pour sortir du métal, l'électron doit fournir au minimum le travail d'extraction $W_0$. Les électrons les plus favorisés (ceux de la couche superficielle) ne perdent rien d'autre.
\item \textbf{Énergie restante} : le surplus est emporté sous forme d'énergie cinétique. Pour les électrons les mieux placés, cette énergie cinétique est maximale.
\item \textbf{Bilan} : énergie reçue = énergie dépensée + énergie restante, soit
\[
h\nu = W_0 + E_{c\,\max} \quad \Longrightarrow \quad E_{c\,\max} = h\nu - W_0.
\]
\item \textbf{Condition d'émission} : l'extraction n'est possible que si $h\nu \geq W_0$, c'est-à-dire $\nu \geq \nu_0$ ou de façon équivalente $\lambda \leq \lambda_0$. Si $h\nu < W_0$, le photon est absorbé sans émission : son énergie est dissipée en agitation thermique dans le métal.
\end{enumerate}
\end{experience}

\begin{popfigure}
\begin{center}
\begin{tikzpicture}[scale=1.0,>=Stealth]
\% Niveau du métal
\fill[popblueL] (-0.5,-2.6) rectangle (7.5,0);
\node[popdark] at (3.5,-1.3) {\textbf{Métal} (électrons liés)};
\% Surface
\draw[popdark,very thick] (-0.5,0) -- (7.5,0);
\node[anchor=south west,popdark] at (-0.5,0.05) {\small surface du métal};
\% Electron
\fill[popink] (1.2,-0.9) circle (0.12);
\node[below,popink] at (1.2,-1.0) {\small électron};
\% Photon incident
\draw[->,poporange,very thick,decorate,decoration={snake,amplitude=1.5pt,segment length=6pt}] (-0.4,1.6) -- (1.2,-0.75);
\node[poporange,anchor=south east] at (-0.3,1.6) {\small photon $h\nu$};
\% W0 arrow
\draw[->,poppurple,very thick] (2.6,-0.9) -- (2.6,0.35);
\node[poppurple,right] at (2.7,-0.25) {\small $W_0$ : extraction};
\% Ec arrow
\draw[->,popgreen,very thick] (2.6,0.35) -- (2.6,1.9);
\node[popgreen,right] at (2.7,1.2) {\small $E_{c\,\max}$ : énergie cinétique};
\% Electron sortant
\fill[popink] (2.6,2.15) circle (0.12);
\draw[->,popink] (2.75,2.15) -- (3.6,2.15);
\node[right,popink] at (3.6,2.15) {\small $\vec{v}_{\max}$};
\% Bilan
\node[align=center,popdark] at (5.7,1.2) {\small Bilan :\\[2pt] $h\nu = W_0 + E_{c\,\max}$};
\end{tikzpicture}
\end{center}
\legende{Diagramme énergétique de l'effet photoélectrique : l'énergie $h\nu$ du photon se partage entre le travail d'extraction $W_0$ (franchissement de la surface) et l'énergie cinétique $E_{c\,\max}$ de l'électron émis.}
\end{popfigure}

\courssub{Conséquences fondamentales : le rôle respectif de la fréquence et de l'intensité}

La relation d'Einstein tranche définitivement le débat entre la physique ondulatoire classique et l'expérience. Deux grandeurs de l'onde lumineuse jouent des rôles \textbf{totalement distincts} :

\begin{aretenir}[Fréquence et intensité : qui fait quoi ?]
\begin{itemize}
\item La \textbf{fréquence} $\nu$ de la radiation fixe l'\textbf{énergie} de chaque photoélectron : $E_{c\,\max} = h\nu - W_0$ ne dépend que de $\nu$ et du métal. Augmenter l'intensité sans changer $\nu$ ne rend \textbf{aucun} électron plus rapide.
\item L'\textbf{intensité} lumineuse fixe le \textbf{nombre de photons} par seconde, donc le \textbf{nombre de photoélectrons émis} par seconde (le courant de saturation, que nous étudierons à la section suivante).
\item L'émission est \textbf{instantanée} (délai inférieur à $10^{-9}$ s) dès que $\nu \geq \nu_0$, même à très faible intensité : un seul photon suffit à extraire un électron.
\end{itemize}
\end{aretenir}

\begin{methode}[Savoir si une radiation extrait des électrons, puis calculer $E_{c\,\max}$]
\begin{enumerate}
\item Identifier le métal et relever $W_0$ (ou $\nu_0$, ou $\lambda_0$). Convertir $W_0$ en joules si nécessaire.
\item Calculer la longueur d'onde seuil : $\lambda_0 = \dfrac{hc}{W_0}$.
\item \textbf{Comparer} : si $\lambda \leq \lambda_0$ (équivalentment $\nu \geq \nu_0$), il y a émission ; sinon, aucun photoélectron, quelle que soit l'intensité.
\item Si émission, calculer : $E_{c\,\max} = hc\left(\dfrac{1}{\lambda} - \dfrac{1}{\lambda_0}\right)$, puis convertir en eV si demandé.
\item Vérifier la cohérence : $E_{c\,\max}$ doit être positive et de l'ordre de quelques eV.
\end{enumerate}
\end{methode}

\begin{attention}[Erreur fréquente : mélanger $\lambda$ et $\nu$]
Le piège classique du Bac : on te donne $\lambda$ et tu appliques $E = h\nu$ avec $\nu = \lambda$ ! Retiens que $\nu = \dfrac{c}{\lambda}$ : ce sont des grandeurs \textbf{inverses l'une de l'autre}. Une grande longueur d'onde signifie une \textbf{petite} fréquence et un photon \textbf{peu énergétique}. Autre piège : la condition d'émission s'écrit $\lambda \leq \lambda_0$ (inégalité \textbf{inversée} par rapport à $\nu \geq \nu_0$). Enfin, n'oublie jamais de convertir les nm en m ($\SI{1}{nm} = 10^{-9}$ m) avant d'utiliser $h = \SI{6,63e-34}{J.s}$.
\end{attention}

\courssub{Exemple chiffré complet : le sodium éclairé en lumière violette}

\begin{exoresolu}[Effet photoélectrique sur le sodium]
Le sodium a un travail d'extraction $W_0 = \SI{2,28}{eV}$. On l'éclaire avec une radiation de longueur d'onde $\lambda = \SI{400}{nm}$ (lumière violette).
\begin{enumerate}
\item Calculer la longueur d'onde seuil $\lambda_0$ du sodium. Y a-t-il émission de photoélectrons ?
\item Calculer l'énergie cinétique maximale des photoélectrons émis, en joules puis en eV.
\item Que se passe-t-il si on éclaire le sodium avec une lumière rouge de longueur d'onde $\lambda' = \SI{700}{nm}$, même très intense ?
\end{enumerate}
On donne : $h = \SI{6,63e-34}{J.s}$ ; $c = \SI{3,00e8}{m/s}$ ; $\SI{1}{eV} = \SI{1,6e-19}{J}$.
\tcblower
\textbf{1. Longueur d'onde seuil.} Convertissons d'abord $W_0$ en joules :
\[
W_0 = 2{,}28 \times 1{,}6 \times 10^{-19} = \SI{3,648e-19}{J}.
\]
\[
\lambda_0 = \dfrac{hc}{W_0} = \dfrac{6{,}63 \times 10^{-34} \times 3{,}00 \times 10^{8}}{3{,}648 \times 10^{-19}} \approx \SI{5,45e-7}{m} = \SI{545}{nm}.
\]
Comme $\lambda = \SI{400}{nm} < \lambda_0 = \SI{545}{nm}$, la condition $\lambda \leq \lambda_0$ est satisfaite : \textbf{il y a émission} de photoélectrons.

\textbf{2. Énergie cinétique maximale.} Appliquons la relation d'Einstein :
\begin{align*}
E_{c\,\max} &= hc\left(\dfrac{1}{\lambda} - \dfrac{1}{\lambda_0}\right)\\
&= 6{,}63 \times 10^{-34} \times 3{,}00 \times 10^{8} \times \left(\dfrac{1}{4{,}00 \times 10^{-7}} - \dfrac{1}{5{,}45 \times 10^{-7}}\right)\\
&= 1{,}989 \times 10^{-25} \times (2{,}50 \times 10^{6} - 1{,}835 \times 10^{6})\\
&\approx \SI{1,32e-19}{J} \approx \SI{0,83}{eV}.
\end{align*}
Vérification rapide en eV : $h\nu = \dfrac{1240}{400} \approx \SI{3,10}{eV}$, donc $E_{c\,\max} = 3{,}10 - 2{,}28 = \SI{0,82}{eV}$. Les deux méthodes concordent (la légère différence vient de l'arrondi de $hc \approx 1240\ \mathrm{eV.nm}$).

\textbf{3. Lumière rouge.} Pour $\lambda' = \SI{700}{nm} > \lambda_0 = \SI{545}{nm}$, l'énergie du photon vaut $h\nu' = \dfrac{1240}{700} \approx \SI{1,77}{eV} < W_0 = \SI{2,28}{eV}$. Chaque photon est individuellement trop pauvre : \textbf{aucun électron n'est émis}, même si on multiplie l'intensité par mille. Augmenter l'intensité augmente le nombre de photons, pas leur énergie individuelle.
\end{exoresolu}

\courssub{Interprétation graphique : la droite $E_{c\,\max} = f(\nu)$ (complément Bac)}

La relation $E_{c\,\max} = h\nu - W_0$ est une fonction \textbf{affine} de la fréquence. Sa représentation graphique est remarquable et très souvent exploitée au Baccalauréat :

\begin{aretenir}[Lecture du graphe $E_{c\,\max} = f(\nu)$]
La courbe $E_{c\,\max} = h\nu - W_0$ est une \textbf{droite} :
\begin{itemize}
\item de \textbf{pente} $h$ (la constante de Planck !) : c'est ainsi que Millikan a mesuré $h$ avec précision en 1916, validant la théorie d'Einstein ;
\item d'\textbf{ordonnée à l'origine} $-W_0$ ;
\item d'\textbf{abscisse à l'origine} $\nu_0 = \dfrac{W_0}{h}$ : la fréquence seuil, point à partir duquel l'émission devient possible.
\end{itemize}
Toutes les droites obtenues pour des métaux différents sont \textbf{parallèles} (même pente $h$) ; seul le seuil $\nu_0$ change.
\end{aretenir}

\begin{popfigure}
\begin{center}
\begin{tikzpicture}
\begin{axis}[
  width=0.85\linewidth, height=7cm,
  axis lines=middle,
  xlabel={$\nu$ ($10^{14}$ Hz)}, ylabel={$E_{c\,\max}$ (eV)},
  xmin=0, xmax=14, ymin=-2.8, ymax=3.2,
  xtick={5.5}, xticklabels={$\nu_0$},
  ytick=\empty,
  every axis x label/.style={at={(ticklabel* cs:1.02)},anchor=west},
  every axis y label/.style={at={(ticklabel* cs:1.02)},anchor=south},
  clip=false
]
\% Zone interdite
\fill[poporangeL,opacity=0.5] (axis cs:0,0) rectangle (axis cs:5.5,3.2);
\node[poporange!70!popdark,align=center,font=\small] at (axis cs:2.75,1.9) {pas d'émission\\ $\nu < \nu_0$};
\% Droite d'Einstein (sodium : W0 = 2.28 eV, nu0 = 5.5e14 Hz)
\addplot[popcurve,domain=5.5:13.2,samples=50]{0.414*x - 2.28};
\% Prolongement en pointillés
\addplot[popblue,dashed,domain=0:5.5,samples=20]{0.414*x - 2.28};
\% Points remarquables
\addplot[only marks,mark=*,popink] coordinates {(5.5,0) (0,-2.28)};
\node[popink,anchor=west] at (axis cs:0.2,-2.28) {$-W_0$};
\node[popblue,anchor=south west,font=\small] at (axis cs:8.5,1.6) {$E_{c\,\max} = h\nu - W_0$};
\node[popblue,anchor=north west,font=\small] at (axis cs:8.5,1.3) {pente $= h$};
\end{axis}
\end{tikzpicture}
\end{center}
\legende{Énergie cinétique maximale des photoélectrons en fonction de la fréquence (cas du sodium, $W_0 = \SI{2,28}{eV}$, $\nu_0 \approx \SI{5,5e14}{Hz}$). La droite, de pente $h$, coupe l'axe des fréquences en $\nu_0$ ; en dessous du seuil, aucune émission n'est possible.}
\end{popfigure}

\begin{methode}[Exploiter un graphe $E_{c\,\max} = f(\nu)$ au Bac]
\begin{enumerate}
\item Lire l'abscisse à l'origine : c'est $\nu_0$ ; en déduire $W_0 = h\nu_0$ et $\lambda_0 = c/\nu_0$.
\item Calculer la pente à partir de deux points éloignés : $h = \dfrac{\Delta E_{c\,\max}}{\Delta \nu}$. Attention aux unités des axes (souvent eV et $10^{14}$ Hz) : convertis pour retrouver $h \approx \SI{6,63e-34}{J.s}$.
\item Pour une fréquence $\nu$ donnée, lire $E_{c\,\max}$ directement ou calculer $h(\nu - \nu_0)$.
\end{enumerate}
\end{methode}

\courssub{Ordres de grandeur et choix des métaux}

\begin{center}
\begin{tabularx}{\linewidth}{|l|c|c|X|}
\hline
\rowcolor{popblueL} \textbf{Métal} & $W_0$ (eV) & $\lambda_0$ (nm) & \textbf{Domaine seuil} \\
\hline
Césium & 1,9 & 653 & rouge / visible \\
\hline
Potassium & 2,3 & 539 & vert / visible \\
\hline
Sodium & 2,28 & 544 & vert / visible \\
\hline
Zinc & 4,3 & 289 & ultraviolet \\
\hline
Cuivre & 4,7 & 264 & ultraviolet \\
\hline
\end{tabularx}
\end{center}

On constate que les \textbf{métaux alcalins} (césium, potassium, sodium) possèdent les plus faibles travaux d'extraction : leur électron périphérique unique est faiblement lié au réseau. C'est pourquoi ils constituent la cathode des cellules photoélectriques sensibles à la lumière visible. À l'inverse, le zinc ou le cuivre n'émettent que sous ultraviolet.

\begin{savaistu}[Pourquoi le césium règne dans les cellules photoélectriques]
Avec $W_0 = \SI{1,9}{eV}$, le césium répond à toute la lumière visible jusqu'au rouge : c'est le métal le plus « photo-sensible » du tableau périodique. Les cellules photoélectriques des anciens cinémas sonores (qui lisaient la piste optique des films), les photomultiplicateurs des hôpitaux de Yaoundé et de Douala utilisés en imagerie médicale, et les détecteurs de flamme des installations pétrolières de Kribi exploitent des couches alcalines à très faible $W_0$. Notons aussi que c'est pour son interprétation de l'effet photoélectrique — et non pour la relativité — qu'Einstein a reçu le prix Nobel de physique en 1921 : la communauté scientifique y voyait la preuve irréfutable de l'existence des photons.
\end{savaistu}

\begin{exercice}[Fréquence seuil et graphe]
Une cellule photoélectrique à cathode de potassium ($W_0 = \SI{2,3}{eV}$) est éclairée successivement par trois radiations : $\lambda_1 = \SI{400}{nm}$, $\lambda_2 = \SI{540}{nm}$, $\lambda_3 = \SI{650}{nm}$.
\begin{enumerate}
\item Calculer la fréquence seuil $\nu_0$ et la longueur d'onde seuil $\lambda_0$ du potassium.
\item Pour chaque radiation, dire s'il y a émission et, le cas échéant, calculer $E_{c\,\max}$ en eV.
\item Tracer l'allure de $E_{c\,\max} = f(\nu)$ pour ce métal et placer les points remarquables.
\end{enumerate}
\end{exercice}

\begin{corrige}[Exercice : Fréquence seuil et graphe]
\textbf{1.} $W_0 = 2{,}3 \times 1{,}6 \times 10^{-19} = \SI{3,68e-19}{J}$.
\[
\nu_0 = \dfrac{W_0}{h} = \dfrac{3{,}68 \times 10^{-19}}{6{,}63 \times 10^{-34}} \approx \SI{5,55e14}{Hz} \ ; \quad \lambda_0 = \dfrac{c}{\nu_0} \approx \SI{540}{nm}.
\]
\textbf{2.} Comparaison à $\lambda_0 = \SI{540}{nm}$ :
\begin{itemize}
\item $\lambda_1 = \SI{400}{nm} < \lambda_0$ : émission, $E_{c\,\max} = \dfrac{1240}{400} - 2{,}3 = 3{,}10 - 2{,}3 = \SI{0,80}{eV}$.
\item $\lambda_2 = \SI{540}{nm} \approx \lambda_0$ : on est pratiquement au seuil, $E_{c\,\max} \approx 0$ (électrons émis sans vitesse).
\item $\lambda_3 = \SI{650}{nm} > \lambda_0$ : $h\nu = \dfrac{1240}{650} \approx \SI{1,91}{eV} < W_0$ : \textbf{pas d'émission}.
\end{itemize}
\textbf{3.} Droite de pente $h$ passant par $(\nu_0\,;\,0)$ avec $\nu_0 = \SI{5,55e14}{Hz}$ et d'ordonnée à l'origine $-W_0 = \SI{-2,3}{eV}$ ; le point $(\SI{7,5e14}{Hz}\,;\,\SI{0,80}{eV})$ correspond à $\lambda_1$.
\end{corrige}

En résumé, le modèle du photon d'Einstein rend compte de tous les faits expérimentaux : l'existence d'un seuil ($W_0 = h\nu_0$), l'instantanéité de l'émission, l'indépendance de $E_{c\,\max}$ vis-à-vis de l'intensité ($E_{c\,\max} = h\nu - W_0$). Reste une question : comment mesurer expérimentalement cette énergie cinétique maximale et le débit des photoélectrons ? C'est précisément l'objet de la section suivante, consacrée à la caractéristique $I = f(U)$ de la cellule photoélectrique.

\coursec{Caractéristique I = f(U) d'une cellule photoélectrique}

Dans la section précédente, nous avons établi que l'effet photoélectrique ne se produit que si la fréquence $\nu$ de la radiation incidente dépasse le seuil $\nu_0$ du métal, et que l'énergie cinétique maximale des photoélectrons obéit à la relation d'Einstein $E_{c\max} = h\nu - W_0$. Une question naturelle se pose alors : comment \emph{mesurer} concrètement ces grandeurs ? Comment vérifier expérimentalement le modèle du photon ? La réponse tient en un instrument remarquable : la \cle{caractéristique} $I = f(U)$ de la cellule photoélectrique. C'est elle qui a permis à Millikan, en 1916, de confirmer la théorie d'Einstein et de mesurer la constante de Planck avec une précision remarquable. Suivons sa démarche.

\courssub{Montage expérimental}

\begin{experience}[Montage d'étude de la cellule photoélectrique]
\textbf{Matériel :} une cellule photoélectrique (ampoule de verre sous vide contenant une cathode $C$ recouverte d'un métal photosensible et une anode $A$), un générateur de tension continue \emph{réglable et réversible} délivrant la tension $U_{AC}$, un ampèremètre sensible (microampèremètre), un voltmètre, une source lumineuse monochromatique de fréquence $\nu$ connue.

\textbf{Protocole :}
\begin{enumerate}
\item Éclairer la cathode $C$ avec la radiation de fréquence $\nu > \nu_0$ : des photoélectrons sont émis.
\item Faire varier la tension $U_{AC}$ de valeurs négatives (tension \emph{inverse} : l'anode est repoussante) à des valeurs positives (tension \emph{directe} : l'anode attire les électrons).
\item Pour chaque valeur de $U$, relever l'intensité $I$ du courant dans le circuit.
\item Tracer la courbe $I = f(U)$, dite \cle{caractéristique} de la cellule.
\end{enumerate}

\textbf{Observations :} pour $U$ suffisamment négative, le courant est nul ; quand $U$ augmente, $I$ croît puis atteint un \emph{palier} : le courant sature. En augmentant l'éclairement de la cathode, le palier monte, mais la tension d'annulation du courant ne change pas.
\end{experience}

\begin{popfigure}
\begin{center}
\begin{tikzpicture}[european, scale=1.0, transform shape]
% Cellule
\draw[thick, popblue] (6,1.5) circle (1.5);
\draw[very thick, popdark] (5.4,1.0) -- (5.4,2.0);
\draw[very thick, popdark] (5.15,1.0) -- (5.15,2.0);
\node[popdark] at (5.0,0.6) {$C$};
\draw[very thick, popdark] (6.6,1.5) circle (0.12);
\node[popdark] at (6.9,1.15) {$A$};
% Circuit
\draw (5.3,1.0) -- (5.3,0) -- (2,0) to[battery1, l_={$U$}] (2,-2) -- (8,-2) to[ammeter, l={$\mu$}A] (8,0) -- (6.6,0) -- (6.6,1.38);
\draw (3.6,0) to[voltmeter, l=V] (3.6,-2);
% Lumière
\draw[-{Stealth}, very thick, poporange] (3.6,2.6) -- (4.9,1.9);
\draw[-{Stealth}, very thick, poporange] (3.9,3.0) -- (5.1,2.3);
\draw[-{Stealth}, very thick, poporange] (4.2,3.4) -- (5.3,2.7);
\node[poporange] at (3.4,3.2) {$h\nu$};
% Sens du courant conventionnel (polarisation directe U>0)
\draw[-{Stealth}, thick, popgreen] (7.6,-1.6) -- (6.8,-1.6);
\node[popgreen] at (7.2,-1.25) {$I$};
% Electron
\draw[-{Stealth}, thick, poppurple] (5.6,1.7) -- (6.4,1.55);
\node[poppurple] at (6.0,2.05) {$e^-$};
\end{tikzpicture}
\end{center}
\legende{Montage d'étude de la caractéristique d'une cellule photoélectrique. La lumière de fréquence $\nu$ arrache des électrons à la cathode $C$ ; la tension réglable $U$ contrôle leur trajet vers l'anode $A$.}
\end{popfigure}

\begin{attention}[Polarité de la tension]
On note $U = U_{AC}$ la tension de l'anode par rapport à la cathode. Si $U > 0$, l'anode est positive : elle \emph{attire} les électrons (tension directe, champ accélérateur). Si $U < 0$, l'anode est négative : elle \emph{repousse} les électrons (tension inverse, champ freinateur). Le générateur doit donc pouvoir délivrer les deux polarités : c'est le point délicat du montage.
\end{attention}

\courssub{Allure de la caractéristique : croissance puis saturation}

À éclairement et fréquence fixés, la courbe $I = f(U)$ présente trois régimes :

\begin{itemize}
\item \textbf{Régime 1 — tension fortement négative} ($U < -U_0$) : le champ électrique entre $C$ et $A$ est si freinateur qu'\emph{aucun} électron, même le plus rapide, n'atteint l'anode : $I = 0$.
\item \textbf{Régime 2 — zone de croissance} ($-U_0 < U < U_{\text{sat}}$) : les électrons les plus énergétiques franchissent la barrière de potentiel ; leur nombre augmente avec $U$, donc $I$ croît.
\item \textbf{Régime 3 — palier de saturation} ($U$ suffisamment positif) : \emph{tous} les photoélectrons émis par la cathode sont collectés par l'anode. Augmenter encore $U$ ne sert à rien : le courant plafonne à la valeur $I_{\text{sat}}$, le \cle{courant de saturation}.
\end{itemize}

\begin{popfigure}
\begin{center}
\begin{tikzpicture}
\begin{axis}[
width=0.85\linewidth, height=7cm,
xlabel={$U$ (V)}, ylabel={$I$ ($\mu$A)},
xmin=-2.4, xmax=4.2, ymin=-0.4, ymax=5.2,
axis lines=middle,
xtick={-1.2}, xticklabels={$-U_0$},
ytick={2,4}, yticklabels={$I_{\text{sat}1}$,$I_{\text{sat}2}$},
samples=200, domain=-2.4:4.2,
every axis x label/.style={at={(ticklabel* cs:1.02)}, anchor=west},
every axis y label/.style={at={(ticklabel* cs:1.02)}, anchor=south},
]
\addplot[popcurve, popblue, very thick] {4/(1+exp(-2.2*(x+0.4)))};
\addplot[popcurve, poporange, very thick] {2/(1+exp(-2.2*(x+0.4)))};
\addplot[dashed, popmuted] coordinates {(-1.2,0) (-1.2,4.6)};
\node[popblue, anchor=west] at (axis cs:2.2,4.15) {fort éclairement};
\node[poporange, anchor=west] at (axis cs:2.2,2.15) {faible éclairement};
\node[popmuted, anchor=south west] at (axis cs:-2.3,0.1) {};
\end{axis}
\end{tikzpicture}
\end{center}
\legende{Caractéristiques $I = f(U)$ pour deux éclairements, à fréquence $\nu$ fixée. Le courant de saturation dépend de l'éclairement ; le potentiel d'arrêt $-U_0$ est le même pour les deux courbes.}
\end{popfigure}

\courssub{Courant de saturation $I_{\text{sat}}$}

Sur le palier, chaque électron émis par la cathode rejoint l'anode : le débit d'électrons dans le circuit est exactement le débit d'émission photoélectrique. Si $N_e$ électrons sont émis pendant la durée $\Delta t$, alors :
\[ I_{\text{sat}} = \frac{N_e \, e}{\Delta t}. \]

Or, dans l'interprétation photonique, chaque photon absorbé libère (au plus) un électron. Doubler l'intensité lumineuse, c'est doubler le nombre de photons incidents par seconde, donc doubler le nombre d'électrons émis par seconde.

\begin{propriete}[Courant de saturation]
Le courant de saturation $I_{\text{sat}}$ est \cle{proportionnel à l'intensité lumineuse} (à l'éclairement) reçue par la cathode, à fréquence fixée. Il ne dépend \emph{pas} de la tension $U$ une fois le palier atteint.
\end{propriete}

C'est un résultat que la physique ondulatoire classique expliquait déjà, mais elle échouait sur le point suivant.

\courssub{Potentiel d'arrêt $U_0$}

\begin{definition}[Potentiel d'arrêt]
Le \cle{potentiel d'arrêt} $U_0$ est la \emph{valeur absolue} de la tension inverse $U_{AC} = -U_0$ qui annule le photocourant. Pour $U_{AC} \leq -U_0$, aucun électron n'atteint l'anode : $I = 0$.
\end{definition}

Physiquement, quand le courant s'annule, c'est que même l'électron émis avec l'énergie cinétique \emph{maximale} $E_{c\max}$ arrive à l'anode avec une vitesse nulle : il est stoppé juste avant de la toucher. Le travail de la force électrique freinatrice a exactement compensé son énergie cinétique initiale.

\begin{propriete}[Relation entre potentiel d'arrêt et énergie cinétique maximale]
\[ e\,U_0 = E_{c\max} = h\nu - W_0 = h(\nu - \nu_0). \]
Cette relation est \emph{exigible} au Baccalauréat ; sa démonstration par le théorème de l'énergie cinétique est détaillée ci-dessous.
\end{propriete}

\begin{experience}[Démonstration guidée : $eU_0 = E_{c\max}$]
Considérons un photoélectron émis de la cathode $C$ avec la vitesse maximale $v_{\max}$, donc avec l'énergie cinétique $E_{c\max} = \dfrac{1}{2} m v_{\max}^2$. On applique la tension inverse $U_{AC} = -U_0$ : l'anode est au potentiel le plus bas, le champ électrique $\vec{E}$ est dirigé de $C$ vers $A$, et la force électrique $\vec{F} = -e\vec{E}$ sur l'électron (charge $q = -e$) est dirigée de $A$ vers $C$ : elle s'oppose au mouvement.

\textbf{Étape 1 — bilan des forces.} Dans le tube à vide, le poids de l'électron ($P = mg \approx 9 \times 10^{-30}$ N) est négligeable devant la force électrique. Seule $\vec{F} = q\vec{E} = -e\vec{E}$ travaille.

\textbf{Étape 2 — travail de la force électrique.} Le travail de la force électrique entre $C$ et $A$ ne dépend que des potentiels :
\[ W_{C \to A}(\vec{F}) = q\,U_{CA} = (-e)(V_C - V_A) = (-e)(-U_{AC}) = e\,U_{AC} = -e\,U_0. \]
Le travail est négatif : la force est bien freinatrice.

\textbf{Étape 3 — théorème de l'énergie cinétique.} Entre la cathode (énergie $E_{c\max}$) et l'anode que l'électron atteint tout juste avec une vitesse nulle (énergie $0$) :
\[ \Delta E_c = W_{C \to A}(\vec{F}) \quad\Longrightarrow\quad 0 - E_{c\max} = -e\,U_0. \]

\textbf{Étape 4 — conclusion.}
\[ \boxed{\; e\,U_0 = E_{c\max} \;} \]
Le potentiel d'arrêt mesure donc directement l'énergie cinétique maximale des photoélectrons : c'est un véritable « énergiemètre » à électrons.
\end{experience}

\begin{exemplebox}[Du potentiel d'arrêt à la vitesse des photoélectrons]
On mesure un potentiel d'arrêt $U_0 = \num{1,2}$ V pour une cathode éclairée par une radiation ultraviolette.
\begin{itemize}
\item Énergie cinétique maximale : $E_{c\max} = eU_0 = \num{1,2}$ eV, soit en joules :
\[ E_{c\max} = 1,2 \times 1,60 \times 10^{-19} \approx \num{1,92e-19} \text{ J}. \]
\item Vitesse maximale des photoélectrons :
\[ v_{\max} = \sqrt{\frac{2 E_{c\max}}{m}} = \sqrt{\frac{2 \times 1,92 \times 10^{-19}}{9,11 \times 10^{-31}}} \approx \num{6,5e5} \text{ m/s}, \]
soit environ \SI{650}{km/s} : plus de 60\,000 fois la vitesse d'un mototaxi à pleine vitesse sur le boulevard de la Réunification !
\end{itemize}
\end{exemplebox}

\courssub{Mesure expérimentale de la constante de Planck}

Combinons la relation précédente avec la formule d'Einstein :
\[ eU_0 = h\nu - W_0 \quad\Longrightarrow\quad U_0 = \frac{h}{e}\,\nu - \frac{W_0}{e}. \]

\begin{propriete}[Droite de Millikan]
Pour un métal donné, le potentiel d'arrêt est une fonction \emph{affine croissante} de la fréquence $\nu$ de la radiation incidente. La droite $U_0 = f(\nu)$ a pour :
\begin{itemize}
\item \cle{pente} $a = \dfrac{h}{e}$ (indépendante du métal : droites parallèles pour tous les métaux) ;
\item \cle{ordonnée à l'origine} $b = -\dfrac{W_0}{e}$ ;
\item \cle{intersection avec l'axe des fréquences} : $\nu_0$, la fréquence de seuil du métal.
\end{itemize}
En mesurant $U_0$ pour plusieurs fréquences connues (raies du mercure, LED, lasers), on détermine expérimentalement $h = e \times a$. C'est ainsi que Millikan valida la théorie d'Einstein en 1916.
\end{propriete}

\courssub{Rendement quantique}

En réalité, tous les photons incidents ne parviennent pas à extraire un électron : certains sont réfléchis par la surface, d'autres communiquent leur énergie à des électrons qui la perdent en collisions avant de sortir du métal.

\begin{definition}[Rendement quantique]
Le \cle{rendement quantique} d'une cellule photoélectrique est le rapport
\[ \eta = \frac{N_e}{N_{ph}} = \frac{\text{nombre d'électrons émis}}{\text{nombre de photons incidents}}. \]
C'est un nombre sans dimension, toujours inférieur à 1. Pour les métaux usuels, il est de l'ordre de $10^{-2}$ à $10^{-1}$ : il faut typiquement 10 à 100 photons pour extraire un électron.
\end{definition}

Le rendement quantique relie la puissance lumineuse reçue au courant de saturation. Si la cathode reçoit une puissance $\mathcal{P}$ sous une radiation de fréquence $\nu$, le nombre de photons incidents par seconde est $\dfrac{\mathcal{P}}{h\nu}$, d'où :
\[ I_{\text{sat}} = \eta \, e \, \frac{\mathcal{P}}{h\nu}. \]
Cette formule montre clairement que $I_{\text{sat}}$ est proportionnel à la puissance lumineuse, comme annoncé plus haut.

\courssub{Influence de l'intensité lumineuse : le test décisif du modèle photonique}

Reprenons l'expérience en doublant l'éclairement, à fréquence $\nu$ constante. On observe (figure ci-dessus) :
\begin{itemize}
\item le courant de saturation \emph{double} : $I_{\text{sat}2} = 2\,I_{\text{sat}1}$ ;
\item le potentiel d'arrêt \emph{ne change pas} : $U_0$ reste identique.
\end{itemize}

L'interprétation photonique est limpide : augmenter l'intensité, c'est augmenter le \emph{nombre} de photons (donc le nombre d'électrons émis, donc $I_{\text{sat}}$), sans changer l'\emph{énergie} $h\nu$ de chaque photon (donc sans changer $E_{c\max}$, donc sans changer $U_0$). La théorie ondulatoire classique prévoyait au contraire qu'une lumière plus intense communiquerait plus d'énergie à chaque électron et augmenterait $U_0$ : l'expérience la contredit formellement.

\begin{attention}[Ne pas confondre $U_0$ et $I_{\text{sat}}$]
C'est l'erreur la plus fréquente au Baccalauréat. Retiens bien :
\begin{itemize}
\item $U_0$ (potentiel d'arrêt) dépend de la \cle{fréquence} $\nu$ de la lumière et de la nature du métal ($W_0$) : $eU_0 = h\nu - W_0$. Il est \emph{indépendant} de l'intensité lumineuse.
\item $I_{\text{sat}}$ (courant de saturation) dépend de l'\cle{intensité} lumineuse (et du rendement quantique). Il est \emph{indépendant} de la fréquence, tant que $\nu > \nu_0$.
\end{itemize}
Ces deux grandeurs sont indépendantes l'une de l'autre. Dire « plus la lumière est intense, plus les électrons sont rapides » est \textbf{faux} : ils sont seulement plus \emph{nombreux}.
\end{attention}

\begin{methode}[Exploiter une caractéristique $I = f(U)$]
\begin{enumerate}
\item \textbf{Lire $I_{\text{sat}}$} sur le palier de la courbe (ordonnée de la partie horizontale).
\item \textbf{Lire $U_0$} à l'intersection de la courbe avec l'axe des tensions : c'est la valeur absolue de l'abscisse du point où $I$ s'annule.
\item \textbf{En déduire l'énergie cinétique maximale} : $E_{c\max} = e\,U_0$ (en joules si $U_0$ en volts, ou directement $E_{c\max}$ en eV égale numériquement $U_0$ en volts).
\item \textbf{En déduire la fréquence de seuil} si $\nu$ est connue : $\nu_0 = \nu - \dfrac{eU_0}{h}$, puis la longueur d'onde de seuil $\lambda_0 = \dfrac{c}{\nu_0}$.
\item \textbf{Comparer deux courbes} : même $U_0$ $\Rightarrow$ même fréquence ; paliers différents $\Rightarrow$ éclairements différents.
\end{enumerate}
\end{methode}

\begin{exoresolu}[Exploitation complète d'une caractéristique]
Une cellule photoélectrique à cathode de césium est éclairée par une radiation de longueur d'onde $\lambda = \num{450}$ nm. Sa caractéristique $I = f(U)$ présente un palier à $I_{\text{sat}} = \num{8,0}$ $\mu$A et annule le courant pour $U_{AC} = -\num{0,95}$ V.

\textbf{Données :} $h = \num{6,63e-34}$ J·s ; $c = \num{3,00e8}$ m/s ; $e = \num{1,60e-19}$ C ; $\num{1}$ eV $= \num{1,60e-19}$ J.

\begin{enumerate}
\item Déterminer l'énergie cinétique maximale des photoélectrons.
\item Déterminer le travail d'extraction $W_0$ et la longueur d'onde de seuil $\lambda_0$ du césium.
\item On double l'éclairement sans changer $\lambda$. Que deviennent $I_{\text{sat}}$ et $U_0$ ?
\item Calculer le nombre d'électrons émis par seconde sur le palier.
\end{enumerate}

\tcblower

\textbf{1. Énergie cinétique maximale.} Le potentiel d'arrêt est $U_0 = \num{0,95}$ V, donc :
\[ E_{c\max} = eU_0 = \num{0,95} \text{ eV} = 0,95 \times 1,60 \times 10^{-19} = \num{1,52e-19} \text{ J}. \]

\textbf{2. Travail d'extraction et seuil.} L'énergie du photon incident :
\[ E_{\text{photon}} = \frac{hc}{\lambda} = \frac{6,63 \times 10^{-34} \times 3,00 \times 10^{8}}{450 \times 10^{-9}} = \num{4,42e-19} \text{ J} \approx \num{2,76} \text{ eV}. \]
Relation d'Einstein :
\[ W_0 = h\nu - E_{c\max} = 4,42 \times 10^{-19} - 1,52 \times 10^{-19} = \num{2,90e-19} \text{ J} \approx \num{1,81} \text{ eV}. \]
Longueur d'onde de seuil :
\[ \lambda_0 = \frac{hc}{W_0} = \frac{6,63 \times 10^{-34} \times 3,00 \times 10^{8}}{2,90 \times 10^{-19}} \approx \num{6,86e-7} \text{ m} = \num{686} \text{ nm}. \]
Cohérent : le césium est photosensible dans le visible, c'est pourquoi on l'utilise dans les cellules des automatismes d'éclairage public.

\textbf{3. Doublement de l'éclairement.} Le nombre de photons incidents par seconde double, donc le nombre d'électrons émis double : $I_{\text{sat}}' = \num{16}$ $\mu$A. L'énergie de chaque photon est inchangée, donc $E_{c\max}$ est inchangée : $U_0' = \num{0,95}$ V, \emph{inchangé}.

\textbf{4. Débit d'électrons.} Sur le palier, $I_{\text{sat}} = \dfrac{N_e e}{\Delta t}$, d'où :
\[ \frac{N_e}{\Delta t} = \frac{I_{\text{sat}}}{e} = \frac{8,0 \times 10^{-6}}{1,60 \times 10^{-19}} = \num{5,0e13} \text{ électrons par seconde}. \]
\end{exoresolu}

\begin{aretenir}[L'essentiel de la section]
\begin{itemize}
\item La caractéristique $I = f(U)$ d'une cellule photoélectrique croît puis atteint un \cle{palier de saturation} $I_{\text{sat}}$, proportionnel à l'intensité lumineuse.
\item Le \cle{potentiel d'arrêt} $U_0$ annule le photocourant ; le théorème de l'énergie cinétique donne $eU_0 = E_{c\max} = h\nu - W_0$.
\item $U_0$ dépend de la fréquence (droite $U_0 = \dfrac{h}{e}\nu - \dfrac{W_0}{e}$, dont la pente permet de mesurer $h$) ; $I_{\text{sat}}$ dépend de l'éclairement. Les deux sont indépendants.
\item Le \cle{rendement quantique} $\eta = \dfrac{N_e}{N_{ph}} < 1$ mesure l'efficacité de conversion photons $\to$ électrons.
\end{itemize}
\end{aretenir}

\begin{savaistu}[Des cellules photoélectriques aux portes automatiques]
Les cellules photoélectriques équipent les barrières de sécurité des portes automatiques, les compteurs de production et les systèmes d'alarme : un faisceau lumineux interrompu fait chuter brutalement le photocourant, ce qui déclenche le signal. Les photomultiplicateurs, descendants sophistiqués de la cellule de Hertz, amplifient un unique photoélectron en une avalanche de millions d'électrons : ils équipent les caméras des télescopes et les appareils d'imagerie médicale (tomographie par émission de positons) utilisés dans les grands hôpitaux. Quant aux panneaux solaires qui électrifient les villages du Nord-Cameroun, ils reposent sur l'\emph{effet photovoltaïque}, cousin de l'effet photoélectrique, que nous étudierons dans la section suivante.
\end{savaistu}

\coursec{Applications de l'effet photoélectrique}

Dans la section précédente, nous avons tracé et exploité la caractéristique $I = f(U)$ d'une cellule photoélectrique : courant de saturation, potentiel d'arrêt $U_0$ tel que $eU_0 = E_{c\max}$, rendement quantique. Cette cellule n'est pas qu'un instrument de laboratoire : c'est un \emph{détecteur de lumière} extrêmement sensible et rapide. Toute la physique que nous venons d'établir — un photon d'énergie $h\nu$ suffisante libère un électron, et le courant qui en résulte est proportionnel au flux lumineux — trouve des applications immenses, de la porte automatique du supermarché au kit solaire qui éclaire un village de l'Adamaoua. C'est l'objet de cette section : montrer comment une découverte de physique fondamentale (Hertz, 1887 ; Einstein, 1905) est devenue une technologie quotidienne, et apprendre à en faire les bilans énergétiques exigibles au Baccalauréat.

\courssub{La cellule photoélectrique comme détecteur de lumière}

\textbf{Principe général.} Une cellule photoélectrique (à émission externe : cathode en métal alcalin dans une ampoule sous vide, anode collectrice) convertit un \cle{flux lumineux} en \cle{courant électrique}. D'après l'étude de la caractéristique, lorsque la tension est suffisante pour saturer le courant, l'intensité $I_{\text{sat}}$ est proportionnelle au nombre de photons reçus par seconde, donc à l'éclairement. La cellule est donc un \emph{luxmètre électrique} : mesurer $I$, c'est mesurer la lumière.

\textbf{Barrière photoélectrique.} L'application la plus intuitive est la détection de passage. On place face à face un émetteur (diode électroluminescente ou lampe) et une cellule réceptrice. Tant que le faisceau arrive sur la cellule, un courant circule et maintient un relais dans un état donné. Si un objet — une personne, un colis, un véhicule — coupe le faisceau, le courant s'annule brutalement : le relais bascule et déclenche une action (ouverture de porte, arrêt d'urgence d'une machine, incrémentation d'un compteur).

\begin{popfigure}
\begin{center}
\begin{tikzpicture}[scale=1.0, >=Stealth]
\draw[fill=popblueL, draw=popblue] (0,-0.6) rectangle (1.2,0.6);
\node[popdark] at (0.6,0) {\small Émetteur};
\node[popdark, below] at (0.6,-0.65) {\scriptsize (DEL)};
\draw[fill=poporangeL, draw=poporange] (6.8,-0.6) rectangle (8,0.6);
\node[popdark] at (7.4,0) {\small Cellule};
\node[popdark, below] at (7.4,-0.65) {\scriptsize réceptrice};
\draw[popgold, line width=1.6pt, ->] (1.2,0) -- (6.8,0);
\foreach \x in {2.0,3.2,4.4,5.6} {\draw[popgold, ->] (\x,0.12) -- (\x+0.45,0.12);}
\node[popgold!70!popdark, above] at (4,0.15) {\small faisceau lumineux};
\draw[fill=poppurpleL, draw=poppurple] (7.0,-2.2) rectangle (7.8,-1.5);
\node[popdark] at (7.4,-1.85) {\small Relais};
\draw[popline, ->] (7.4,-0.6) -- (7.4,-1.5);
\draw[fill=popgreenL, draw=popgreen] (9.0,-2.2) rectangle (11.0,-1.5);
\node[popdark, align=center] at (10,-1.85) {\small Action : porte,\\ \small compteur, alarme};
\draw[popline, ->] (7.8,-1.85) -- (9.0,-1.85);
\draw[fill=popmuted!40, draw=popmuted, dashed] (3.6,-0.9) rectangle (4.4,0.9);
\node[popmuted, below, align=center] at (4,-1.0) {\scriptsize obstacle : le faisceau\\ \scriptsize est coupé $\Rightarrow I = 0$};
\end{tikzpicture}
\end{center}
\legende{Principe d'une barrière photoélectrique : tant que le faisceau atteint la cellule, le photocourant maintient le relais ; la coupure du faisceau annule le courant et déclenche l'action.}
\end{popfigure}

\begin{exemplebox}[Portes automatiques et sécurité des machines]
Les portes coulissantes des supermarchés et des banques à Douala ou Yaoundé utilisent des barrières photoélectriques (souvent dans l'infrarouge, invisible pour l'œil) : dès qu'une personne interrompt le faisceau, la commande d'ouverture est activée. Dans l'industrie, une barrière immatérielle protège les presses : si la main de l'opérateur coupe le faisceau, la machine s'arrête en une fraction de seconde — la cellule photoélectrique répond en moins d'une milliseconde, ce qu'aucun interrupteur mécanique ne saurait faire.
\end{exemplebox}

\textbf{Comptage et lecture optique.} Le même principe sert à compter : flacons sur une chaîne de production, véhicules à un péage, tours d'une hélice. Historiquement, la cellule photoélectrique a aussi permis le \emph{son du cinéma} : la bande sonore optique du film module un faisceau lumineux, et la cellule transforme ces variations de lumière en signal électrique amplifié.

\courssub{Les panneaux photovoltaïques : de la lumière à l'électricité}

\textbf{De l'effet externe à l'effet interne.} Dans la cellule étudiée jusqu'ici, les électrons sont \emph{arrachés} d'un métal et circulent \emph{dans le vide} : c'est l'\cle{effet photoélectrique externe}. Dans un \cle{semi-conducteur} (silicium), un photon d'énergie suffisante peut libérer un électron \emph{à l'intérieur du matériau}, en l'arrachant à une liaison cristalline : c'est l'\cle{effet photoélectrique interne}. L'électron libéré est mis en mouvement par un champ électrique interne (jonction P-N), ce qui crée une tension et un courant utilisables dans un circuit extérieur. Aucun vide, aucune éjection : l'électron reste dans le matériau, mais le bilan d'énergie est le même — un photon absorbé transfère son énergie $h\nu$ à un électron.

\begin{definition}[Cellule photovoltaïque]
Une \cle{cellule photovoltaïque} est un composant à semi-conducteur (le plus souvent du silicium) qui convertit directement l'énergie lumineuse en énergie électrique par effet photoélectrique interne : chaque photon absorbé d'énergie $h\nu$ supérieure au seuil du matériau libère un porteur de charge mis en mouvement par la jonction. Un \cle{panneau solaire} est un assemblage de cellules montées en série et en parallèle.
\end{definition}

\begin{attention}[Effet externe ou effet interne ?]
Erreur très fréquente au Bac : confondre les deux effets. Dans l'\textbf{effet photoélectrique externe} (expérience de Hertz, cellule à vide), les électrons sont \emph{éjectés du métal} et circulent dans le vide ; il faut $h\nu \geq W_0$, travail d'extraction du métal. Dans l'\textbf{effet interne} (photovoltaïque, photorésistances), les électrons sont \emph{libérés à l'intérieur} du semi-conducteur et y circulent ; le seuil est l'énergie de la bande interdite (\emph{gap}) du matériau. Dans les deux cas, un photon = un électron mis en mouvement, mais le cadre est différent. Ne dis jamais qu'un panneau solaire « éjecte des électrons dans le vide » !
\end{attention}

\textbf{Bilan énergétique d'un panneau.} Un panneau ne convertit qu'une fraction de l'énergie lumineuse reçue : le reste est réfléchi ou transformé en chaleur. On définit le \cle{rendement} $\eta$ du panneau, et le bilan de puissance s'écrit :
\[ P_{\text{élec}} = \eta \times P_{\text{lum}} = \eta \times E \times S \]
où $E$ est l'éclairement (puissance lumineuse reçue par unité de surface, en $\mathrm{W/m^2}$) et $S$ la surface du panneau. Vérification dimensionnelle : $[\eta] = 1$ (sans dimension), $[E \times S] = \mathrm{W\,m^{-2}} \times \mathrm{m^2} = \mathrm{W}$ : la relation est homogène.

\begin{exemplebox}[Ordre de grandeur fondamental]
Par un beau jour de midi, l'éclairement solaire au sol vaut environ $E = \SI{1000}{W/m^2}$ (conditions standard d'essai). Un panneau de surface $S = \SI{1}{m^2}$ et de rendement $\eta = 15\,\%$ fournit :
\[ P_{\text{élec}} = 0{,}15 \times 1000 \times 1 = \SI{150}{W} \]
soit de quoi alimenter une dizaine de lampes LED ou recharger plusieurs téléphones. Pour alimenter toute une maison, il faut donc plusieurs mètres carrés de panneaux.
\end{exemplebox}

\begin{savaistu}[Le Cameroun, un pays fait pour le solaire]
L'ensoleillement moyen au Cameroun est d'environ $4$ à $5{,}5\ \mathrm{kWh/m^2/jour}$ selon les régions — les valeurs les plus élevées se trouvant dans le Grand Nord (Maroua, Garoua) —, ce qui est très favorable au photovoltaïque, comparable aux meilleurs sites d'Europe du Sud et supérieur à l'Allemagne, pourtant leader européen du solaire ! C'est pourquoi les \cle{kits solaires} (petit panneau + batterie + lampes LED) se diffusent rapidement dans les zones rurales non raccordées au réseau ENEO, et que de nombreux \cle{lampadaires solaires autonomes} éclairent désormais les villes : chaque lampadaire porte son panneau, sa batterie et sa lampe, sans aucun câble. L'énergie solaire complète ainsi les grands barrages hydroélectriques (Lom Pangar, Nachtigal) dans la diversification du mix énergétique national.
\end{savaistu}

\begin{popfigure}
\begin{center}
\begin{tikzpicture}[scale=1.0, >=Stealth]
\draw[fill=popgold!60, draw=popgold] (0,0) circle (0.55);
\foreach \a in {0,45,...,315} {\draw[popgold, line width=1pt] (\a:0.65) -- (\a:0.95);}
\node[popdark, below] at (0,-0.75) {\small Soleil};
\draw[fill=popblue, draw=popdark] (2.6,-0.5) rectangle (4.0,0.5);
\draw[popcream, line width=0.6pt] (2.6,0) -- (4.0,0);
\draw[popcream, line width=0.6pt] (3.3,-0.5) -- (3.3,0.5);
\node[popdark, below] at (3.3,-0.6) {\small Panneau};
\node[popdark, below] at (3.3,-1.05) {\small photovoltaïque};
\draw[fill=poppurpleL, draw=poppurple] (5.4,-0.5) rectangle (6.8,0.5);
\node[popdark, align=center] at (6.1,0) {\small Régulateur\\ \small de charge};
\draw[fill=popgreenL, draw=popgreen] (8.2,-0.5) rectangle (9.4,0.5);
\node[popdark] at (8.8,0) {\small Batterie};
\draw[fill=poporangeL, draw=poporange] (10.8,-0.5) rectangle (12.0,0.5);
\node[popdark, align=center] at (11.4,0) {\small Lampe,\\ \small téléphone};
\draw[popgold, line width=1.4pt, ->] (0.7,0) -- (2.5,0) node[midway, above, popdark]{\scriptsize $P_{\text{lum}}$};
\draw[popblue, line width=1.4pt, ->] (4.0,0) -- (5.4,0) node[midway, above, popdark]{\scriptsize $P_{\text{élec}}$};
\draw[popline, line width=1.2pt, ->] (6.8,0) -- (8.2,0);
\draw[popline, line width=1.2pt, ->] (9.4,0) -- (10.8,0);
\end{tikzpicture}
\end{center}
\legende{Chaîne solaire domestique (kit solaire) : le panneau convertit la puissance lumineuse en puissance électrique, le régulateur protège la batterie qui stocke l'énergie pour la nuit.}
\end{popfigure}

\begin{exoresolu}[Dimensionnement d'un kit solaire villageois]
Une maison d'un village de la région du Nord utilise un panneau de $S = \SI{2}{m^2}$ et de rendement $\eta = 16\,\%$. L'ensoleillement moyen y est de $\SI{5}{kWh/m^2/jour}$. \textbf{(1)} Calculer l'énergie lumineuse reçue par le panneau en une journée. \textbf{(2)} En déduire l'énergie électrique disponible chaque jour. \textbf{(3)} Peut-on alimenter chaque soir 4 lampes LED de $\SI{5}{W}$ pendant 5 h et recharger 2 téléphones ($\SI{10}{Wh}$ chacun) ?
\tcblower
\textbf{(1)} Énergie lumineuse reçue : $E_{\text{lum}} = 5\ \mathrm{kWh/m^2/jour} \times 2\ \mathrm{m^2} = \SI{10}{kWh}$ par jour.
\textbf{(2)} Énergie électrique produite : $E_{\text{élec}} = \eta \times E_{\text{lum}} = 0{,}16 \times 10 = \SI{1,6}{kWh}$ par jour.
\textbf{(3)} Besoins quotidiens : lampes : $4 \times 5\ \mathrm{W} \times 5\ \mathrm{h} = \SI{100}{Wh} = \SI{0,10}{kWh}$ ; téléphones : $2 \times 10 = \SI{20}{Wh} = \SI{0,02}{kWh}$. Total : $\SI{0,12}{kWh}$, très inférieur à $\SI{1,6}{kWh}$ produits. Le kit est largement suffisant : l'excédent stocké dans la batterie compense les jours nuageux. Remarque : en pratique on retranche les pertes du régulateur et de la batterie (rendement global de stockage $\approx 80\,\%$), ce qui laisse encore $\approx \SI{1,3}{kWh}$ disponibles.
\end{exoresolu}

\courssub{Le photomultiplicateur : voir un photon à la fois}

\textbf{Le problème.} Certaines lumières sont si faibles — scintillation d'un cristal traversé par une particule, lueur d'une galaxie lointaine — que le photocourant d'une simple cellule serait indétectable (quelques électrons par seconde !). Comment amplifier sans noyer le signal dans le bruit ?

\textbf{La solution : l'amplification en cascade.} Le \cle{photomultiplicateur} est un tube sous vide dans lequel le photoélectron arraché à la photocathode est accéléré par une tension vers une première électrode appelée \cle{dynode}. En la frappant, il en arrache \emph{plusieurs} électrons secondaires (typiquement 3 à 5), qui sont à leur tour accélérés vers une deuxième dynode, et ainsi de suite sur une dizaine d'étages. Si chaque étage multiplie le nombre d'électrons par $g$, le gain total est $G = g^n$ : avec $g = 4$ et $n = 10$, $G = 4^{10} \approx 10^6$. Un seul photon produit donc une impulsion d'environ un million d'électrons, facilement mesurable.

\begin{exemplebox}[Où trouve-t-on des photomultiplicateurs ?]
\begin{itemize}
\item \textbf{Médecine nucléaire} : dans les caméras à scintillation et les scanners TEP utilisés en imagerie hospitalière, le cristal scintillateur convertit les rayons gamma en flashs lumineux minuscules, détectés par des photomultiplicateurs.
\item \textbf{Astronomie} : détection des photons uniques venus d'étoiles lointaines.
\item \textbf{Physique des particules} : comptage des scintillations dans les détecteurs.
\end{itemize}
\end{exemplebox}

\courssub{Autres applications quotidiennes}

\begin{itemize}
\item \textbf{Lecteurs de codes-barres} : un faisceau lumineux balaye le code ; les barres noires absorbent, les espaces blancs réfléchissent. La cellule réceptrice transforme cette succession d'éclairs et d'ombres en signal électrique binaire, décodé par la caisse du magasin.
\item \textbf{Capteurs des appareils photo numériques et téléphones} : le capteur CCD ou CMOS est une mosaïque de millions de minuscules cellules photoélectriques à semi-conducteur (effet interne). Chaque pixel accumule des charges proportionnellement à la lumière reçue : c'est littéralement la relation d'Einstein qui écrit vos photos.
\item \textbf{Télécommandes infrarouges} : la télécommande émet des impulsions de lumière infrarouge ($\lambda \approx \SI{940}{nm}$) ; une photodiode réceptrice dans le téléviseur les convertit en signal électrique. L'infrarouge est invisible mais transporte bien l'information par photons.
\item \textbf{Photorésistances (LDR)} : leur résistance chute quand l'éclairement augmente ; elles commandent l'allumage automatique des lampadaires à la tombée de la nuit.
\end{itemize}

\begin{methode}[Exploiter un bilan énergétique photovoltaïque au Bac]
\begin{enumerate}
\item Identifier les données : éclairement $E$ (en $\mathrm{W/m^2}$), surface $S$, rendement $\eta$ (convertir le pourcentage en fraction), durée $t$.
\item Calculer la puissance lumineuse reçue : $P_{\text{lum}} = E \times S$.
\item Appliquer le rendement : $P_{\text{élec}} = \eta \, P_{\text{lum}}$.
\item Si une énergie est demandée : $\mathcal{E} = P \times t$ (attention aux unités : $t$ en secondes pour des joules, en heures pour des Wh).
\item Conclure avec un nombre raisonnable de chiffres significatifs et l'unité SI.
\end{enumerate}
\end{methode}

\begin{attention}[Pièges de rédaction et de calcul]
\begin{itemize}
\item Ne pas oublier de convertir le rendement : $15\,\% = 0{,}15$, pas $15$ !
\item Le rendement s'applique à la puissance \emph{reçue}, pas à la puissance fournie : c'est $P_{\text{élec}} = \eta P_{\text{lum}}$, jamais l'inverse.
\item Un rendement est toujours inférieur à 1 : si ton calcul donne $P_{\text{élec}} > P_{\text{lum}}$, il y a une erreur.
\item Ne pas écrire qu'un panneau solaire « produit de l'énergie » : il \emph{convertit} l'énergie lumineuse en énergie électrique (conservation de l'énergie).
\end{itemize}
\end{attention}

\begin{aretenir}[L'essentiel de la section]
\begin{itemize}
\item L'effet photoélectrique est un \textbf{convertisseur lumière $\to$ électricité} : base de tous les détecteurs de lumière (barrières de sécurité, compteurs, codes-barres, télécommandes, capteurs photo).
\item \textbf{Effet externe} : électrons éjectés d'un métal dans le vide (cellule à vide, photomultiplicateur). \textbf{Effet interne} : électrons libérés dans un semi-conducteur (cellule photovoltaïque, photodiode, LDR).
\item Le \textbf{photomultiplicateur} amplifie en cascade (dynodes) : un photon $\to$ jusqu'à $10^6$ électrons.
\item Bilan énergétique d'un panneau : $P_{\text{élec}} = \eta \times E \times S$ ; l'ensoleillement camerounais ($4$ à $5{,}5\ \mathrm{kWh/m^2/jour}$) rend le solaire particulièrement pertinent pour l'électrification rurale.
\end{itemize}
\end{aretenir}

\begin{exercice}[Lampadaire solaire autonome]
Un lampadaire solaire à Yaoundé porte un panneau de $S = \SI{0,50}{m^2}$ de rendement $\eta = 18\,\%$, recevant en moyenne $\SI{4,5}{kWh/m^2/jour}$. Sa lampe LED consomme $\SI{30}{W}$ et fonctionne 12 h par nuit. \textbf{(1)} Calculer l'énergie électrique produite chaque jour. \textbf{(2)} La production couvre-t-elle les besoins de la lampe ? \textbf{(3)} Quel est le rôle de la batterie ?
\end{exercice}

\begin{corrige}[Lampadaire solaire autonome]
\textbf{(1)} $E_{\text{élec}} = \eta \times 4{,}5 \times 0{,}50 = 0{,}18 \times 2{,}25 \approx \SI{0,41}{kWh}$ par jour, soit environ $\SI{405}{Wh}$.
\textbf{(2)} Besoins : $30\ \mathrm{W} \times 12\ \mathrm{h} = \SI{360}{Wh}$. La production ($\approx \SI{405}{Wh}$) couvre les besoins, avec une marge modeste qui justifie un dimensionnement prudent (jours nuageux, pertes de stockage).
\textbf{(3)} La batterie stocke l'énergie produite le jour pour la restituer la nuit : elle décale dans le temps production et consommation, et lisse les variations d'ensoleillement.
\end{corrige}

\coursec{Quantité de mouvement du photon}

Dans la section précédente, nous avons vu que l'effet photoélectrique — utilisé dans les cellules photoélectriques et les panneaux photovoltaïques qui équipent de plus en plus de villages camerounais en électrification rurale — s'explique parfaitement si l'on admet que la lumière est constituée de grains d'énergie, les \cle{photons}, chacun transportant l'énergie $E = h\nu$. Mais si le photon se comporte comme une particule, une question surgit naturellement : possède-t-il, comme toute particule matérielle, une \cle{quantité de mouvement} ? C'est la question que nous allons traiter ici, et sa réponse ouvrira la porte à l'effet Compton, preuve éclatante de la nature corpusculaire du rayonnement.

\courssub{Dualité onde--corpuscule : un grain de lumière sans masse}

\textbf{Un constat troublant.}\par
D'un côté, les phénomènes d'interférences et de diffraction (franges d'Young, tache de Poisson) montrent que la lumière est une \emph{onde} caractérisée par une fréquence $\nu$ et une longueur d'onde $\lambda$. De l'autre, l'effet photoélectrique montre qu'elle échange son énergie par \emph{paquets} localisés : les photons. La lumière présente donc un \cle{double aspect}, ondulatoire et corpusculaire : c'est la \cle{dualité onde--corpuscule}.

\begin{definition}[Le photon]
Le \cle{photon} est le grain d'énergie associé à une onde électromagnétique de fréquence $\nu$. Il possède les caractéristiques suivantes :
\begin{itemize}
\item une énergie $E = h\nu = \dfrac{hc}{\lambda}$, où $h = \num{6,63e-34}\ \mathrm{J\,s}$ est la constante de Planck ;
\item une \textbf{masse nulle} : $m_{\text{photon}} = 0$ ;
\item une vitesse dans le vide égale à $c = \num{3,00e8}\ \mathrm{m\,s^{-1}}$.
\end{itemize}
\end{definition}

\begin{attention}[Le photon n'est pas une petite bille classique]
Le photon n'a \textbf{pas de masse}. Il est donc impossible de lui appliquer la mécanique de Newton : on ne peut pas écrire $E_c = \dfrac{1}{2}mv^2$ ni $\vec{p} = m\vec{v}$ pour un photon. Toutes les relations du photon ($E = h\nu$, $p = h/\lambda$) proviennent de la théorie de la relativité d'Einstein, et non de la mécanique classique.
\end{attention}

\courssub{Quantité de mouvement du photon : $p = \dfrac{h}{\lambda}$}

\textbf{Intuition.}\par
En mécanique classique, une particule de masse $m$ animée de la vitesse $\vec{v}$ possède une quantité de mouvement $\vec{p} = m\vec{v}$. Cette formule ne peut pas s'appliquer au photon, dont la masse est nulle : elle donnerait $p = 0$, ce qui signifierait que la lumière ne pourrait jamais pousser quoi que ce soit... Or on observe le contraire (nous y reviendrons avec la pression de radiation). Il faut donc une autre relation, fournie par la relativité restreinte.

En relativité, l'énergie totale d'une particule libre s'écrit $E^2 = p^2c^2 + m^2c^4$. Pour le photon, $m = 0$, donc $E = pc$, c'est-à-dire $p = \dfrac{E}{c}$. Comme $E = h\nu = \dfrac{hc}{\lambda}$, on obtient la relation fondamentale de cette section.

\begin{propriete}[Quantité de mouvement du photon — admise]
Tout photon de fréquence $\nu$ et de longueur d'onde $\lambda$ possède une quantité de mouvement de norme :
\[
p = \frac{E}{c} = \frac{h\nu}{c} = \frac{h}{\lambda}
\]
avec $p$ en $\mathrm{kg\,m\,s^{-1}}$, $h = \num{6,63e-34}\ \mathrm{J\,s}$ et $\lambda$ en mètres.
Cette relation est \textbf{admise} (elle découle de la relativité restreinte) ; sa démonstration n'est pas exigible au Baccalauréat.
\end{propriete}

\textbf{Vérification par analyse dimensionnelle.}\par
La constante de Planck $h$ s'exprime en $\mathrm{J\,s} = \mathrm{kg\,m^2\,s^{-1}}$. Donc :
\[
\left[\frac{h}{\lambda}\right] = \frac{\mathrm{kg\,m^2\,s^{-1}}}{\mathrm{m}} = \mathrm{kg\,m\,s^{-1}}
\]
qui est bien l'unité d'une quantité de mouvement. La relation est homogène.

\begin{definition}[Vecteur quantité de mouvement du photon]
Le vecteur quantité de mouvement $\vec{p}$ du photon est un vecteur :
\begin{itemize}
\item de \textbf{direction} : celle du rayon lumineux ;
\item de \textbf{sens} : celui de la propagation de la lumière ;
\item de \textbf{norme} : $p = \dfrac{h}{\lambda}$.
\end{itemize}
\end{definition}

\begin{methode}[Calculer la quantité de mouvement d'un photon]
\begin{enumerate}
\item Identifier la longueur d'onde $\lambda$ (ou la fréquence $\nu$) du photon dans l'énoncé.
\item Convertir $\lambda$ en \textbf{mètres} : $1\ \mathrm{nm} = 10^{-9}\ \mathrm{m}$, $1\ \mathrm{pm} = 10^{-12}\ \mathrm{m}$, $1\ \mathrm{\AA} = 10^{-10}\ \mathrm{m}$.
\item Appliquer $p = \dfrac{h}{\lambda}$ avec $h = \num{6,63e-34}\ \mathrm{J\,s}$.
\item Vérifier que le résultat est exprimé en $\mathrm{kg\,m\,s^{-1}}$ et contrôler l'ordre de grandeur (très inférieur à 1).
\item Si l'énoncé donne l'énergie $E$ du photon, on peut aussi utiliser directement $p = \dfrac{E}{c}$.
\end{enumerate}
\end{methode}

\begin{exemplebox}[Photon X utilisé en radiographie]
Un appareil de radiographie d'un hôpital de Yaoundé émet des photons X de longueur d'onde $\lambda = \num{0,07}\ \mathrm{nm}$. Calculons la quantité de mouvement d'un tel photon.
\begin{align*}
\lambda &= \num{0,07}\ \mathrm{nm} = \num{7e-11}\ \mathrm{m}\\
p &= \frac{h}{\lambda} = \frac{\num{6,63e-34}}{\num{7e-11}} \approx \num{9,5e-24}\ \mathrm{kg\,m\,s^{-1}}
\end{align*}
Cette valeur, bien que très faible à notre échelle, est \textbf{comparable} à la quantité de mouvement d'un électron quasi au repos : c'est pourquoi le choc photon--électron (effet Compton) modifie sensiblement le photon.
\end{exemplebox}

\courssub{Ordres de grandeur et conséquence : la pression de radiation}

\textbf{À quel point la quantité de mouvement d'un photon est-elle petite ?}\par
Pour un photon visible de longueur d'onde $\lambda \approx \num{0,5}\ \mathrm{\mu m}$ (lumière verte) :
\[
p = \frac{h}{\lambda} = \frac{\num{6,63e-34}}{\num{5e-7}} \approx \num{1,3e-27}\ \mathrm{kg\,m\,s^{-1}}
\]
soit de l'ordre de $10^{-27}\ \mathrm{kg\,m\,s^{-1}}$. À titre de comparaison, un grain de poussière de masse $10^{-10}\ \mathrm{kg}$ dérivant à $1\ \mathrm{mm\,s^{-1}}$ possède une quantité de mouvement de $10^{-13}\ \mathrm{kg\,m\,s^{-1}}$, soit $10^{14}$ fois plus grande !

\begin{aretenir}[Ordre de grandeur à connaître]
La quantité de mouvement d'un photon visible est de l'ordre de $10^{-27}\ \mathrm{kg\,m\,s^{-1}}$ : elle est \textbf{extrêmement faible} à l'échelle macroscopique, mais \textbf{mesurable et déterminante} dans les interactions individuelles photon--particule (effet Compton, diffusion de la lumière par les électrons).
\end{aretenir}

\textbf{Conséquence macroscopique : la pression de radiation.}\par
Pris isolément, le choc d'un photon est insignifiant. Mais un faisceau lumineux transporte un nombre gigantesque de photons par seconde. En frappant une surface, chaque photon lui cède sa quantité de mouvement (totalement s'il est absorbé, davantage s'il est réfléchi). D'après la deuxième loi de Newton généralisée ($\vec{F} = \dfrac{d\vec{p}}{dt}$), ce transfert global de quantité de mouvement se traduit par une \textbf{force} exercée par la lumière sur la surface, donc par une \textbf{pression} : c'est la \cle{pression de radiation}. Elle est très faible sur Terre (de l'ordre de quelques micropascals pour la lumière solaire), mais réelle et mesurable.

\begin{savaistu}[Les voiles solaires : naviguer grâce à la lumière]
La pression de radiation solaire est envisagée pour propulser des sondes spatiales équipées d'immenses \textbf{voiles solaires} réfléchissantes : les photons du Soleil, en rebondissant sur la voile, la poussent sans aucun carburant. La sonde japonaise \emph{IKAROS} (2010) a ainsi navigué vers Vénus, et le projet \emph{LightSail} a confirmé la faisabilité du concept. La lumière, sans masse, pousse réellement les objets : une conséquence directe de $p = \dfrac{h}{\lambda}$ !
\end{savaistu}

\courssub{Schéma : le photon incident et son vecteur quantité de mouvement}

\begin{popfigure}
\begin{center}
\begin{tikzpicture}[scale=1.1, >=Stealth]
% Onde incidente (sinusoïde)
\draw[popblue, thick, domain=0:3.2, samples=100] plot (\x, {0.35*sin(360*\x/1.2)});
% Vecteur p du photon
\draw[->, poporange, very thick] (1.6,0) -- (3.6,0) node[midway, above=2pt] {$\vec{p} = \dfrac{h}{\lambda}$};
% Photon (grain)
\fill[poporange] (3.6,0) circle (2.5pt);
\node[poporange, below] at (3.6,-0.15) {photon $(E = h\nu)$};
% Électron au repos
\fill[popink] (5.6,0) circle (4pt);
\node[popink, below] at (5.6,-0.25) {électron au repos};
\node[popink, above] at (5.6,0.25) {$e^-$};
% Longueur d'onde
\draw[<->, popgreen, thick] (0.3,-0.7) -- (1.5,-0.7) node[midway, below] {$\lambda$};
\draw[popgreen, dashed] (0.3,-0.35) -- (0.3,-0.75);
\draw[popgreen, dashed] (1.5,-0.35) -- (1.5,-0.75);
% Direction de propagation
\node[popmuted] at (1.6,0.75) {direction de propagation};
\end{tikzpicture}
\end{center}
\legende{Photon incident de longueur d'onde $\lambda$ se dirigeant vers un électron au repos. Le vecteur quantité de mouvement $\vec{p}$ a la direction et le sens du rayon lumineux, et pour norme $p = \dfrac{h}{\lambda}$.}
\end{popfigure}

\courssub{Vers l'effet Compton : le photon comme un véritable projectile}

Résumons la situation. Le photon possède :
\begin{itemize}
\item une énergie $E = h\nu = \dfrac{hc}{\lambda}$ ;
\item une quantité de mouvement $\vec{p}$ de norme $p = \dfrac{h}{\lambda}$, dirigée selon le rayon lumineux.
\end{itemize}
Il possède donc \emph{tous les attributs dynamiques d'une particule}. Dès lors, si un photon heurte un électron (par exemple un électron faiblement lié d'un atome, quasi libre et au repos), ce choc doit obéir aux deux grandes \textbf{lois de conservation} de la physique des collisions :
\begin{itemize}
\item conservation de l'\textbf{énergie totale} du système \{photon + électron\} ;
\item conservation de la \textbf{quantité de mouvement totale} (vectorielle) du système.
\end{itemize}
C'est exactement le raisonnement qu'Arthur Compton appliqua en 1923 à la diffusion des rayons X par les électrons de la matière. En traitant le photon comme un projectile de quantité de mouvement $\vec{p} = \dfrac{h}{\lambda}$, il montra que le photon diffusé ressort avec une longueur d'onde $\lambda'$ \emph{plus grande} que $\lambda$, selon la relation :
\[
\lambda' - \lambda = \frac{h}{m_e c}\,(1 - \cos\theta)
\]
que nous établirons et exploiterons dans la prochaine section. L'accord parfait entre cette formule et l'expérience constitue la preuve définitive que le photon porte bien une quantité de mouvement.

\begin{exoresolu}[Quantité de mouvement et énergie d'un photon]
Énoncé. Une station de télécommunications utilise un laser émettant des photons de longueur d'onde $\lambda = \num{1,55}\ \mathrm{\mu m}$ (fenêtre des fibres optiques).
\begin{enumerate}
\item Calculer la quantité de mouvement $p$ d'un photon.
\item En déduire son énergie $E$ en joules, puis en électronvolts.
\end{enumerate}
Données : $h = \num{6,63e-34}\ \mathrm{J\,s}$ ; $c = \num{3,00e8}\ \mathrm{m\,s^{-1}}$ ; $1\ \mathrm{eV} = \num{1,60e-19}\ \mathrm{J}$.
\tcblower
Solution détaillée.
\begin{enumerate}
\item Conversion : $\lambda = \num{1,55}\ \mathrm{\mu m} = \num{1,55e-6}\ \mathrm{m}$.
\[
p = \frac{h}{\lambda} = \frac{\num{6,63e-34}}{\num{1,55e-6}} \approx \num{4,28e-28}\ \mathrm{kg\,m\,s^{-1}}
\]
\item On utilise $E = pc$ :
\[
E = pc = \num{4,28e-28} \times \num{3,00e8} \approx \num{1,28e-19}\ \mathrm{J}
\]
Conversion en électronvolts :
\[
E = \frac{\num{1,28e-19}}{\num{1,60e-19}} \approx \num{0,80}\ \mathrm{eV}
\]
Vérification : on retrouve le même résultat avec $E = \dfrac{hc}{\lambda}$, ce qui confirme la cohérence des relations $E = pc$ et $p = \dfrac{h}{\lambda}$.
\end{enumerate}
\end{exoresolu}

\begin{attention}[Pièges fréquents au Baccalauréat]
\begin{itemize}
\item \textbf{Ne jamais écrire} $p = mv$ pour un photon : sa masse est nulle. Seule la relation $p = \dfrac{h}{\lambda}$ (ou $p = \dfrac{E}{c}$) est valable.
\item \textbf{Convertir systématiquement} $\lambda$ en mètres avant le calcul : oublier de convertir des nanomètres conduit à une erreur d'un facteur $10^9$ !
\item Ne pas confondre $p = \dfrac{h}{\lambda}$ (photon) avec la relation de de Broglie $\lambda = \dfrac{h}{p}$ pour une particule matérielle : c'est la même écriture, mais le sens physique diffère.
\item L'unité de $p$ est le $\mathrm{kg\,m\,s^{-1}}$ ; vérifier l'homogénéité avant de conclure.
\end{itemize}
\end{attention}

\coursec{Effet Compton}

Dans la section précédente, nous avons vu que le photon, grain de lumière d'énergie $E = h\nu$, transporte aussi une \cle{quantité de mouvement} $p = \dfrac{h\nu}{c} = \dfrac{h}{\lambda}$. Cette propriété, qui peut paraître abstraite, a été mise en évidence de façon spectaculaire en 1923 par le physicien américain Arthur Holly Compton : en bombardant de la matière avec des rayons X, il a observé que les photons « rebondissent » sur les électrons \emph{comme des boules de billard}, en leur cédant une partie de leur énergie et de leur quantité de mouvement. Cette expérience, qui lui valut le prix Nobel en 1927, constitue la preuve définitive du caractère corpusculaire de la lumière. Étudions-la en détail.

\courssub{Expérience de Compton (1923)}

\begin{experience}[L'expérience historique de Compton]
\textbf{Dispositif.} Un faisceau monochromatique de rayons X de longueur d'onde $\lambda$ bien connue (Compton utilisa la raie $K_\alpha$ du molybdène, $\lambda = \SI{0,071}{nm}$) est dirigé sur une cible de \cle{graphite} (carbone). Les électrons des couches externes du carbone sont \emph{faiblement liés} : leur énergie de liaison (quelques eV) est négligeable devant l'énergie des photons X (environ $\SI{17}{keV}$), si bien qu'on peut les considérer comme des électrons \textbf{libres et au repos}. On analyse, à l'aide d'un spectromètre à cristal, le rayonnement diffusé dans une direction faisant un angle $\theta$ avec la direction incidente.

\textbf{Observations.} Le rayonnement diffusé contient :
\begin{itemize}
\item une raie de longueur d'onde $\lambda$ \emph{inchangée} (diffusion sur les électrons fortement liés, où c'est l'atome entier qui recule) ;
\item une raie \emph{décalée} de longueur d'onde $\lambda' > \lambda$, dont le décalage $\Delta\lambda = \lambda' - \lambda$ \textbf{augmente quand l'angle $\theta$ augmente}, et ne dépend ni du matériau ni de $\lambda$.
\end{itemize}

\textbf{Interprétation.} Aucune théorie ondulatoire classique ne peut expliquer ce décalage : une onde électromagnétique diffusée conserve sa fréquence. Compton l'interprète comme un \emph{choc élastique} entre un photon et un électron au repos : le photon cède une partie de son énergie à l'électron, donc sa fréquence diminue ($\nu' < \nu$) et sa longueur d'onde augmente ($\lambda' > \lambda$).
\end{experience}

\begin{definition}[Effet Compton]
L'\cle{effet Compton} est la \textbf{diffusion d'un photon par un électron libre (ou faiblement lié) au repos}, au cours de laquelle le photon cède une partie de son énergie à l'électron : le photon diffusé a une fréquence plus faible, donc une longueur d'onde $\lambda'$ \textbf{plus grande} que la longueur d'onde incidente $\lambda$. L'électron, lui, est éjecté avec une énergie cinétique de recul.
\end{definition}

\courssub{Le choc photon--électron : bilans de conservation}

Modélisons la collision comme un choc élastique relativiste entre deux particules : le photon incident (énergie $h\nu$, quantité de mouvement $\vec{p}$ de norme $p = \dfrac{h}{\lambda}$) et l'électron supposé au repos (énergie au repos $m_ec^2$, quantité de mouvement nulle). Après le choc, le photon part dans une direction faisant l'angle $\theta$ avec la direction incidente, avec l'énergie $h\nu'$ et la quantité de mouvement $\vec{p}\,'$ ($p' = \dfrac{h}{\lambda'}$), tandis que l'électron de recul part avec l'angle $\varphi$ et la quantité de mouvement $\vec{p}_e$.

\begin{popfigure}
\begin{center}
\begin{tikzpicture}[scale=1.0, >=Stealth]
% Avant le choc
\draw[popmuted, dashed] (-3.2,0) -- (5.2,0);
\draw[->, very thick, poporange, decorate, decoration={snake, amplitude=0.8pt, segment length=6pt}] (-3,0) -- (-0.25,0);
\node[poporange, above] at (-1.8,0.15) {$\vec{p}$, $\lambda$, $h\nu$};
\node[popmuted, left] at (-3.2,0) {photon incident};
\fill[popblue] (0,0) circle (0.09);
\node[popblue, below left=1pt and -2pt] at (0,0) {$e^-$ au repos};
% Après le choc : photon diffusé
\draw[->, very thick, poporange, decorate, decoration={snake, amplitude=0.8pt, segment length=6pt}] (0.25,0.12) -- (3.4,1.9);
\node[poporange, above right] at (3.0,1.55) {$\vec{p}\,'$, $\lambda' > \lambda$, $h\nu'$};
% Électron de recul
\draw[->, very thick, popblue] (0.15,-0.12) -- (2.9,-1.55);
\fill[popblue] (2.9,-1.55) circle (0.09);
\node[popblue, below right] at (2.7,-1.45) {$e^-$ de recul, $\vec{p}_e$, $E_c$};
% Angles
\draw[->, popdark] (1.5,0) arc (0:29:1.5);
\node[popdark] at (1.85,0.32) {$\theta$};
\draw[->, popdark] (1.1,0) arc (0:-28:1.1);
\node[popdark] at (1.45,-0.33) {$\varphi$};
% Triangle des vecteurs p (encart)
\begin{scope}[shift={(5.6,1.9)}, scale=0.85]
\draw[->, very thick, poporange] (0,0) -- (2.6,0) node[midway, below] {$\vec{p}$};
\draw[->, very thick, poporange!70!popdark] (0,0) -- (1.7,1.15) node[midway, above left] {$\vec{p}\,'$};
\draw[->, very thick, popblue] (1.7,1.15) -- (2.6,0) node[midway, right] {$\vec{p}_e$};
\draw[popdark] (0.75,0) arc (0:34:0.75);
\node[popdark] at (1.0,0.22) {$\theta$};
\node[popdark, below] at (1.3,-0.55) {$\vec{p} = \vec{p}\,' + \vec{p}_e$};
\end{scope}
\end{tikzpicture}
\end{center}
\legende{Choc de Compton : le photon incident cède de l'énergie et de la quantité de mouvement à l'électron. À droite, le triangle des quantités de mouvement traduisant $\vec{p} = \vec{p}\,' + \vec{p}_e$.}
\end{popfigure}

\textbf{Conservation de l'énergie.} L'énergie totale avant le choc est la somme de l'énergie du photon et de l'énergie au repos de l'électron ; après le choc, elle est la somme de l'énergie du photon diffusé, de l'énergie au repos et de l'énergie cinétique de recul de l'électron :
\[
h\nu + m_ec^2 = h\nu' + m_ec^2 + E_c
\quad\Longrightarrow\quad
\boxed{E_c = h\nu - h\nu'}
\]
Le terme $m_ec^2$ se simplifie : l'énergie cinétique de l'électron de recul est simplement la différence entre les énergies des photons incident et diffusé.

\textbf{Conservation de la quantité de mouvement.} Vectoriellement :
\[
\vec{p} = \vec{p}\,' + \vec{p}_e
\quad\Longrightarrow\quad
\vec{p}_e = \vec{p} - \vec{p}\,'
\]
En projetant sur la direction incidente et sur la direction perpendiculaire, on obtient deux équations scalaires. En combinant ces deux équations avec la conservation de l'énergie écrite sous forme relativiste, et après élimination de $\varphi$ et de $p_e$ (calcul admis, non exigible au Baccalauréat), on aboutit à la célèbre relation de Compton.

\begin{propriete}[Relation de Compton (admise)]
Lors de la diffusion d'un photon de longueur d'onde $\lambda$ par un électron libre au repos, la longueur d'onde $\lambda'$ du photon diffusé sous l'angle $\theta$ vérifie :
\[
\boxed{\lambda' - \lambda = \frac{h}{m_e c}\,(1 - \cos\theta)}
\]
où $h = \SI{6,63e-34}{J.s}$ (constante de Planck), $m_e = \SI{9,11e-31}{kg}$ (masse de l'électron) et $c = \SI{3,00e8}{m/s}$. La démonstration complète (combinaison des deux conservations en mécanique relativiste) n'est \textbf{pas exigible} au Baccalauréat ; il faut en revanche savoir \emph{exploiter} la relation.
\end{propriete}

\begin{aretenir}[Longueur d'onde de Compton de l'électron]
La constante $\lambda_C = \dfrac{h}{m_e c}$ est appelée \cle{longueur d'onde de Compton} de l'électron :
\[
\lambda_C = \frac{\SI{6,63e-34}{J.s}}{\SI{9,11e-31}{kg} \times \SI{3,00e8}{m/s}} \approx \SI{2,43e-12}{m} = \SI{0,00243}{nm}
\]
Vérification dimensionnelle : $\left[\dfrac{h}{m c}\right] = \dfrac{\mathrm{kg\,m^2\,s^{-1}}}{\mathrm{kg}\times \mathrm{m\,s^{-1}}} = \mathrm{m}$. C'est bien une longueur.
\end{aretenir}

\courssub{Étude des cas limites et courbe $\Delta\lambda(\theta)$}

La relation de Compton $\Delta\lambda = \lambda_C(1-\cos\theta)$ permet d'analyser trois directions remarquables :

\begin{itemize}
\item $\theta = 0$ : $\cos\theta = 1$, donc $\Delta\lambda = 0$. Le photon continue tout droit : il n'y a pas eu d'interaction effective, pas de transfert d'énergie.
\item $\theta = 90°$ : $\cos\theta = 0$, donc $\Delta\lambda = \lambda_C \approx \SI{2,43e-12}{m}$. C'est le cas de l'exemple historique de Compton.
\item $\theta = 180°$ : $\cos\theta = -1$, donc $\Delta\lambda = 2\lambda_C \approx \SI{4,85e-12}{m}$. C'est la \textbf{rétrodiffusion} : le photon « rebondit » vers l'arrière et le décalage est \emph{maximal}. L'électron emporte alors l'énergie cinétique maximale.
\end{itemize}

\begin{popfigure}
\begin{center}
\begin{tikzpicture}
\begin{axis}[
width=0.85\linewidth, height=6.2cm,
xlabel={$\theta$ (degrés)}, ylabel={$\Delta\lambda$ (pm)},
xmin=0, xmax=180, ymin=0, ymax=5.4,
xtick={0,30,60,90,120,150,180},
ytick={0,1.215,2.43,3.645,4.86},
yticklabels={0, $\lambda_C/2$, $\lambda_C$, $3\lambda_C/2$, $2\lambda_C$},
grid=both, grid style={popline!40},
axis line style={popdark},
label style={popdark},
]
\addplot[popcurve, domain=0:180, samples=200] {2.43*(1-cos(x))};
\addplot[dashed, popmuted, domain=0:180] {4.86};
\node[popmuted, anchor=south west] at (axis cs:5,4.86) {\small $\Delta\lambda_{\max} = 2\lambda_C$};
\addplot[only marks, mark=*, poporange, mark size=2pt] coordinates {(90,2.43) (180,4.86)};
\node[poporange, anchor=south] at (axis cs:90,2.43) {\small $\theta=90°$};
\node[poporange, anchor=south east] at (axis cs:178,4.86) {\small $\theta=180°$};
\end{axis}
\end{tikzpicture}
\end{center}
\legende{Décalage Compton $\Delta\lambda = \lambda_C(1-\cos\theta)$ en fonction de l'angle de diffusion $\theta$ (en picomètres, $1\ \mathrm{pm} = 10^{-12}\ \mathrm{m}$). Le décalage croît de $0$ à $2\lambda_C$.}
\end{popfigure}

\begin{attention}[Pièges fréquents]
\begin{itemize}
\item $\theta$ est l'angle de diffusion \textbf{du photon} (entre $\vec{p}$ et $\vec{p}\,'$), \emph{pas} l'angle $\varphi$ de recul de l'électron. Les confondre fausse tout le calcul.
\item Travaille en \textbf{mètres} : $\lambda_C = \SI{2,43e-12}{m} = \SI{2,43e-3}{nm}$. Une longueur d'onde X est souvent donnée en nanomètres : convertis avant d'additionner.
\item Le décalage $\Delta\lambda$ ne dépend \emph{que} de $\theta$ : il est indépendant de $\lambda$ et du matériau. C'est $\lambda'$ qui dépend de $\lambda$.
\item $\Delta\lambda$ est toujours \textbf{positif} : un photon ne peut pas gagner d'énergie en frappant un électron au repos.
\end{itemize}
\end{attention}

\courssub{Exploitation quantitative : méthode et exemple}

\begin{methode}[Exploiter la relation de Compton]
\begin{enumerate}
\item \textbf{Calculer la longueur d'onde diffusée} : $\lambda' = \lambda + \lambda_C(1-\cos\theta)$, avec toutes les longueurs dans la même unité.
\item \textbf{En déduire l'énergie du photon diffusé} : $E' = h\nu' = \dfrac{hc}{\lambda'}$.
\item \textbf{En déduire l'énergie cinétique de recul de l'électron} (complément Bac) :
\[
E_c = h\nu - h\nu' = hc\left(\frac{1}{\lambda} - \frac{1}{\lambda'}\right)
\]
\item Convertir en électronvolts si demandé : $1\ \mathrm{eV} = \SI{1,602e-19}{J}$.
\end{enumerate}
\end{methode}

\begin{exemplebox}[L'expérience historique chiffrée]
Pour la raie du molybdène $\lambda = \SI{0,071}{nm}$ diffusée à $\theta = 90°$ :
\[
\Delta\lambda = \lambda_C(1-\cos 90°) = \lambda_C = \SI{0,00243}{nm}
\quad\Longrightarrow\quad
\lambda' = 0,071 + 0,00243 \approx \SI{0,0734}{nm}
\]
C'est exactement le décalage mesuré par Compton en 1923 : l'accord entre la théorie corpusculaire et l'expérience est remarquable.
\end{exemplebox}

\begin{exoresolu}[Rétrodiffusion d'un photon X]
Un photon X de longueur d'onde $\lambda = \SI{0,050}{nm}$ frappe un électron libre au repos et est rétrodiffusé ($\theta = 180°$). Déterminer : 1) la longueur d'onde $\lambda'$ du photon diffusé ; 2) l'énergie cinétique de l'électron de recul, en joules puis en keV.

Données : $h = \SI{6,63e-34}{J.s}$ ; $c = \SI{3,00e8}{m/s}$ ; $m_e = \SI{9,11e-31}{kg}$ ; $1\ \mathrm{eV} = \SI{1,602e-19}{J}$.
\tcblower
\textbf{1) Longueur d'onde diffusée.} Pour $\theta = 180°$, $\cos\theta = -1$, donc $1 - \cos\theta = 2$ :
\[
\lambda' = \lambda + 2\lambda_C = \SI{0,050}{nm} + 2 \times \SI{0,00243}{nm} = \SI{0,0549}{nm}
\]

\textbf{2) Énergie cinétique de recul.} On applique la conservation de l'énergie avec les longueurs en mètres :
\begin{align*}
E_c &= hc\left(\frac{1}{\lambda} - \frac{1}{\lambda'}\right)
= \SI{6,63e-34}{} \times \SI{3,00e8}{} \times \left(\frac{1}{\SI{5,0e-11}{}} - \frac{1}{\SI{5,49e-11}{}}\right)\\
&= \SI{1,989e-25}{J.m} \times \left(\SI{2,00e10}{m^{-1}} - \SI{1,821e10}{m^{-1}}\right)\\
&= \SI{1,989e-25}{} \times \SI{1,79e9}{} \approx \SI{3,56e-16}{J}
\end{align*}
Conversion en keV :
\[
E_c = \frac{\SI{3,56e-16}{J}}{\SI{1,602e-19}{J/eV}} \approx \SI{2,22e3}{eV} \approx \SI{2,2}{keV}
\]
Vérification : le photon incident valait $E = \dfrac{hc}{\lambda} \approx \SI{3,98e-15}{J} \approx \SI{24,8}{keV}$ ; le photon diffusé emporte $E' \approx \SI{22,6}{keV}$. La différence, $\SI{2,2}{keV}$, est bien l'énergie de l'électron : le bilan est cohérent.
\end{exoresolu}

\courssub{Portée de l'effet Compton et applications}

L'effet Compton a une portée historique immense. L'effet photoélectrique (1905) avait montré que la lumière \emph{échange} son énergie par quanta ; l'effet Compton (1923) montre en plus que le photon possède une \emph{quantité de mouvement} $p = h/\lambda$ et se comporte, dans une collision, exactement comme une particule matérielle obéissant aux lois de conservation de la mécanique. C'est la preuve définitive du \cle{caractère corpusculaire de la lumière}, qui parachève la dualité onde--corpuscule : la lumière se propage comme une onde (interférences, diffraction) mais interagit avec la matière comme un corpuscule (effets photoélectrique et Compton).

\begin{savaistu}[L'effet Compton en imagerie médicale]
L'effet Compton n'est pas qu'une curiosité de laboratoire : c'est le \textbf{mécanisme dominant d'interaction des rayons X avec les tissus mous} du corps humain aux énergies utilisées en radiographie et en scanner. La probabilité de diffusion Compton dépend de la densité électronique des tissus : les os, riches en calcium, interagissent autrement que les muscles ou les graisses, ce qui crée le contraste des images radiographiques. En radiothérapie, pratiquée notamment dans les hôpitaux de référence de Yaoundé et Douala, l'énergie déposée dans la tumeur par les faisceaux de photons de haute énergie provient en grande partie des électrons de recul produits par effet Compton : ce sont eux qui ionisent et détruisent les cellules cancéreuses. Comprendre Compton, c'est comprendre la physique qui sauve des vies.
\end{savaistu}

\begin{exercice}[À toi de jouer]
Un photon de longueur d'onde $\lambda = \SI{0,030}{nm}$ est diffusé par un électron libre au repos sous un angle $\theta = 60°$.
\begin{enumerate}
\item Calculer le décalage de Compton $\Delta\lambda$ et la longueur d'onde $\lambda'$ du photon diffusé.
\item Calculer l'énergie cinétique de recul de l'électron en keV.
\item Sous quel angle faudrait-il observer la diffusion pour obtenir un décalage $\Delta\lambda = \dfrac{\lambda_C}{2}$ ?
\end{enumerate}
\end{exercice}

\begin{corrige}[Exercice 1]
\textbf{1)} Pour $\theta = 60°$, $\cos\theta = 0,5$, donc $1-\cos\theta = 0,5$ :
\[
\Delta\lambda = 0,5\,\lambda_C = \SI{1,215e-3}{nm} \quad\Longrightarrow\quad \lambda' = 0,030 + 0,001215 \approx \SI{0,0312}{nm}
\]
\textbf{2)} $E_c = hc\left(\dfrac{1}{\lambda}-\dfrac{1}{\lambda'}\right) = \SI{1,989e-25}{} \times \left(\dfrac{1}{\SI{3,0e-11}{}} - \dfrac{1}{\SI{3,1215e-11}{}}\right) \approx \SI{1,989e-25}{} \times \SI{1,29e9}{} \approx \SI{2,57e-16}{J} \approx \SI{1,6}{keV}$.

\textbf{3)} On veut $\lambda_C(1-\cos\theta) = \dfrac{\lambda_C}{2}$, soit $\cos\theta = \dfrac{1}{2}$, d'où $\theta = 60°$. (Cohérent avec la question 1.)
\end{corrige}

\newpage

\coursec{Activité d'intégration}

\courssub{Objectifs de l'activité}

À l'issue de cette activité d'intégration, tu seras capable de :
\begin{itemize}
\item tracer et exploiter la caractéristique $I = f(U)$ d'une cellule photoélectrique et d'en déduire le \cle{potentiel d'arrêt} et le \cle{courant de saturation} ;
\item vérifier expérimentalement la \cle{relation d'Einstein} $E_{c\max} = h\nu - W_0$ et la relation $eU_0 = E_{c\max}$ ;
\item calculer un \cle{rendement quantique} à partir du courant de saturation et du flux de photons incidents ;
\item exploiter la \cle{relation de Compton} $\lambda' - \lambda = \dfrac{h}{m_e c}(1 - \cos\theta)$ pour analyser un spectre de diffusion ;
\item rédiger une note technique argumentée, compétence clé du physicien et de l'ingénieur.
\end{itemize}

\courssub{Situation}

Le lycée bilingue de Maroua (région de l'Extrême-Nord) souhaite équiper sa salle informatique de lampes LED alimentées par des panneaux solaires photovoltaïques. Avant l'achat du matériel, le club scientifique du lycée est chargé de deux missions de contrôle qualité.

\textbf{Mission A : étalonnage d'une cellule photoélectrique.} Pour comprendre le principe de conversion lumière--électricité, les élèves étudient d'abord une cellule photoélectrique étalon dont la photocathode est en \cle{césium}, de travail d'extraction $W_0 = \SI{1,9}{eV}$. La cellule est éclairée par une lumière monochromatique de longueur d'onde $\lambda = \SI{500}{nm}$. Un flux de $N = \num{1,0e15}$ photons par seconde arrive sur la photocathode. Les élèves relèvent, point par point, l'intensité $I$ du courant en fonction de la tension $U$ entre anode et cathode :

\begin{center}
\begin{tabularx}{\linewidth}{|l|X|X|X|X|X|X|X|}\hline
\rowcolor{popblueL}
$U$ (V) & $-1,0$ & $-0,6$ & $-0,4$ & $-0,2$ & $0$ & $1$ & $3$ \\ \hline
$I$ (\si{\micro\ampere}) & $0$ & $0$ & $1,5$ & $4,0$ & $6,5$ & $7,6$ & $8,0$ \\ \hline
\end{tabularx}
\end{center}

Pour $U \geq \SI{3}{V}$, l'intensité reste constante : $I_{\text{sat}} = \SI{8,0}{\micro\ampere}$.

\textbf{Mission B : sécurité autour d'un appareil de radiographie.} L'hôpital régional de Maroua utilise un appareil à rayons X émettant des photons de longueur d'onde $\lambda = \SI{0,030}{nm}$. Lors de la traversée des tissus, certains photons subissent une \cle{diffusion Compton} sous un angle $\theta = 60^{\circ}$. Le club doit évaluer l'énergie cédée aux électrons des tissus afin de justifier les mesures de protection radiologique (tabliers de plomb, écrans).

\textbf{Données :} $h = \SI{6,63e-34}{J.s}$ ; $c = \SI{3,00e8}{m/s}$ ; $e = \SI{1,60e-19}{C}$ ; $m_e = \SI{9,11e-31}{kg}$ ; $\SI{1}{eV} = \SI{1,60e-19}{J}$ ; longueur d'onde de Compton $\lambda_C = \dfrac{h}{m_e c} = \SI{2,43e-12}{m}$.

\courssub{Consignes}

\textbf{Partie A — La cellule photoélectrique au césium}
\begin{enumerate}
\item Calculer l'énergie $E = h\nu$ d'un photon incident, en joules puis en électronvolts. Vérifier que l'effet photoélectrique est possible.
\item Tracer la caractéristique $I = f(U)$ à partir du relevé. En déduire graphiquement le potentiel d'arrêt $U_0$ et le courant de saturation $I_{\text{sat}}$.
\item Calculer l'énergie cinétique maximale $E_{c\max} = eU_0$ des photoélectrons, puis comparer à la valeur prédite par la relation d'Einstein $E_{c\max} = h\nu - W_0$. Conclure.
\item Calculer le nombre d'électrons émis par seconde à saturation, puis le rendement quantique $\eta$ de la cellule. Commenter.
\end{enumerate}

\textbf{Partie B — Diffusion Compton des rayons X}
\begin{enumerate}
\setcounter{enumi}{4}
\item Calculer la longueur d'onde $\lambda'$ du photon diffusé à $\theta = 60^{\circ}$.
\item En déduire l'énergie $E_c$ cédée à l'électron de recul, en joules puis en keV.
\item Comparer $E_c$ à l'énergie du photon incident et expliquer pourquoi une protection radiologique est indispensable.
\end{enumerate}

\textbf{Restitution :} rédiger une note technique d'une page, adressée au proviseur et au médecin-chef, présentant les résultats, les vérifications effectuées et les recommandations.

\courssub{Ressources à mobiliser}

\begin{aretenir}[Boîte à outils de l'activité]
\begin{itemize}
\item Énergie d'un photon : $E = h\nu = \dfrac{hc}{\lambda}$.
\item Relation d'Einstein : $E_{c\max} = h\nu - W_0$, valable seulement si $\nu \geq \nu_0 = \dfrac{W_0}{h}$.
\item Potentiel d'arrêt : $eU_0 = E_{c\max}$, où $U_0$ est la valeur absolue de la tension qui annule le courant.
\item Courant de saturation : $I_{\text{sat}} = n_e \, e$ où $n_e$ est le nombre d'électrons émis par seconde.
\item Rendement quantique : $\eta = \dfrac{n_e}{N}$ (électrons émis / photons incidents).
\item Relation de Compton : $\lambda' - \lambda = \dfrac{h}{m_e c}(1 - \cos\theta)$.
\item Conservation de l'énergie dans la diffusion : $E_c(\text{électron}) = h\nu - h\nu' = hc\left(\dfrac{1}{\lambda} - \dfrac{1}{\lambda'}\right)$.
\end{itemize}
\end{aretenir}

\begin{attention}[Pièges fréquents à éviter]
\begin{itemize}
\item Ne pas confondre les unités : convertis toujours les longueurs d'onde en \textbf{mètres} avant tout calcul ($\SI{500}{nm} = \SI{5,00e-7}{m}$ ; $\SI{0,030}{nm} = \SI{3,0e-11}{m}$).
\item Le potentiel d'arrêt se lit comme une tension \emph{négative} sur le graphique ($U = -U_0$) ; la relation $eU_0 = E_{c\max}$ utilise sa \textbf{valeur absolue}.
\item La relation d'Einstein ne s'applique que si $h\nu \geq W_0$ : vérifie cette condition \emph{avant} de calculer $E_{c\max}$.
\item Dans la relation de Compton, $\theta$ est l'angle de diffusion du \textbf{photon}, pas celui de l'électron de recul.
\end{itemize}
\end{attention}

\courssub{Démarche et corrigé commenté}

\begin{exoresolu}[Note technique — corrigé complet]
Énoncé : reprendre les huit consignes ci-dessus avec les données de la situation.
\tcblower
\textbf{Partie A.}

\textbf{1. Énergie du photon incident.}
\[ E = \frac{hc}{\lambda} = \frac{6,63\times 10^{-34}\times 3,00\times 10^{8}}{500\times 10^{-9}} = \SI{3,98e-19}{J} = \frac{3,98\times 10^{-19}}{1,60\times 10^{-19}} = \SI{2,49}{eV}. \]
Comme $E = \SI{2,49}{eV} > W_0 = \SI{1,9}{eV}$, l'extraction d'électrons est possible : l'effet photoélectrique a lieu. (À titre de vérification, la fréquence seuil du césium vaut $\nu_0 = \dfrac{W_0}{h} = \dfrac{1,9\times 1,60\times 10^{-19}}{6,63\times 10^{-34}} \approx \SI{4,6e14}{Hz}$, soit $\lambda_0 \approx \SI{650}{nm} > \SI{500}{nm}$ : la condition $\lambda \leq \lambda_0$ est bien satisfaite.)

\textbf{2. Tracé et exploitation de la caractéristique.} On porte $U$ en abscisse, $I$ en ordonnée. La courbe est nulle jusqu'à $U \approx \SI{-0,6}{V}$, croît fortement entre $-0,6$ et $0$ V, puis plafonne à $I_{\text{sat}} = \SI{8,0}{\micro\ampere}$ pour $U \geq \SI{3}{V}$. Le potentiel d'arrêt, tension pour laquelle le courant s'annule, se lit sur l'axe des abscisses : $U_0 \approx \SI{0,6}{V}$ (en valeur absolue).

\begin{popfigure}
\begin{center}
\begin{tikzpicture}
\begin{axis}[width=0.85\linewidth,height=6cm,
xlabel={$U$ (V)},ylabel={$I$ (\si{\micro\ampere})},
xmin=-1.2,xmax=3.5,ymin=-0.5,ymax=9,
grid=both,grid style={popline!40},
axis lines=left,thick]
\addplot[popcurve,mark=*,mark size=1.5pt] coordinates {(-1,0) (-0.6,0) (-0.4,1.5) (-0.2,4) (0,6.5) (1,7.6) (3,8)};
\addplot[popcurve,domain=3:3.4,samples=2]{8};
\draw[dashed,poporange] (axis cs:-0.6,0) -- (axis cs:-0.6,8);
\node[poporange,anchor=south,font=\small] at (axis cs:-0.6,8) {$-U_0$};
\draw[dashed,popgreen] (axis cs:-1.2,8) -- (axis cs:3.4,8);
\node[popgreen,anchor=south west,font=\small] at (axis cs:1.2,8.1) {$I_{\text{sat}}$};
\end{axis}
\end{tikzpicture}
\end{center}
\legende{Caractéristique $I = f(U)$ de la cellule au césium éclairée à $\lambda = \SI{500}{nm}$ : lecture du potentiel d'arrêt $U_0 \approx \SI{0,6}{V}$ et du courant de saturation $I_{\text{sat}} = \SI{8,0}{\micro\ampere}$.}
\end{popfigure}

\textbf{3. Vérification de la relation d'Einstein.}
\[ E_{c\max} = eU_0 = 1,60\times 10^{-19}\times 0,6 = \SI{9,6e-20}{J} \approx \SI{0,60}{eV}. \]
Prédiction d'Einstein : $E_{c\max} = h\nu - W_0 = 2,49 - 1,9 = \SI{0,59}{eV}$. Les deux valeurs coïncident à la précision de la lecture graphique près : la relation d'Einstein est vérifiée expérimentalement. Cet accord conforte le modèle corpusculaire de la lumière : chaque photon cède \emph{intégralement} son énergie $h\nu$ à un seul électron.

\textbf{4. Rendement quantique.} À saturation, chaque électron émis est collecté par l'anode :
\[ n_e = \frac{I_{\text{sat}}}{e} = \frac{8,0\times 10^{-6}}{1,60\times 10^{-19}} = \SI{5,0e13}{\text{électrons/s}}. \]
\[ \eta = \frac{n_e}{N} = \frac{5,0\times 10^{13}}{1,0\times 10^{15}} = 0,050 = \SI{5,0}{\%}. \]
Seul un photon sur vingt arrache un électron : c'est typique d'une photocathode réelle, et cela justifie la recherche de matériaux performants pour les panneaux solaires du lycée.

\textbf{Partie B.}

\textbf{5. Longueur d'onde diffusée.}
\[ \lambda' = \lambda + \lambda_C (1 - \cos 60^{\circ}) = 3,00\times 10^{-11} + 2,43\times 10^{-12}\times (1 - 0,5) = \SI{3,12e-11}{m} = \SI{0,0312}{nm}. \]
Le photon diffusé a une longueur d'onde \emph{plus grande} que le photon incident : il a perdu de l'énergie, preuve que le choc photon--électron se comporte comme un choc entre deux particules.

\textbf{6. Énergie cédée à l'électron.}
\begin{align*}
E_c &= hc\left(\frac{1}{\lambda} - \frac{1}{\lambda'}\right) = 6,63\times 10^{-34}\times 3,00\times 10^{8}\left(\frac{1}{3,00\times 10^{-11}} - \frac{1}{3,12\times 10^{-11}}\right)\\
&= 1,989\times 10^{-25}\times\left(3,333\times 10^{10} - 3,205\times 10^{10}\right) = \SI{2,55e-16}{J} = \SI{1,6}{keV}.
\end{align*}

\textbf{7. Commentaire.} Le photon incident vaut $E = \dfrac{hc}{\lambda} = \dfrac{1,989\times 10^{-25}}{3,00\times 10^{-11}} = \SI{6,6e-15}{J} \approx \SI{41}{keV}$ ; l'électron de recul emporte environ $4\,\%$ de cette énergie à $\theta = 60^{\circ}$ (la fraction croît avec l'angle et devient maximale en rétrodiffusion, $\theta = 180^{\circ}$). Ces électrons énergétiques ionisent les tissus et endommagent l'ADN des cellules : d'où les tabliers de plomb, les écrans et la limitation des expositions. La diffusion Compton est d'ailleurs le mode dominant d'interaction des rayons X de radiographie avec la matière.
\end{exoresolu}

\begin{methode}[Exploiter une caractéristique $I = f(U)$ en quatre étapes]
\begin{enumerate}
\item \textbf{Tracer} la courbe point par point, $U$ en abscisse (tensions négatives à gauche), $I$ en ordonnée.
\item \textbf{Repérer le potentiel d'arrêt} : abscisse du point où la courbe rejoint l'axe ($I = 0$) ; prendre sa valeur absolue $U_0$.
\item \textbf{Repérer le palier de saturation} : ordonnée $I_{\text{sat}}$ de la partie horizontale.
\item \textbf{Exploiter} : $E_{c\max} = eU_0$ à comparer à $h\nu - W_0$ ; $n_e = \dfrac{I_{\text{sat}}}{e}$ puis $\eta = \dfrac{n_e}{N}$.
\end{enumerate}
\end{methode}

\courssub{Bilan de l'activité}

Cette activité a permis de mobiliser l'ensemble des savoir-faire de la leçon dans un cadre authentique : le tracé et l'exploitation de la caractéristique $I = f(U)$ ont donné un sens concret au potentiel d'arrêt et au courant de saturation ; la vérification de la relation d'Einstein a montré qu'une théorie se valide par la mesure ; le calcul du rendement quantique a relié le monde microscopique (photons, électrons) à une grandeur macroscopique mesurable (le courant) ; enfin la relation de Compton a confirmé la nature corpusculaire de la lumière et fondé scientifiquement les règles de radioprotection. Tu as ainsi agi comme un véritable technicien-physicien : mesurer, modéliser, vérifier, puis conseiller.

\newpage

\coursec{Méthodes et exercices résolus}

\coursec{Mise en évidence de l'effet photoélectrique}

\begin{experience}[Expérience de Hertz (1887) : découverte de l'effet photoélectrique]
\textbf{Matériel} : une source d'étincelles (bobine de Ruhmkorff), un récepteur constitué d'une boucle métallique avec un petit intervalle, une lampe à arc (source UV) ou un prisme pour filtrer la lumière.
\textbf{Protocole} : On produit des étincelles dans l'émetteur. On observe l'apparition d'étincelles dans la boucle réceptrice. On interpose ensuite une plaque de verre (qui absorbe les UV) ou un prisme pour ne laisser passer que la lumière visible.
\textbf{Observations} :
\begin{itemize}
\item Les étincelles dans le récepteur sont plus faciles à obtenir lorsque la lumière ultraviolette (UV) de l'émetteur atteint la boucle réceptrice.
\item Si l'on place une plaque de verre (opaque aux UV) entre l'émetteur et le récepteur, les étincelles cessent ou deviennent très difficiles à observer.
\item La lumière visible seule (filtrée par un prisme) ne provoque pas l'effet.
\end{itemize}
\textbf{Interprétation} : La lumière UV arrache des électrons de la surface métallique de la boucle réceptrice (cathode), ce qui ionise l'air dans l'intervalle et facilite le passage de l'étincelle. L'énergie de la lumière est transférée aux électrons du métal. L'expérience montre qu'il existe une \textbf{condition de fréquence} (seuil) : les UV fonctionnent, le visible non.
\end{experience}

\begin{definition}[Cellule photoélectrique (tube à vide)]
Dispositif constitué d'un bulbe sous vide contenant deux électrodes métalliques :
\begin{itemize}
\item une \textbf{cathode} (émetteur), sensible à la lumière, maintenue à un potentiel $U_C$ ;
\item une \textbf{anode} (collecteur), maintenue à un potentiel $U_A$.
\end{itemize}
La tension appliquée est $U = U_A - U_C$. On éclaire la cathode par une radiation monochromatique de fréquence $\nu$ et de puissance $P$. On mesure le courant $I$ traversant le circuit en fonction de $U$.
\end{definition}

\begin{savaistu}[Contexte camerounais : l'énergie solaire]
Au Cameroun, l'électrification rurale repose en partie sur les \textbf{panneaux photovoltaïques} (effet photovoltaïque, cousin de l'effet photoélectrique). Contrairement aux barrages de \textbf{Lom Pangar}, \textbf{Nachtigal} ou \textbf{Edéa} (hydroélectricité gérée par \textbf{ENEO}/\textbf{SONATREL}), le solaire permet une production décentralisée. Les cellules en silicium absorbent les photons du soleil (spectre visible et proche IR) et génèrent un courant continu sans pièce mobile, idéal pour les centres de santé isolés, les pompages d'eau agricoles ou les relais de télécommunications \textbf{MTN}/\textbf{Orange}/\textbf{CAMTEL} en zone rurale.
\end{savaistu}

\coursec{Seuil photoélectrique et relation d'Einstein}

\begin{definition}[Travail d'extraction $W_0$, fréquence seuil $\nu_0$, longueur d'onde seuil $\lambda_0$]
Pour qu'un électron soit extrait du métal, il doit recevoir une énergie au moins égale au \textbf{travail d'extraction} $W_0$ (énergie de liaison de l'électron le moins lié dans le métal, aussi appelé \emph{fonction de travail}).
\begin{itemize}
\item Fréquence seuil : $\nu_0 = \dfrac{W_0}{h}$.
\item Longueur d'onde seuil : $\lambda_0 = \dfrac{c}{\nu_0} = \dfrac{hc}{W_0}$.
\end{itemize}
\cle{Condition d'émission} : L'effet photoélectrique se produit \textbf{si et seulement si} $\nu \geq \nu_0$ (équivalent à $\lambda \leq \lambda_0$).
\end{definition}

\begin{propriete}[Relation d'Einstein (1905) — \textbf{Exigible au Bac}]
L'énergie cinétique maximale des photoélectrons émis est donnée par :
\[ E_{c\max} = h\nu - W_0 = \frac{hc}{\lambda} - W_0. \]
\begin{itemize}
\item $h = \num{6,63e-34}\ \text{J.s}$ (constante de Planck).
\item $W_0$ en joules (ou en eV : $1\ \text{eV} = \num{1,60e-19}\ \text{J}$).
\item Si $h\nu < W_0$, alors $E_{c\max} < 0$ : \textbf{aucune émission}, quelle que soit l'intensité lumineuse.
\end{itemize}
\cle{Analyse dimensionnelle} : $[h\nu] = \text{J.s} \times \text{s}^{-1} = \text{J}$ (homogène à une énergie).
\end{propriete}

\begin{methode}[Seuil photoélectrique et travail d'extraction]
Pour déterminer si l'effet photoélectrique se produit et exploiter le seuil :
\begin{enumerate}
\item \textbf{Convertir les données en unités cohérentes} : longueur d'onde en mètres, énergie en joules (ou en eV avec $1\ \text{eV} = \num{1,60e-19}\ \text{J}$).
\item \textbf{Calculer l'énergie du photon incident} : $E = h\nu = \dfrac{hc}{\lambda}$.
\item \textbf{Comparer au travail d'extraction} $W_0 = h\nu_0 = \dfrac{hc}{\lambda_0}$ :
\begin{itemize}
\item si $E \geq W_0$ (c'est-à-dire $\lambda \leq \lambda_0$) : émission de photoélectrons ;
\item si $E < W_0$ (c'est-à-dire $\lambda > \lambda_0$) : aucune émission, quelle que soit l'intensité lumineuse.
\end{itemize}
\item \textbf{Conclure par une phrase} précisant la condition sur $\nu$ ou $\lambda$.
\end{enumerate}
\textbf{Réflexe} : la condition $\lambda \leq \lambda_0$ va dans le sens \emph{inverse} de la condition sur les fréquences, car $\lambda = c/\nu$.
\end{methode}

\begin{exoresolu}[Seuil photoélectrique du zinc]
Une plaque de zinc est éclairée par une radiation de longueur d'onde $\lambda = \num{0,25e-6}\ \text{m}$. Le travail d'extraction du zinc est $W_0 = \num{4,3}\ \text{eV}$.
\begin{enumerate}
\item Calculer la longueur d'onde seuil $\lambda_0$ du zinc.
\item L'effet photoélectrique se produit-il ? Justifier.
\end{enumerate}
On donne : $h = \num{6,63e-34}\ \text{J.s}$, $c = \num{3,00e8}\ \text{m/s}$, $1\ \text{eV} = \num{1,60e-19}\ \text{J}$.
\tcblower
\textbf{1. Longueur d'onde seuil.}

Par définition, $W_0 = h\nu_0 = \dfrac{hc}{\lambda_0}$, donc :
\[ \lambda_0 = \frac{hc}{W_0}. \]
Conversion : $W_0 = 4{,}3 \times \num{1,60e-19} = \num{6,88e-19}\ \text{J}$.
\[ \lambda_0 = \frac{\num{6,63e-34} \times \num{3,00e8}}{\num{6,88e-19}} = \num{2,89e-7}\ \text{m} = \num{0,289e-6}\ \text{m} = \num{289}\ \text{nm}. \]

\textbf{2. Existence de l'effet photoélectrique.}

On compare : $\lambda = \num{0,25e-6}\ \text{m} = \num{250}\ \text{nm} < \lambda_0 = \num{289}\ \text{nm}$.

La radiation incidente est plus énergétique que le seuil : \textbf{l'effet photoélectrique se produit}. (Vérification : $E = hc/\lambda = \num{7,96e-19}\ \text{J} = \num{4,97}\ \text{eV} > W_0 = \num{4,3}\ \text{eV}$.)
\end{exoresolu}

\begin{methode}[Énergie cinétique maximale des photoélectrons]
\begin{enumerate}
\item Écrire la \textbf{relation d'Einstein} : $E_{c\max} = h\nu - W_0 = \dfrac{hc}{\lambda} - W_0$.
\item Mettre toutes les énergies dans la \textbf{même unité} (joules de préférence pour les calculs, eV pour la présentation).
\item Calculer $E_{c\max}$ ; vérifier qu'elle est \textbf{positive} (sinon pas d'émission).
\item Si la vitesse maximale est demandée : $E_{c\max} = \dfrac{1}{2}m_e v_{\max}^2$ d'où $v_{\max} = \sqrt{\dfrac{2E_{c\max}}{m_e}}$.
\end{enumerate}
\textbf{Analyse dimensionnelle} : $[h\nu] = \text{J.s} \times \text{s}^{-1} = \text{J}$ : homogène à une énergie.
\end{methode}

\begin{exoresolu}[Photoélectrons émis par le césium]
La cathode d'une cellule photoélectrique est en césium, de travail d'extraction $W_0 = \num{1,9}\ \text{eV}$. Elle est éclairée par une lumière monochromatique de longueur d'onde $\lambda = \num{450}\ \text{nm}$.
\begin{enumerate}
\item Calculer l'énergie cinétique maximale des photoélectrons émis, en joules puis en eV.
\item En déduire leur vitesse maximale.
\end{enumerate}
On donne : $h = \num{6,63e-34}\ \text{J.s}$, $c = \num{3,00e8}\ \text{m/s}$, $m_e = \num{9,11e-31}\ \text{kg}$, $1\ \text{eV} = \num{1,60e-19}\ \text{J}$.
\tcblower
\textbf{1. Énergie cinétique maximale.}

Énergie du photon incident :
\[ E = \frac{hc}{\lambda} = \frac{\num{6,63e-34} \times \num{3,00e8}}{\num{450e-9}} = \num{4,42e-19}\ \text{J} = \frac{\num{4,42e-19}}{\num{1,60e-19}} = \num{2,76}\ \text{eV}. \]
Comme $E = \num{2,76}\ \text{eV} > W_0 = \num{1,9}\ \text{eV}$, il y a émission. Relation d'Einstein :
\[ E_{c\max} = h\nu - W_0 = 2{,}76 - 1{,}9 = \num{0,86}\ \text{eV} = 0{,}86 \times \num{1,60e-19} = \num{1,38e-19}\ \text{J}. \]

\textbf{2. Vitesse maximale.}
\[ v_{\max} = \sqrt{\frac{2E_{c\max}}{m_e}} = \sqrt{\frac{2 \times \num{1,38e-19}}{\num{9,11e-31}}} = \sqrt{\num{3,03e11}} = \num{5,5e5}\ \text{m/s}. \]
Les photoélectrons sont émis avec une vitesse maximale d'environ $\num{5,5e5}\ \text{m/s}$, très inférieure à $c$ : le traitement non relativiste est légitime.
\end{exoresolu}

\coursec{Caractéristique $I = f(U)$ d'une cellule photoélectrique}

\begin{experience}[Caractéristique courant-tension $I = f(U)$]
\textbf{Matériel} : Cellule photoélectrique (cathode alcaline, anode annulaire), source lumineuse monochromatique (laser ou lampe + filtre interférentiel) de puissance variable, générateur de tension variable $U$ (polarité directe et inverse), ampèremètre sensible (pico/ nano-ampèremètre).
\textbf{Protocole} : On éclaire la cathode à fréquence $\nu$ fixée. On fait varier la tension $U = U_A - U_C$ de valeurs négatives (freinage) à valeurs positives (accélération). On relève le courant $I$.
\textbf{Observations} (courbe typique) :
\begin{itemize}
\item Pour $U \geq 0$ (tension accélératrice) : le courant augmente rapidement puis atteint un palier \textbf{courant de saturation $I_S$}. Tous les électrons émis sont collectés.
\item Pour $U < 0$ (tension freinante) : le courant diminue et s'annule pour $U = -U_0$. $U_0 > 0$ est le \textbf{potentiel d'arrêt}.
\item Si l'on augmente l'intensité lumineuse (puissance $P$) à $\nu$ fixée : $I_S$ augmente proportionnellement, $U_0$ \textbf{ne change pas}.
\item Si l'on augmente $\nu$ à $P$ fixée : $U_0$ augmente, $I_S$ \textbf{ne change pas} (nombre de photons $N = P/(h\nu)$ diminue si $P$ constant, mais rendement quantique constant $\Rightarrow I_S$ constant).
\end{itemize}
\end{experience}

\begin{definition}[Potentiel d'arrêt $U_0$, Courant de saturation $I_S$, Rendement quantique $\eta$]
\begin{itemize}
\item \textbf{Potentiel d'arrêt $U_0$} : Tension inverse (valeur positive) qui stoppe les photoélectrons les plus rapides. $eU_0 = E_{c\max}$.
\item \textbf{Courant de saturation $I_S$} : Courant maximal atteint pour $U \geq 0$. $I_S = e \times (\text{nombre d'électrons émis par seconde})$.
\item \textbf{Rendement quantique $\eta$} : Rapport entre le nombre d'électrons émis par seconde et le nombre de photons incidents par seconde.
\[ \eta = \frac{n_e}{n_{ph}} = \frac{I_S/e}{P/(h\nu)}. \]
\end{itemize}
\end{definition}

\begin{propriete}[Influence des paramètres sur la caractéristique $I=f(U)$]
\begin{center}
\begin{tabularx}{\linewidth}{|l|X|X|}\hline
\textbf{Paramètre varié} & \textbf{Effet sur $U_0$} & \textbf{Effet sur $I_S$} \\ \hline
Fréquence $\nu$ $\uparrow$ (puissance $P$ constante) & $U_0$ $\uparrow$ (car $E_{c\max} \uparrow$) & $I_S$ constant (car $n_{ph} = P/h\nu \downarrow$ mais $\eta$ constant) \\ \hline
Intensité/Puissance $P$ $\uparrow$ ($\nu$ constante) & $U_0$ \textbf{inchangé} & $I_S \uparrow$ proportionnellement à $P$ \\ \hline
Travail d'extraction $W_0$ $\uparrow$ (autre métal) & $U_0 \downarrow$ & $I_S$ peut varier (rendement $\eta$ dépend du métal) \\ \hline
\end{tabularx}
\end{center}
\end{propriete}

\begin{methode}[Lecture de la caractéristique d'une cellule photoélectrique]
Pour exploiter le graphe $I = f(U)$ :
\begin{enumerate}
\item \textbf{Identifier les points particuliers} :
\begin{itemize}
\item l'intersection avec l'axe des tensions ($I = 0$) donne $U = -U_0$ (potentiel d'arrêt) ;
\item le palier pour $U$ grand donne le courant de saturation $I_S$.
\end{itemize}
\item \textbf{Déduire $E_{c\max}$} : $E_{c\max} = eU_0$, puis $\nu$ ou $\lambda$ par la relation d'Einstein si $W_0$ est connu.
\item \textbf{Courant de saturation} : tous les électrons émis par unité de temps atteignent l'anode. Si $N$ photons efficaces arrivent par seconde, $I_S = Ne$.
\item \textbf{Rendement quantique} : $\eta = \dfrac{\text{nombre d'électrons émis}}{\text{nombre de photons incidents}}$, soit $\eta = \dfrac{I_S/e}{P/(h\nu)}$ où $P$ est la puissance lumineuse reçue.
\item \textbf{Effet des paramètres} : augmenter l'intensité lumineuse (à $\nu$ fixé) augmente $I_S$ sans changer $U_0$ ; augmenter $\nu$ augmente $U_0$ sans changer $I_S$ (à intensité constante).
\end{enumerate}
\end{methode}

\begin{exoresolu}[Exploitation d'une caractéristique]
Une cellule photoélectrique dont la cathode a un travail d'extraction $W_0 = \num{2,0}\ \text{eV}$ est éclairée par une radiation monochromatique de puissance $P = \num{1,0}\ \text{mW}$. On relève : potentiel d'arrêt $U_0 = \num{0,75}\ \text{V}$ et courant de saturation $I_S = \num{8,0e-6}\ \text{A}$.
\begin{enumerate}
\item Déterminer la longueur d'onde de la radiation utilisée.
\item Calculer le nombre de photoélectrons émis par seconde.
\item En déduire le rendement quantique de la cellule.
\end{enumerate}
On donne : $h = \num{6,63e-34}\ \text{J.s}$, $c = \num{3,00e8}\ \text{m/s}$, $e = \num{1,60e-19}\ \text{C}$.
\tcblower
\textbf{1. Longueur d'onde.}

Relation d'Einstein : $h\nu = W_0 + eU_0 = 2{,}0 + 0{,}75 = \num{2,75}\ \text{eV} = \num{4,40e-19}\ \text{J}$.
\[ \lambda = \frac{hc}{h\nu} = \frac{\num{6,63e-34} \times \num{3,00e8}}{\num{4,40e-19}} = \num{4,52e-7}\ \text{m} = \num{452}\ \text{nm}. \]
La radiation est bleue (visible).

\textbf{2. Nombre de photoélectrons par seconde.}

Au saturation, chaque électron émis est collecté ; le courant correspond au débit de charge :
\[ I_S = n_e\, e \quad\Longrightarrow\quad n_e = \frac{I_S}{e} = \frac{\num{8,0e-6}}{\num{1,60e-19}} = \num{5,0e13}\ \text{électrons/s}. \]

\textbf{3. Rendement quantique.}

Nombre de photons incidents par seconde :
\[ n_{ph} = \frac{P}{h\nu} = \frac{\num{1,0e-3}}{\num{4,40e-19}} = \num{2,27e15}\ \text{photons/s}. \]
D'où :
\[ \eta = \frac{n_e}{n_{ph}} = \frac{\num{5,0e13}}{\num{2,27e15}} = \num{2,2e-2} \approx 2{,}2\ \%. \]
Seulement environ 2 photons sur 100 arrachent un électron : les rendements quantiques des cellules photoémissives sont effectivement faibles.
\end{exoresolu}

\coursec{Applications de l'effet photoélectrique}

\begin{propriete}[Cellule photoémissive vs Cellule photovoltaïque]
\begin{itemize}
\item \textbf{Cellule photoémissive (tube à vide)} : Électrons arrachés du métal dans le vide. Faible rendement quantique ($\sim 10^{-2}$), temps de réponse très court (ns), utilisée en détection rapide (spectrophotométrie, tubes photomultiplicateurs pour scintillateurs en imagerie médicale \emph{TEP} dans les hôpitaux de Yaoundé/Douala).
\item \textbf{Cellule photovoltaïque (jonction $pn$ semi-conducteur)} : Création de paires électron-trou dans le volume (silicium). Rendement quantique élevé ($> 80\%$ en visible), conversion directe lumière $\to$ électricité. Base des \textbf{panneaux solaires} (électrification rurale, pompage, télécoms au Cameroun).
\end{itemize}
\end{propriete}

\begin{savaistu}[Einstein et le prix Nobel 1921]
C'est pour son explication de l'effet photoélectrique (et non pour la relativité) qu'\textbf{Albert Einstein} a reçu le prix Nobel de physique en 1921. Son modèle corpusculaire de la lumière (photons d'énergie $h\nu$) a résolu le paradoxe de l'indépendance de $E_{c\max}$ vis-à-vis de l'intensité, inexplicable par la théorie ondulatoire classique de Maxwell.
\end{savaistu}

\coursec{Quantité de mouvement du photon}

\begin{propriete}[Quantité de mouvement du photon]
Le photon, particule de masse nulle ($m_0 = 0$), d'énergie $E = h\nu$, transporte une quantité de mouvement (impulsion) :
\[ p = \frac{E}{c} = \frac{h\nu}{c} = \frac{h}{\lambda}. \]
\cle{Analyse dimensionnelle} : $[h/\lambda] = \dfrac{\text{J.s}}{\text{m}} = \dfrac{\text{kg.m}^2\text{.s}^{-2}\text{.s}}{\text{m}} = \text{kg.m/s}$.
\end{propriete}

\begin{propriete}[Pression de radiation]
Un faisceau de $N$ photons par seconde, de quantité de mouvement $p$, incident sur une surface :
\begin{itemize}
\item \textbf{Absorption totale} : Force $F = N p$, Pression $P_{\text{rad}} = \dfrac{I}{c}$ ($I$ : intensité lumineuse en $\text{W/m}^2$).
\item \textbf{Réflexion totale (miroir parfait)} : Force $F = 2 N p$, Pression $P_{\text{rad}} = \dfrac{2I}{c}$.
\end{itemize}
Cet effet est mesurable (ex: queue de comètes poussée par le vent solaire, voile solaire \emph{IKAROS}).
\end{propriete}

\begin{methode}[Quantité de mouvement et pression de radiation]
\begin{enumerate}
\item Le photon, de masse nulle, transporte une quantité de mouvement :
\[ p = \frac{h\nu}{c} = \frac{h}{\lambda}. \]
\item Vérifier l'homogénéité : $[h/\lambda] = \dfrac{\text{J.s}}{\text{m}} = \dfrac{\text{kg.m}^2\text{.s}^{-2}\text{.s}}{\text{m}} = \text{kg.m/s}$.
\item Pour un faisceau de $N$ photons absorbés par unité de temps, la force exercée (pression de radiation) est $F = Np$ ; si les photons sont réfléchis, $F = 2Np$.
\item Toujours exprimer $\lambda$ en mètres avant le calcul.
\end{enumerate}
\end{methode}

\begin{exoresolu}[Quantité de mouvement d'un photon X]
\begin{enumerate}
\item Calculer la quantité de mouvement d'un photon de longueur d'onde $\lambda = \num{0,10}\ \text{nm}$ (rayons X utilisés en radiographie dans les hôpitaux).
\item Comparer à la quantité de mouvement d'un électron de vitesse $v = \num{1,0e6}\ \text{m/s}$.
\end{enumerate}
On donne : $h = \num{6,63e-34}\ \text{J.s}$, $m_e = \num{9,11e-31}\ \text{kg}$.
\tcblower
\textbf{1. Quantité de mouvement du photon.}
\[ p = \frac{h}{\lambda} = \frac{\num{6,63e-34}}{\num{0,10e-9}} = \num{6,63e-24}\ \text{kg.m/s}. \]

\textbf{2. Comparaison avec l'électron.}
\[ p_e = m_e v = \num{9,11e-31} \times \num{1,0e6} = \num{9,11e-25}\ \text{kg.m/s}. \]
Rapport : $\dfrac{p}{p_e} = \dfrac{\num{6,63e-24}}{\num{9,11e-25}} \approx 7{,}3$.

La quantité de mouvement du photon X est du même ordre de grandeur que celle d'un électron rapide : c'est exactement pourquoi, dans l'effet Compton, le choc photon--électron peut être traité comme une collision entre deux particules matérielles avec conservation de la quantité de mouvement totale.
\end{exoresolu}

\coursec{Effet Compton}

\begin{experience}[Expérience de Compton (1923) : diffusion photon-électron]
\textbf{Matériel} : Source de rayons X (tube de Coolidge) $\lambda \approx \num{0,07}\ \text{nm}$, cible en graphite (électrons quasi-libres), spectromètre à cristal (mesure $\lambda'$ par diffraction de Bragg) mobile en angle $\theta$.
\textbf{Observation} : Le spectre diffusé présente deux pics : un pic non modifié ($\lambda' = \lambda$, diffusion sur électrons fortement liés) et un pic modifié ($\lambda' > \lambda$, diffusion sur électrons quasi-libres). Le décalage $\Delta\lambda = \lambda' - \lambda$ dépend de $\theta$ mais \textbf{pas} de $\lambda$ ni du matériau.
\textbf{Interprétation} : Choc élastique entre un photon (énergie $h\nu$, impulsion $h/\lambda$) et un électron libre au repos (masse $m_e$). Conservation de l'énergie et de la quantité de mouvement vectorielle.
\end{experience}

\begin{propriete}[Relation de Compton — \textbf{Exigible au Bac}]
Lors de la diffusion d'un photon sur un électron quasi libre au repos, la longueur d'onde du photon diffusé augmente :
\[ \lambda' - \lambda = \frac{h}{m_e c}(1 - \cos\theta) = \lambda_C (1 - \cos\theta), \]
où :
\begin{itemize}
\item $\theta$ : angle de diffusion (angle entre direction incidente et direction du photon diffusé).
\item $\lambda_C = \dfrac{h}{m_e c} = \num{2,43e-12}\ \text{m} = \num{2,43}\ \text{pm}$ : \textbf{longueur d'onde de Compton de l'électron}.
\end{itemize}
\cle{Cas limites} :
\begin{itemize}
\item $\theta = 0$ : $\lambda' = \lambda$ (pas de diffusion effective, photon continue tout droit).
\item $\theta = \pi$ (rétrodiffusion) : $\lambda' - \lambda = 2\lambda_C$ (variation maximale, perte d'énergie maximale pour le photon).
\end{itemize}
L'énergie cinétique acquise par l'électron de recul : $E_c(e^-) = hc\left(\dfrac{1}{\lambda} - \dfrac{1}{\lambda'}\right)$.
\end{propriete}

\begin{methode}[Formule de Compton]
Lors de la diffusion d'un photon sur un électron quasi libre, la longueur d'onde du photon diffusé augmente :
\begin{enumerate}
\item Écrire la \textbf{relation de Compton} :
\[ \lambda' - \lambda = \frac{h}{m_e c}(1 - \cos\theta) = \lambda_C (1 - \cos\theta), \]
où $\theta$ est l'angle de diffusion et $\lambda_C = \dfrac{h}{m_e c} = \num{2,43e-12}\ \text{m}$ la longueur d'onde de Compton de l'électron.
\item Vérifier que $\theta$ est bien l'angle entre la direction incidente et la direction du photon diffusé.
\item Calculer $\lambda'$, puis l'énergie du photon diffusé $E' = hc/\lambda'$ si demandée.
\item \textbf{Cas limites à connaître} : $\theta = 0 \Rightarrow \lambda' = \lambda$ (pas de diffusion) ; $\theta = \pi \Rightarrow \lambda' - \lambda = 2\lambda_C$ (variation maximale, rétrodiffusion).
\item L'énergie perdue par le photon est cédée à l'électron : $E_c(\text{électron}) = hc\left(\dfrac{1}{\lambda} - \dfrac{1}{\lambda'}\right)$.
\end{enumerate}
\end{methode}

\begin{exoresolu}[Diffusion Compton d'un photon X]
Un photon X de longueur d'onde $\lambda = \num{0,0711}\ \text{nm}$ frappe un électron quasi libre au repos. Le photon diffusé est observé sous un angle $\theta = 90^{\circ}$.
\begin{enumerate}
\item Calculer la longueur d'onde $\lambda'$ du photon diffusé.
\item Calculer l'énergie cinétique acquise par l'électron de recul, en eV.
\end{enumerate}
On donne : $h = \num{6,63e-34}\ \text{J.s}$, $c = \num{3,00e8}\ \text{m/s}$, $m_e = \num{9,11e-31}\ \text{kg}$, $1\ \text{eV} = \num{1,60e-19}\ \text{J}$.
\tcblower
\textbf{1. Longueur d'onde du photon diffusé.}

Longueur d'onde de Compton :
\[ \lambda_C = \frac{h}{m_e c} = \frac{\num{6,63e-34}}{\num{9,11e-31} \times \num{3,00e8}} = \num{2,43e-12}\ \text{m}. \]
Relation de Compton avec $\cos 90^{\circ} = 0$ :
\[ \lambda' = \lambda + \lambda_C(1 - \cos\theta) = \num{7,11e-11} + \num{2,43e-12} = \num{7,35e-11}\ \text{m} = \num{0,0735}\ \text{nm}. \]
On constate $\lambda' > \lambda$ : le photon diffusé est moins énergétique que le photon incident.

\textbf{2. Énergie cinétique de l'électron de recul.}

Conservation de l'énergie : l'énergie perdue par le photon est transférée à l'électron :
\[ E_c = \frac{hc}{\lambda} - \frac{hc}{\lambda'} = hc\left(\frac{1}{\lambda} - \frac{1}{\lambda'}\right). \]
\begin{align*}
E_c &= \num{6,63e-34} \times \num{3,00e8} \times \left(\frac{1}{\num{7,11e-11}} - \frac{1}{\num{7,35e-11}}\right) \\
&= \num{1,989e-25} \times (\num{1,406e10} - \num{1,360e10}) \\
&= \num{1,989e-25} \times \num{4,6e8} = \num{9,1e-17}\ \text{J}.
\end{align*}
Conversion : $E_c = \dfrac{\num{9,1e-17}}{\num{1,60e-19}} \approx \num{5,7e2}\ \text{eV} = \num{0,57}\ \text{keV}$.

L'électron de recul emporte environ $\num{0,57}\ \text{keV}$. L'augmentation de longueur d'onde, inexplicable par l'optique ondulatoire classique, confirme le comportement corpusculaire du photon : c'est la preuve expérimentale décisive apportée par Compton (1923).
\end{exoresolu}

\begin{savaistu}[Arthur Holly Compton — Prix Nobel 1927]
Arthur Holly Compton a reçu le prix Nobel de physique en 1927 (partagé avec C.T.R. Wilson) pour sa découverte de l'effet qui porte son nom. Cette expérience a définitivement convaincu la communauté scientifique de la réalité corpusculaire du photon (quantité de mouvement $p = h/\lambda$), 22 ans après l'hypothèse d'Einstein. Au Cameroun, la diffusion Compton est à la base du fonctionnement des \textbf{scanners à rayons X} (tomodensitométrie) dans les grands hôpitaux : les photons X traversent le corps, subissent des diffusions Compton sur les électrons des tissus, et l'atténuation du faisceau permet de reconstruire l'image 3D des organes.
\end{savaistu}

\begin{attention}[Les 5 erreurs qui coûtent des points au Bac]
\begin{enumerate}
\item \textbf{Oublier de convertir les eV en joules} avant d'utiliser $h$, $c$, $e$ en unités SI. Un calcul de $\lambda_0 = hc/W_0$ avec $W_0$ laissé en eV donne un résultat faux d'un facteur $\num{1,6e-19}$. Convertis \emph{systématiquement} et écris la conversion sur ta copie.
\item \textbf{Inverser la condition de seuil} : l'émission exige $\nu \geq \nu_0$, c'est-à-dire $\lambda \leq \lambda_0$ (et non l'inverse !). Une radiation \emph{plus longue} que $\lambda_0$ est \emph{moins} énergétique. Vérifie toujours avec l'énergie : $E = hc/\lambda$ doit dépasser $W_0$.
\item \textbf{Croire que l'intensité lumineuse modifie $E_{c\max}$ ou $U_0$} : à fréquence fixée, augmenter l'intensité augmente uniquement le \emph{nombre} de photoélectrons (courant de saturation $I_S$), jamais leur énergie maximale. C'est une question de cours classique au Bac.
\item \textbf{Confondre $U_0$ et $-U_0$ sur la caractéristique} : la caractéristique coupe l'axe des tensions en $U = -U_0$ (tension inverse). Le potentiel d'arrêt $U_0$ est une valeur \emph{positive} définie par $eU_0 = E_{c\max}$.
\item \textbf{Dans l'effet Compton : mal gérer les puissances de dix et l'angle}. Les longueurs d'onde sont de l'ordre du picomètre ($\num{e-12}\ \text{m}$) : une erreur d'exposant fausse tout. Vérifie aussi que $\theta$ est l'angle de \emph{diffusion} (entre photon incident et photon diffusé), et rappelle-toi que $\cos 90^{\circ} = 0$ et $\cos 180^{\circ} = -1$. Enfin, n'écris jamais $\lambda' < \lambda$ : le photon diffusé est \emph{toujours} moins énergétique, donc $\lambda' \geq \lambda$.
\end{enumerate}
\end{attention}

\newpage

\coursec{Exercices}

\courssub{Je vérifie mes connaissances}

\begin{exercice}[QCM : les bases de l'effet photoélectrique]
Pour chaque question, choisir la (ou les) bonne(s) réponse(s) en justifiant brièvement.
\begin{enumerate}
\item L'effet photoélectrique a été mis en évidence expérimentalement pour la première fois par :
\begin{enumerate}
\item Einstein \quad b) Hertz \quad c) Compton \quad d) Planck
\end{enumerate}
\item L'extraction d'un électron d'un métal par la lumière n'est possible que si :
\begin{enumerate}
\item l'intensité lumineuse dépasse une valeur seuil ;
\item la fréquence de la radiation dépasse la fréquence seuil $\nu_0$ du métal ;
\item la durée d'éclairage est suffisamment longue ;
\item la longueur d'onde est inférieure à $\lambda_0 = c/\nu_0$.
\end{enumerate}
\item Si l'on augmente l'intensité d'une radiation de fréquence $\nu > \nu_0$ éclairant une cellule photoélectrique :
\begin{enumerate}
\item l'énergie cinétique maximale des photoélectrons augmente ;
\item le nombre de photoélectrons émis par seconde augmente ;
\item le potentiel d'arrêt augmente ;
\item le courant de saturation augmente.
\end{enumerate}
\item La quantité de mouvement d'un photon de longueur d'onde $\lambda$ vaut :
\begin{enumerate}
\item $p = \dfrac{h\lambda}{c}$ \quad b) $p = \dfrac{h}{\lambda}$ \quad c) $p = \dfrac{hc}{\lambda}$ \quad d) $p = \dfrac{h\nu}{c}$
\end{enumerate}
\end{enumerate}
\end{exercice}

\begin{exercice}[Vrai ou faux ? Justifier]
Répondre par vrai ou faux en justifiant chaque réponse par une phrase ou un calcul.
\begin{enumerate}
\item Toute radiation lumineuse, aussi intense soit-elle, peut arracher des électrons à n'importe quel métal.
\item Le potentiel d'arrêt $U_0$ d'une cellule photoélectrique dépend de la fréquence de la radiation incidente mais pas de son intensité.
\item Dans l'effet Compton, le photon diffusé a une longueur d'onde plus petite que celle du photon incident.
\item Le modèle ondulatoire de la lumière permet d'expliquer l'existence du seuil photoélectrique.
\item La constante $\Lambda_C = \dfrac{h}{m_e c}$ qui intervient dans la formule de Compton est homogène à une longueur.
\end{enumerate}
\end{exercice}

\begin{exercice}[Définitions et relations fondamentales]
\begin{enumerate}
\item Définir : effet photoélectrique ; fréquence seuil $\nu_0$ ; travail d'extraction $W_0$ ; potentiel d'arrêt $U_0$ ; rendement quantique d'une cellule photoélectrique.
\item Écrire la relation d'Einstein de l'effet photoélectrique en précisant la signification et l'unité SI de chaque terme.
\item Écrire la relation de Compton reliant $\lambda'$, $\lambda$ et l'angle de diffusion $\theta$. Que représente la grandeur $\dfrac{h}{m_e c}$ ? Donner sa valeur numérique en picomètres.
\end{enumerate}
\end{exercice}

\begin{exercice}[Calculs directs]
Données : $h = \num{6,63e-34}$ J$\cdot$s ; $c = \num{3,00e8}$ m/s ; $e = \num{1,60e-19}$ C ; $m_e = \num{9,11e-31}$ kg.
\begin{enumerate}
\item Le travail d'extraction du césium vaut $W_0 = \num{1,9}$ eV. Calculer sa fréquence seuil $\nu_0$ et sa longueur d'onde seuil $\lambda_0$. Le césium est-il sensible à la lumière visible ($\num{400}$ nm $< \lambda <$ $\num{800}$ nm) ?
\item Une radiation de longueur d'onde $\lambda = \num{400}$ nm éclaire une surface de potassium ($W_0 = \num{2,3}$ eV). Calculer l'énergie cinétique maximale des photoélectrons en joules puis en électronvolts.
\item Un photon X de longueur d'onde $\lambda = \num{0,020}$ nm est diffusé sous l'angle $\theta = 90\degree$. Calculer la variation de longueur d'onde $\Delta\lambda = \lambda' - \lambda$.
\end{enumerate}
\end{exercice}

\courssub{Je m'entraîne}

\begin{exercice}[Métal éclairé par une radiation ultraviolette]
Une plaque de zinc, de longueur d'onde seuil $\lambda_0 = \num{290}$ nm, est éclairée par une radiation ultraviolette de longueur d'onde $\lambda = \num{254}$ nm (lampe à vapeur de mercure).
Données : $h = \num{6,63e-34}$ J$\cdot$s ; $c = \num{3,00e8}$ m/s ; $e = \num{1,60e-19}$ C ; $m_e = \num{9,11e-31}$ kg.
\begin{enumerate}
\item Calculer le travail d'extraction $W_0$ du zinc, en joules puis en électronvolts.
\item Vérifier que l'effet photoélectrique a lieu, puis calculer l'énergie cinétique maximale $E_{c\max}$ des photoélectrons.
\item En déduire la vitesse maximale $v_{\max}$ d'émission des photoélectrons.
\item Calculer le potentiel d'arrêt $U_0$ de la cellule.
\item On remplace la radiation par une lumière rouge de longueur d'onde $\lambda = \num{650}$ nm, d'intensité dix fois plus grande. Y a-t-il émission de photoélectrons ? Justifier.
\end{enumerate}
\end{exercice}

\begin{exercice}[Détermination expérimentale de la constante de Planck]
Une cellule photoélectrique à cathode de potassium est éclairée successivement par quatre radiations monochromatiques. Pour chacune, on mesure le potentiel d'arrêt $U_0$ :
\begin{center}
\begin{tabular}{|c|c|c|c|c|}
\hline
\rowcolor{popblueL} $\lambda$ (nm) & 580 & 500 & 436 & 405 \\
\hline $\nu$ ($\times 10^{14}$ Hz) & 5,17 & 6,00 & 6,88 & 7,41 \\
\hline $U_0$ (V) & 0,24 & 0,58 & 0,93 & 1,15 \\
\hline
\end{tabular}
\end{center}
\begin{enumerate}
\item Établir, à partir de la relation d'Einstein, l'expression du potentiel d'arrêt en fonction de la fréquence : $U_0 = f(\nu)$. Montrer que la courbe obtenue est une droite dont on précisera le coefficient directeur et l'ordonnée à l'origine.
\item Tracer la courbe $U_0 = f(\nu)$ sur papier millimétré (échelles : 1 cm pour $0,5 \times 10^{14}$ Hz ; 1 cm pour $0,2$ V).
\item En déduire une valeur expérimentale de la constante de Planck $h$. Comparer à la valeur de référence $h = \num{6,63e-34}$ J$\cdot$s (calculer l'écart relatif).
\item Déterminer graphiquement la fréquence seuil $\nu_0$ du potassium, puis son travail d'extraction $W_0$ en électronvolts.
\end{enumerate}
\end{exercice}

\begin{exercice}[Exploitation d'une caractéristique $I = f(U)$]
Une cellule photoélectrique est éclairée par un faisceau monochromatique de longueur d'onde $\lambda = \num{450}$ nm et de puissance lumineuse $P = \num{2,0}$ mW. La caractéristique $I = f(U)$ relevée présente un courant de saturation $I_{\text{sat}} = \num{12}$ \textmu A et un potentiel d'arrêt $U_0 = \num{0,90}$ V.
Données : $h = \num{6,63e-34}$ J$\cdot$s ; $c = \num{3,00e8}$ m/s ; $e = \num{1,60e-19}$ C.
\begin{enumerate}
\item Faire le schéma du montage permettant de tracer la caractéristique $I = f(U)$ de la cellule.
\item Calculer l'énergie cinétique maximale des photoélectrons et le travail d'extraction du métal constituant la cathode.
\item Calculer le nombre $n_e$ d'électrons émis par seconde lorsque le courant est saturé.
\item Calculer le nombre $n_p$ de photons incidents sur la cathode par seconde.
\item En déduire le rendement quantique $\eta$ de la cellule. Commenter la valeur obtenue.
\item On double la puissance du faisceau sans changer sa longueur d'onde. Que deviennent $I_{\text{sat}}$ et $U_0$ ? Justifier.
\end{enumerate}
\end{exercice}

\begin{exercice}[L'effet Compton est-il observable en lumière visible ?]
On considère la diffusion d'un photon par un électron libre supposé au repos.
Données : $\Lambda_C = \dfrac{h}{m_e c} = \num{2,43e-12}$ m ; $h = \num{6,63e-34}$ J$\cdot$s ; $c = \num{3,00e8}$ m/s.
\begin{enumerate}
\item Rappeler la relation de Compton et calculer la variation maximale possible de longueur d'onde $\Delta\lambda_{\max}$. Pour quel angle de diffusion est-elle obtenue ?
\item Un photon de lumière verte de longueur d'onde $\lambda = \num{500}$ nm est diffusé sous l'angle $\theta = 180\degree$. Calculer la variation relative $\dfrac{\Delta\lambda}{\lambda}$. Conclure sur la possibilité d'observer l'effet Compton dans le domaine visible.
\item Reprendre la question précédente pour un photon X de longueur d'onde $\lambda = \num{0,050}$ nm diffusé à $\theta = 180\degree$. Conclure sur le domaine du spectre électromagnétique où l'effet Compton est significatif.
\item Expliquer en quelques lignes pourquoi l'effet Compton ne peut pas être interprété par le modèle ondulatoire classique, qui prévoit une diffusion sans changement de fréquence.
\end{enumerate}
\end{exercice}

\courssub{Je me prépare au Bac}

\begin{exercice}[Type Bac — La cellule photoélectrique du laboratoire de lycée]
Au laboratoire de physique d'un lycée de Yaoundé, un professeur étudie avec ses élèves de Terminale C une cellule photoélectrique dont la cathode est recouverte d'un métal de longueur d'onde seuil $\lambda_0 = \num{275}$ nm. La cathode est éclairée par la radiation ultraviolette de longueur d'onde $\lambda = \num{254}$ nm issue d'une lampe à vapeur de mercure.
Données : $h = \num{6,63e-34}$ J$\cdot$s ; $c = \num{3,00e8}$ m/s ; $e = \num{1,60e-19}$ C ; $m_e = \num{9,11e-31}$ kg ; $\num{1}$ eV $= \num{1,60e-19}$ J.
\begin{enumerate}
\item Définir l'effet photoélectrique. Pourquoi le modèle ondulatoire de la lumière ne permet-il pas d'expliquer l'existence d'une longueur d'onde seuil ?
\item Calculer le travail d'extraction $W_0$ du métal, en joules puis en électronvolts.
\item Énoncer la relation d'Einstein de l'effet photoélectrique. Calculer l'énergie cinétique maximale $E_{c\max}$ des photoélectrons émis.
\item En déduire la vitesse maximale $v_{\max}$ des photoélectrons à la sortie de la cathode. (On justifiera l'approximation non relativiste en comparant $v_{\max}$ à $c$.)
\item On applique entre l'anode et la cathode une tension $U_{AC} = -U_0$ qui annule le courant photoélectrique.
\begin{enumerate}
\item Faire le schéma du montage en représentant la cellule, le générateur de tension réglable, le voltmètre et le milliampèremètre.
\item Établir, à l'aide du théorème de l'énergie cinétique, la relation $eU_0 = E_{c\max}$.
\item Calculer la valeur numérique du potentiel d'arrêt $U_0$.
\end{enumerate}
\item La lampe éclaire maintenant la cathode avec une puissance $P = \num{1,5}$ mW. On mesure un courant de saturation $I_{\text{sat}} = \num{8,0}$ \textmu A.
\begin{enumerate}
\item Calculer le nombre de photons incidents par seconde sur la cathode.
\item Calculer le rendement quantique de la cellule.
\end{enumerate}
\item On interpose entre la lampe et la cellule un filtre qui ne laisse passer que la radiation de longueur d'onde $\lambda' = \num{313}$ nm, d'intensité cinq fois plus grande. Le courant photoélectrique persiste-t-il ? Justifier par un calcul.
\end{enumerate}
\end{exercice}

\begin{exercice}[Type Bac — Effet Compton et imagerie médicale à l'hôpital]
Dans un service de radiologie d'un hôpital de Douala, des photons X de longueur d'onde $\lambda = \num{0,050}$ nm interagissent avec les électrons des tissus. On étudie la diffusion d'un photon par un électron libre supposé initialement au repos : le photon est diffusé sous l'angle $\theta = 120\degree$ par rapport à sa direction incidente.
Données : $h = \num{6,63e-34}$ J$\cdot$s ; $c = \num{3,00e8}$ m/s ; $m_e = \num{9,11e-31}$ kg ; $e = \num{1,60e-19}$ C ; constante de Compton $\Lambda_C = \dfrac{h}{m_e c} = \num{2,43e-12}$ m.
\begin{enumerate}
\item Donner l'expression de la quantité de mouvement d'un photon en fonction de sa longueur d'onde. Calculer la quantité de mouvement $p$ du photon incident.
\item Énoncer la relation de Compton. Quelles grandeurs sont conservées lors de la collision photon-électron ?
\item Calculer la longueur d'onde $\lambda'$ du photon diffusé.
\item Calculer la variation relative de longueur d'onde $\dfrac{\Delta\lambda}{\lambda}$ en pourcentage. Commenter.
\item Calculer l'énergie $E$ du photon incident et l'énergie $E'$ du photon diffusé, en joules puis en kiloélectronvolts (keV).
\item En appliquant la conservation de l'énergie, montrer que l'énergie cinétique de l'électron de recul vaut $E_c = E - E'$. Calculer sa valeur en keV.
\item Expliquer en quelques lignes pourquoi l'effet Compton, préjudiciable à la qualité des images radiographiques (diffusion du rayonnement), impose de protéger le personnel soignant par des tabliers de plomb.
\end{enumerate}
\end{exercice}

\begin{exercice}[Type Bac — Synthèse : effet photoélectrique, effet Compton et dualité de la lumière]
Un professeur propose à ses élèves de Terminale C une étude comparative des deux expériences historiques qui ont imposé le modèle corpusculaire de la lumière.

\textbf{Partie A — Effet photoélectrique.} Une cellule photoélectrique à cathode de sodium ($W_0 = \num{2,28}$ eV) est éclairée par une radiation de fréquence $\nu = \num{8,50e14}$ Hz.
\begin{enumerate}
\item Vérifier que l'effet photoélectrique a lieu et calculer l'énergie cinétique maximale des photoélectrons en eV.
\item Calculer le potentiel d'arrêt $U_0$.
\item Décrire qualitativement l'allure de la caractéristique $I = f(U)$ de la cellule pour deux intensités lumineuses différentes $E_1 < E_2$ (même fréquence), en précisant les points particuliers. On pourra s'aider d'un schéma.
\end{enumerate}

\textbf{Partie B — Effet Compton.} Un photon de longueur d'onde $\lambda = \num{0,0708}$ nm frappe un électron libre au repos ; on détecte un photon diffusé sous l'angle $\theta = 90\degree$.
\begin{enumerate}
\setcounter{enumi}{3}
\item Calculer la longueur d'onde $\lambda'$ du photon diffusé et l'énergie cinétique de l'électron de recul en keV.
\end{enumerate}

\textbf{Partie C — Comparaison.}
\begin{enumerate}
\setcounter{enumi}{4}
\item Recopier et compléter le tableau comparatif suivant :
\begin{center}
\begin{tabularx}{\linewidth}{|l|X|X|}
\hline
\rowcolor{popblueL} \textbf{Critère} & \textbf{Effet photoélectrique} & \textbf{Effet Compton} \\
\hline Cible du photon & & \\
\hline Destin du photon (absorbé ou diffusé) & & \\
\hline Grandeurs conservées & & \\
\hline Domaine spectral typique & & \\
\hline
\end{tabularx}
\end{center}
\item Conclure en deux ou trois phrases : en quoi ces deux effets, inexplicables par le modèle ondulatoire, illustrent-ils la dualité onde-corpuscule de la lumière ? Citer une application technologique de chaque effet dans le contexte camerounais ou africain.
\end{enumerate}
Données : $h = \num{6,63e-34}$ J$\cdot$s ; $c = \num{3,00e8}$ m/s ; $e = \num{1,60e-19}$ C ; $\Lambda_C = \num{2,43e-12}$ m ; $\num{1}$ eV $= \num{1,60e-19}$ J.
\end{exercice}



\newpage

\coursec{Corrigés des exercices}

\coursec{Exercices corrigés — Effet photoélectrique et effet Compton}

\begin{corrige}[Exercice 1 — QCM : les bases de l'effet photoélectrique]
\begin{enumerate}
\item \textbf{Réponse b) : Hertz.} C'est Heinrich Hertz qui observe accidentellement le phénomène en 1887 : une étincelle éclate plus facilement entre deux sphères métalliques lorsqu'elles sont éclairées par un rayonnement ultraviolet. \textbf{Einstein} (1905) en donne l'interprétation corpusculaire (quantas de lumière), \textbf{Planck} (1900) introduit les quantas pour le corps noir, et \textbf{Compton} (1923) étudie la diffusion des rayons X.
\item \textbf{Réponses b) et d).} L'extraction d'un électron exige que l'énergie du photon soit supérieure au travail d'extraction : $h\nu > W_0 = h\nu_0$, soit $\nu > \nu_0$, ce qui équivaut à $\lambda < \lambda_0 = c/\nu_0$. Ni l'intensité (a) ni la durée d'éclairage (c) ne peuvent compenser une fréquence insuffisante : c'est précisément ce que le modèle ondulatoire ne parvenait pas à expliquer.
\item \textbf{Réponses b) et d).} L'énergie cinétique maximale $E_{c,\max} = h\nu - W_0$ ne dépend que de la fréquence $\nu$ et du métal : elle est \textbf{inchangée} quand l'intensité varie, donc le potentiel d'arrêt $U_0 = E_{c,\max}/e$ reste constant (réponses a et c fausses). En revanche, augmenter l'intensité revient à augmenter le nombre de photons incidents par seconde ; le nombre de photoélectrons émis par seconde augmente donc, et avec lui le courant de saturation $I_{\text{sat}} = n_e e$.
\item \textbf{Réponses b) et d).} La quantité de mouvement du photon vaut
\[ p = \frac{h}{\lambda} = \frac{h\nu}{c} \quad (\text{car } \lambda = c/\nu). \]
La réponse c) $hc/\lambda$ est l'\emph{énergie} du photon $E = h\nu$, et non sa quantité de mouvement : attention à ne pas confondre.
\end{enumerate}
\end{corrige}

\begin{corrige}[Exercice 2 — Vrai ou faux ?]
\begin{enumerate}
\item \textbf{Faux.} L'effet photoélectrique ne se produit que si la fréquence de la radiation dépasse la fréquence seuil du métal : $\nu > \nu_0$ (c'est-à-dire $\lambda < \lambda_0$). En dessous du seuil, aucun électron n'est extrait, \emph{quelle que soit l'intensité} du rayonnement. C'est l'observation qui a mis en échec le modèle ondulatoire.
\item \textbf{Vrai.} D'après la relation d'Einstein : $eU_0 = E_{c,\max} = h\nu - W_0$, donc $U_0 = \dfrac{h\nu - W_0}{e}$ ne dépend que de la fréquence $\nu$ et du métal (via $W_0$), jamais de l'intensité lumineuse.
\item \textbf{Faux.} Dans l'effet Compton, le photon cède une partie de son énergie à l'électron de recul ; sa fréquence diminue donc et sa longueur d'onde \textbf{augmente} : $\lambda' = \lambda + \Lambda_C(1-\cos\theta) > \lambda$ puisque $1-\cos\theta \geq 0$.
\item \textbf{Faux.} Le modèle ondulatoire prévoit que l'énergie transportée par une onde est proportionnelle à son intensité et s'accumule progressivement dans le métal : une onde de faible fréquence mais très intense (ou éclairant longtemps) devrait finir par extraire des électrons. Il ne prévoit donc \emph{aucun seuil en fréquence}, contrairement à l'expérience.
\item \textbf{Vrai.} Analyse dimensionnelle : $[h] = \text{J}\cdot\text{s} = \text{kg}\cdot\text{m}^2\cdot\text{s}^{-1}$ et $[m_e c] = \text{kg}\cdot\text{m}\cdot\text{s}^{-1}$, d'où
\[ \left[\frac{h}{m_e c}\right] = \frac{\text{kg}\cdot\text{m}^2\cdot\text{s}^{-1}}{\text{kg}\cdot\text{m}\cdot\text{s}^{-1}} = \text{m}. \]
$\Lambda_C$ est bien homogène à une longueur (c'est la longueur d'onde de Compton de l'électron).
\end{enumerate}
\end{corrige}

\begin{corrige}[Exercice 3 — Définitions et relations fondamentales]
\begin{enumerate}
\item \textbf{Effet photoélectrique :} phénomène d'émission d'électrons (appelés photoélectrons) par un métal soumis à un rayonnement électromagnétique de fréquence suffisante. \textbf{Fréquence seuil $\nu_0$ :} fréquence minimale que doit avoir le rayonnement incident pour que l'effet photoélectrique se produise ; elle est caractéristique du métal. \textbf{Travail d'extraction $W_0$ :} énergie minimale qu'il faut fournir à un électron pour l'extraire du métal ; $W_0 = h\nu_0$. \textbf{Potentiel d'arrêt $U_0$ :} valeur absolue de la tension (anode négative par rapport à la cathode) qu'il faut appliquer pour annuler le courant photoélectrique ; elle vérifie $eU_0 = E_{c,\max}$. \textbf{Rendement quantique $\eta$ :} rapport du nombre d'électrons émis par seconde au nombre de photons incidents par seconde : $\eta = n_e/n_p$ (grandeur sans unité, toujours inférieure à 1).
\item \textbf{Relation d'Einstein :}
\[ E_{c,\max} = h\nu - W_0 \]
avec $E_{c,\max}$ énergie cinétique maximale des photoélectrons (en J), $h = \SI{6,63e-34}{\joule\second}$ la constante de Planck, $\nu$ la fréquence du rayonnement incident (en Hz) et $W_0$ le travail d'extraction du métal (en J). Interprétation : le photon, d'énergie $h\nu$, est absorbé par un électron ; une partie $W_0$ sert à extraire l'électron, le reste devient son énergie cinétique.
\item \textbf{Relation de Compton :}
\[ \lambda' - \lambda = \frac{h}{m_e c}\,(1 - \cos\theta) = \Lambda_C(1-\cos\theta) \]
où $\lambda$ et $\lambda'$ sont les longueurs d'onde du photon incident et du photon diffusé, et $\theta$ l'angle de diffusion. La grandeur $\Lambda_C = \dfrac{h}{m_e c}$ est la \textbf{longueur d'onde de Compton} de l'électron :
\[ \Lambda_C = \frac{\SI{6,63e-34}{\joule\second}}{\SI{9,11e-31}{\kilogram}\times\SI{3,00e8}{\meter\per\second}} \approx \SI{2,43e-12}{\meter} = \SI{2,43}{\pico\meter}. \]
\end{enumerate}
\end{corrige}

\begin{corrige}[Exercice 4 — Calculs directs]
\begin{enumerate}
\item Conversion : $W_0 = \SI{1,9}{\electronvolt} \times \SI{1,60e-19}{\joule\per\electronvolt} = \SI{3,04e-19}{\joule}$.
\[ \nu_0 = \frac{W_0}{h} = \frac{\SI{3,04e-19}{\joule}}{\SI{6,63e-34}{\joule\second}} = \SI{4,58e14}{\hertz}. \]
\[ \lambda_0 = \frac{c}{\nu_0} = \frac{\SI{3,00e8}{\meter\per\second}}{\SI{4,58e14}{\hertz}} = \SI{6,55e-7}{\meter} = \SI{655}{\nano\meter}. \]
Comme $\SI{400}{\nano\meter} < \lambda_0 = \SI{655}{\nano\meter} < \SI{800}{\nano\meter}$, la longueur d'onde seuil se situe dans le visible (rouge) : \textbf{le césium est sensible à la lumière visible} (pour toute radiation telle que $\lambda < \SI{655}{\nano\meter}$). C'est pourquoi on l'utilise dans les cellules photoélectriques.
\item Énergie du photon incident :
\[ E = \frac{hc}{\lambda} = \frac{\SI{6,63e-34}{\joule\second}\times\SI{3,00e8}{\meter\per\second}}{\SI{400e-9}{\meter}} = \SI{4,97e-19}{\joule} = \frac{\SI{4,97e-19}{\joule}}{\SI{1,60e-19}{\joule\per\electronvolt}} = \SI{3,11}{\electronvolt}. \]
Relation d'Einstein :
\[ E_{c,\max} = E - W_0 = \SI{3,11}{\electronvolt} - \SI{2,30}{\electronvolt} = \SI{0,81}{\electronvolt} = \SI{1,30e-19}{\joule}. \]
\item Pour $\theta = 90\degree$, $\cos\theta = 0$, donc :
\[ \Delta\lambda = \Lambda_C(1-0) = \Lambda_C = \frac{\SI{6,63e-34}{\joule\second}}{\SI{9,11e-31}{\kilogram}\times\SI{3,00e8}{\meter\per\second}} = \SI{2,43e-12}{\meter} = \SI{0,00243}{\nano\meter}. \]
\end{enumerate}
\end{corrige}

\begin{corrige}[Exercice 5 — Métal éclairé par une radiation ultraviolette]
\begin{enumerate}
\item Travail d'extraction du zinc :
\[ W_0 = \frac{hc}{\lambda_0} = \frac{\SI{6,63e-34}{\joule\second}\times\SI{3,00e8}{\meter\per\second}}{\SI{290e-9}{\meter}} = \SI{6,86e-19}{\joule} = \frac{\SI{6,86e-19}{\joule}}{\SI{1,60e-19}{\joule\per\electronvolt}} = \SI{4,29}{\electronvolt}. \]
\item Condition d'existence : $\lambda = \SI{254}{\nano\meter} < \lambda_0 = \SI{290}{\nano\meter}$, donc $\nu > \nu_0$ : \textbf{l'effet photoélectrique a lieu}.
Énergie du photon : $E = \dfrac{hc}{\lambda} = \dfrac{\SI{6,63e-34}{\joule\second}\times\SI{3,00e8}{\meter\per\second}}{\SI{254e-9}{\meter}} = \SI{7,83e-19}{\joule} = \SI{4,89}{\electronvolt}$.
\[ E_{c,\max} = E - W_0 = \SI{4,89}{\electronvolt} - \SI{4,29}{\electronvolt} = \SI{0,60}{\electronvolt} = \SI{9,6e-20}{\joule}. \]
\item De $E_{c,\max} = \dfrac{1}{2}m_e v_{\max}^2$ on tire :
\[ v_{\max} = \sqrt{\frac{2E_{c,\max}}{m_e}} = \sqrt{\frac{2\times\SI{9,6e-20}{\joule}}{\SI{9,11e-31}{\kilogram}}} = \SI{4,6e5}{\meter\per\second}. \]
Vérification : $v_{\max}/c \approx \num{1,5e-3} \ll 1$ : l'approximation non relativiste est parfaitement justifiée.
\item Potentiel d'arrêt :
\[ eU_0 = E_{c,\max} \quad\Rightarrow\quad U_0 = \frac{E_{c,\max}}{e} = \frac{\SI{0,60}{\electronvolt}}{e} = \SI{0,60}{\volt}. \]
\item La lumière rouge a $\lambda = \SI{650}{\nano\meter} > \lambda_0 = \SI{290}{\nano\meter}$, soit $\nu < \nu_0$ : l'énergie des photons ($E = hc/\lambda \approx \SI{1,9}{\electronvolt}$) est inférieure au travail d'extraction ($\SI{4,29}{\electronvolt}$). \textbf{Il n'y a aucune émission de photoélectrons}, même avec une intensité dix fois plus grande : c'est la fréquence, et non l'intensité, qui conditionne l'extraction.
\end{enumerate}
\end{corrige}

\begin{corrige}[Exercice 6 — Détermination expérimentale de la constante de Planck]
\begin{enumerate}
\item D'après la relation d'Einstein et la définition du potentiel d'arrêt :
\[ eU_0 = h\nu - W_0 \quad\Rightarrow\quad U_0 = \frac{h}{e}\,\nu - \frac{W_0}{e}. \]
C'est une fonction affine de $\nu$ : la courbe $U_0 = f(\nu)$ est une \textbf{droite} de coefficient directeur $a = \dfrac{h}{e}$ et d'ordonnée à l'origine $b = -\dfrac{W_0}{e}$ (négative).
\item \textit{Tracé sur papier millimétré :} on porte les quatre points $(\nu ; U_0)$ : $(5,17\times10^{14} ; 0,24)$, $(6,00\times10^{14} ; 0,58)$, $(6,88\times10^{14} ; 0,93)$, $(7,41\times10^{14} ; 1,15)$. Les points sont sensiblement alignés ; on trace la droite moyenne passant au plus près des points.
\item Coefficient directeur mesuré sur la droite (points extrêmes) :
\[ a = \frac{\num{1,15} - \num{0,24}}{(\num{7,41} - \num{5,17})\times 10^{14}} = \frac{\num{0,91}}{\num{2,24e14}} = \num{4,06e-15}~\text{V}\cdot\text{s}. \]
Or $a = h/e$, d'où :
\[ h_{\text{exp}} = e\,a = \SI{1,60e-19}{\coulomb}\times\num{4,06e-15}~\text{V}\cdot\text{s} = \SI{6,50e-34}{\joule\second}. \]
Écart relatif par rapport à la valeur de référence :
\[ \frac{|h_{\text{exp}} - h|}{h} = \frac{|\num{6,50} - \num{6,63}|}{\num{6,63}} \times 100 \approx \num{2,0}\%. \]
Accord très satisfaisant compte tenu de la précision des mesures de $U_0$.
\item La fréquence seuil est l'abscisse du point où la droite coupe l'axe des fréquences ($U_0 = 0$) : lecture graphique $\nu_0 \approx \num{4,6e14}$ Hz (calcul : $\nu_0 = -b/a \approx \num{4,58e14}$ Hz).
\[ W_0 = h\nu_0 = \SI{6,63e-34}{\joule\second}\times\num{4,58e14}~\text{Hz} = \SI{3,04e-19}{\joule} = \frac{\SI{3,04e-19}{\joule}}{\SI{1,60e-19}{\joule\per\electronvolt}} = \SI{1,90}{\electronvolt}. \]
C'est bien la valeur tabulée du travail d'extraction du potassium.
\end{enumerate}
\end{corrige}

\begin{corrige}[Exercice 7 — Exploitation d'une caractéristique $I = f(U)$]
\begin{enumerate}
\item \textbf{Montage :} la cellule photoélectrique est alimentée par un générateur de tension continue réglable $U$ (tension anode-cathode) ; un milliampèremètre (ou microampèremètre) est placé en série pour mesurer $I$, et un voltmètre en dérivation aux bornes de la cellule pour mesurer $U$.
\begin{popfigure}
\begin{tikzpicture}[european, scale=0.9]
\draw (0,0) to[battery1, l_={$U$}] (0,2.5)
to[ammeter, l={$A$}] (3,2.5)
to[lamp, l={$\mathcal{C}$}] (3,0) -- (0,0);
\draw (3,2.5) -- (5,2.5) to[voltmeter, l={$V$}] (5,0) -- (3,0);
\end{tikzpicture}
\legende{Montage pour relever la caractéristique $I = f(U)$ d'une cellule photoélectrique ($\mathcal{C}$).}
\end{popfigure}
\item Énergie cinétique maximale des photoélectrons :
\[ E_{c,\max} = eU_0 = \SI{1,60e-19}{\coulomb}\times\SI{0,90}{\volt} = \SI{1,44e-19}{\joule} = \SI{0,90}{\electronvolt}. \]
Énergie du photon incident :
\[ E = \frac{hc}{\lambda} = \frac{\SI{6,63e-34}{\joule\second}\times\SI{3,00e8}{\meter\per\second}}{\SI{450e-9}{\meter}} = \SI{4,42e-19}{\joule} = \SI{2,76}{\electronvolt}. \]
Travail d'extraction :
\[ W_0 = E - E_{c,\max} = \SI{2,76}{\electronvolt} - \SI{0,90}{\electronvolt} = \SI{1,86}{\electronvolt} = \SI{2,98e-19}{\joule}. \]
\item En régime de saturation, tous les électrons émis atteignent l'anode : $I_{\text{sat}} = n_e e$, donc
\[ n_e = \frac{I_{\text{sat}}}{e} = \frac{\SI{12e-6}{\ampere}}{\SI{1,60e-19}{\coulomb}} = \num{7,5e13}~\text{électrons/s}. \]
\item Nombre de photons incidents par seconde (puissance = énergie délivrée par seconde) :
\[ n_p = \frac{P}{E} = \frac{\SI{2,0e-3}{\watt}}{\SI{4,42e-19}{\joule}} = \num{4,5e15}~\text{photons/s}. \]
\item Rendement quantique :
\[ \eta = \frac{n_e}{n_p} = \frac{\num{7,5e13}}{\num{4,5e15}} = \num{1,7e-2} \approx \SI{1,7}{\percent}. \]
\textit{Commentaire :} ce rendement est faible mais tout à fait typique des cellules photoémissives (effet photoélectrique externe) : la plupart des photons sont réfléchis par la surface ou absorbés sans émission d'électron vers l'extérieur.
\item Doubler la puissance à longueur d'onde constante revient à doubler le nombre de photons incidents par seconde : le nombre de photoélectrons émis par seconde double, donc \textbf{$I_{\text{sat}}$ double} ($I'_{\text{sat}} = \SI{24}{\micro\ampere}$). En revanche \textbf{$U_0$ reste inchangé} ($\SI{0,90}{\volt}$) car il ne dépend que de la fréquence $\nu$ et du métal : $eU_0 = h\nu - W_0$.
\end{enumerate}
\end{corrige}

\begin{corrige}[Exercice 8 — L'effet Compton est-il observable en lumière visible ?]
\begin{enumerate}
\item Relation de Compton : $\lambda' - \lambda = \Lambda_C(1-\cos\theta)$. La variation est maximale lorsque $1-\cos\theta$ est maximal, c'est-à-dire pour $\cos\theta = -1$, soit $\theta = 180\degree$ (diffusion vers l'arrière, « rétrodiffusion ») :
\[ \Delta\lambda_{\max} = 2\Lambda_C = 2\times\SI{2,43e-12}{\meter} = \SI{4,86e-12}{\meter} = \SI{4,86}{\pico\meter}. \]
\item Pour $\lambda = \SI{500}{\nano\meter} = \num{5,00e5}~\text{pm}$ et $\theta = 180\degree$ :
\[ \frac{\Delta\lambda}{\lambda} = \frac{\num{4,86}}{\num{5,00e5}} \approx \num{9,7e-6} \approx \num{1e-3}\%. \]
La variation relative est de l'ordre du millionième : totalement indétectable avec les spectromètres optiques usuels. \textbf{Conclusion : l'effet Compton n'est pas observable dans le domaine visible} ; la lumière visible y est diffusée pratiquement sans changement de longueur d'onde.
\item Pour $\lambda = \SI{0,050}{\nano\meter} = \num{50}~\text{pm}$ et $\theta = 180\degree$ :
\[ \frac{\Delta\lambda}{\lambda} = \frac{\num{4,86}}{\num{50}} \approx \num{9,7e-2} = \num{9,7}\%. \]
La variation est de près de 10 \% : parfaitement mesurable. \textbf{Conclusion : l'effet Compton est significatif dans le domaine des rayons X et $\gamma$} (longueurs d'onde de l'ordre du picomètre au nanomètre), car $\Delta\lambda$ (fixée par $\Lambda_C$) devient comparable à $\lambda$.
\item Dans le modèle ondulatoire classique (diffusion Thomson), l'onde incidente fait osciller l'électron à \emph{sa propre fréquence} ; l'électron réémet alors une onde électromagnétique de \textbf{même fréquence} : la diffusion est élastique, sans changement de longueur d'onde, quel que soit l'angle. Or l'expérience de Compton montre un décalage $\Delta\lambda$ qui dépend de l'angle de diffusion. Seul le modèle corpusculaire — collision entre un photon (énergie $h\nu$, quantité de mouvement $h/\lambda$) et un électron libre, avec conservation de l'énergie et de la quantité de mouvement — rend compte quantitativement de ce décalage.
\end{enumerate}
\end{corrige}

\begin{corrige}[Exercice 9 — Type Bac : la cellule photoélectrique du laboratoire de lycée]
\begin{enumerate}
\item \textbf{Définition :} l'effet photoélectrique est l'émission d'électrons par un métal soumis à un rayonnement électromagnétique de fréquence supérieure à une fréquence seuil $\nu_0$ caractéristique du métal.
\textbf{Échec du modèle ondulatoire :} dans le modèle ondulatoire, l'énergie transportée par une onde est proportionnelle à son intensité et s'accumule continûment dans le temps ; une radiation de n'importe quelle fréquence, suffisamment intense ou éclairant assez longtemps, devrait donc finir par extraire des électrons. L'existence d'un seuil net en fréquence (ou en longueur d'onde), indépendant de l'intensité et de la durée d'éclairage, est inexplicable dans ce cadre.
\item Travail d'extraction :
\[ W_0 = \frac{hc}{\lambda_0} = \frac{\SI{6,63e-34}{\joule\second}\times\SI{3,00e8}{\meter\per\second}}{\SI{275e-9}{\meter}} = \SI{7,24e-19}{\joule} = \frac{\SI{7,24e-19}{\joule}}{\SI{1,60e-19}{\joule\per\electronvolt}} = \SI{4,53}{\electronvolt}. \]
\item \textbf{Relation d'Einstein :} $E_{c,\max} = h\nu - W_0 = hc\left(\dfrac{1}{\lambda} - \dfrac{1}{\lambda_0}\right)$.
Énergie du photon : $E = \dfrac{hc}{\lambda} = \dfrac{\SI{6,63e-34}{\joule\second}\times\SI{3,00e8}{\meter\per\second}}{\SI{254e-9}{\meter}} = \SI{7,83e-19}{\joule} = \SI{4,89}{\electronvolt}$.
\[ E_{c,\max} = \SI{4,89}{\electronvolt} - \SI{4,53}{\electronvolt} = \SI{0,36}{\electronvolt} = \SI{5,8e-20}{\joule}. \]
\item Vitesse maximale :
\[ v_{\max} = \sqrt{\frac{2E_{c,\max}}{m_e}} = \sqrt{\frac{2\times\SI{5,8e-20}{\joule}}{\SI{9,11e-31}{\kilogram}}} = \SI{3,6e5}{\meter\per\second}. \]
Comparaison : $v_{\max}/c \approx \num{1,2e-3} \ll 1$ : l'électron est largement non relativiste, l'expression classique $E_c = \frac{1}{2}m_ev^2$ est légitime.
\item \textbf{a) Schéma :} générateur de tension réglable alimentant la cellule (borne négative vers l'anode pour freiner les électrons : $U_{AC} < 0$), microampèremètre en série, voltmètre en dérivation aux bornes de la cellule (même montage qu'à l'exercice 7).
\textbf{b) Démonstration :} entre la cathode C et l'anode A, l'électron (charge $-e$) ne subit que la force électrique, dont le travail est :
\[ W_{C\to A} = -e\,(V_A - V_C) = -e\,U_{AC} = -e(-U_0) = -eU_0. \]
Au potentiel d'arrêt, même l'électron le plus rapide arrive sur l'anode avec une vitesse nulle. Le théorème de l'énergie cinétique s'écrit :
\[ \Delta E_c = 0 - E_{c,\max} = W_{C\to A} = -eU_0 \quad\Rightarrow\quad \boxed{eU_0 = E_{c,\max}}. \]
\textbf{c)} $U_0 = \dfrac{E_{c,\max}}{e} = \SI{0,36}{\volt}$.
\item \textbf{a) Nombre de photons incidents par seconde :}
\[ n_p = \frac{P}{E} = \frac{\SI{1,5e-3}{\watt}}{\SI{7,83e-19}{\joule}} = \num{1,9e15}~\text{photons/s}. \]
\textbf{b) Nombre d'électrons collectés par seconde en saturation :}
\[ n_e = \frac{I_{\text{sat}}}{e} = \frac{\SI{8,0e-6}{\ampere}}{\SI{1,60e-19}{\coulomb}} = \num{5,0e13}~\text{électrons/s}. \]
Rendement quantique :
\[ \eta = \frac{n_e}{n_p} = \frac{\num{5,0e13}}{\num{1,9e15}} = \num{2,6e-2} \approx \SI{2,6}{\percent}. \]
\item La radiation filtrée a $\lambda' = \SI{313}{\nano\meter} > \lambda_0 = \SI{275}{\nano\meter}$, donc $\nu' < \nu_0$. Vérification énergétique : $E' = hc/\lambda' = \SI{6,36e-19}{\joule} \approx \SI{3,97}{\electronvolt} < W_0 = \SI{4,53}{\electronvolt}$. \textbf{Le courant photoélectrique ne persiste pas} : l'énergie de chaque photon est insuffisante pour extraire un électron, et multiplier l'intensité par cinq n'y change rien (le seuil est un seuil en fréquence, pas en intensité).
\end{enumerate}
\end{corrige}

\begin{attention}[Barème indicatif et points de vigilance — Exercice 9 (sur 10 pts)]
\textbf{Barème :} 1) 1 pt ; 2) 1 pt ; 3) 1,5 pt ; 4) 1 pt ; 5a) 1 pt ; 5b) 1,5 pt ; 5c) 0,5 pt ; 6a) 0,75 pt ; 6b) 0,75 pt ; 7) 1 pt.
\textbf{Points de vigilance :} citer explicitement la condition $\lambda < \lambda_0$ (ou $\nu > \nu_0$) avant tout calcul ; convertir systématiquement les eV en joules avant d'utiliser $E_c = \frac{1}{2}m_ev^2$ ; dans la démonstration 5b), préciser que le travail de la force électrique est résistant ($W = -eU_0 < 0$) et que l'électron arrive à vitesse nulle ; ne pas oublier que le rendement quantique s'exprime sans unité (ou en \%) ; pour la question 7, insister sur le fait que l'intensité ne compense jamais une fréquence inférieure au seuil — c'est la justification attendue par le correcteur.
\end{attention}

\begin{corrige}[Exercice 10 — Type Bac : effet Compton et imagerie médicale à l'hôpital]
\begin{enumerate}
\item La quantité de mouvement d'un photon de longueur d'onde $\lambda$ est $p = \dfrac{h}{\lambda}$.
\[ p = \frac{\SI{6,63e-34}{\joule\second}}{\SI{0,050e-9}{\meter}} = \SI{1,33e-23}{\kilogram\meter\per\second}. \]
\item \textbf{Relation de Compton :} $\lambda' - \lambda = \Lambda_C(1-\cos\theta)$ avec $\Lambda_C = h/(m_e c)$. Lors de la collision photon-électron, il y a conservation de l'\textbf{énergie totale} et de la \textbf{quantité de mouvement totale} (vectorielle) du système \{photon + électron\}.
\item Pour $\theta = 120\degree$, $\cos\theta = -\num{0,5}$, donc $1 - \cos\theta = \num{1,5}$ :
\[ \Delta\lambda = \num{1,5}\times\SI{2,43e-12}{\meter} = \SI{3,6e-12}{\meter} = \SI{0,0036}{\nano\meter}. \]
\[ \lambda' = \lambda + \Delta\lambda = \SI{0,0500}{\nano\meter} + \SI{0,0036}{\nano\meter} = \SI{0,0536}{\nano\meter}. \]
\item Variation relative :
\[ \frac{\Delta\lambda}{\lambda} = \frac{\SI{0,0036}{\nano\meter}}{\SI{0,0500}{\nano\meter}} = \num{7,3e-2} = \SI{7,3}{\percent}. \]
\textit{Commentaire :} la variation est importante et aisément mesurable : dans le domaine des rayons X, l'effet Compton est un phénomène majeur (contrairement au domaine visible, cf. exercice 8).
\item Énergies des photons (on utilise $hc \approx \SI{1240}{\electronvolt\nano\meter}$) :
\[ E = \frac{hc}{\lambda} = \frac{\SI{1240}{\electronvolt\nano\meter}}{\SI{0,0500}{\nano\meter}} = \num{2,48e4}~\text{eV} = \SI{24,8}{\kilo\electronvolt} = \SI{3,97e-15}{\joule}. \]
\[ E' = \frac{hc}{\lambda'} = \frac{\SI{1240}{\electronvolt\nano\meter}}{\SI{0,0536}{\nano\meter}} = \num{2,31e4}~\text{eV} = \SI{23,1}{\kilo\electronvolt} = \SI{3,70e-15}{\joule}. \]
\item Conservation de l'énergie (l'électron est initialement au repos, son énergie cinétique initiale est nulle) :
\[ E = E' + E_c \quad\Rightarrow\quad E_c = E - E' = \SI{24,8}{\kilo\electronvolt} - \SI{23,1}{\kilo\electronvolt} = \SI{1,7}{\kilo\electronvolt}. \]
L'électron de recul emporte donc environ \SI{1,7}{\kilo\electronvolt}.
\item Lors d'une radiographie, une fraction importante des photons X subit la diffusion Compton dans les tissus du patient : ces photons diffusés partent dans toutes les directions. Ceux qui atteignent le capteur n'ont plus de lien géométrique avec la trajectoire directe : ils créent un \textbf{brouillard} qui dégrade le contraste de l'image. Ceux qui se propagent dans la salle irradient le personnel soignant, de façon répétée au fil des examens. Le plomb des tabliers, de numéro atomique $Z$ élevé, absorbe très efficacement les photons X (essentiellement par effet photoélectrique, dont l'efficacité croît fortement avec $Z$) : il protège ainsi les organes du soignant.
\end{enumerate}
\end{corrige}

\begin{attention}[Barème indicatif et points de vigilance — Exercice 10 (sur 10 pts)]
\textbf{Barème :} 1) 1 pt ; 2) 1,5 pt ; 3) 1,5 pt ; 4) 1 pt ; 5) 2 pts ; 6) 1,5 pt ; 7) 1,5 pt.
\textbf{Points de vigilance :} ne pas oublier que $\Delta\lambda$ ne dépend \emph{que} de l'angle $\theta$, jamais de $\lambda$ ; vérifier le signe : le photon diffusé est toujours moins énergétique, donc $\lambda' > \lambda$ ; dans la question 6, exiger la mention explicite de la conservation de l'énergie ($E = E' + E_c$) ; garder au moins trois chiffres significatifs dans $\lambda'$ car $\Delta\lambda$ est petit devant $\lambda$ (risque d'erreur d'arrondi sur $E'$) ; pour la question 7, le correcteur attend le lien « photons diffusés dans toutes les directions $\Rightarrow$ irradiation du personnel $\Rightarrow$ blindage au plomb ».
\end{attention}

\begin{corrige}[Exercice 11 — Type Bac : synthèse, dualité de la lumière]
\textbf{Partie A — Effet photoélectrique.}
\begin{enumerate}
\item Énergie du photon incident :
\[ h\nu = \SI{6,63e-34}{\joule\second}\times\SI{8,50e14}{\hertz} = \SI{5,64e-19}{\joule} = \frac{\SI{5,64e-19}{\joule}}{\SI{1,60e-19}{\joule\per\electronvolt}} = \SI{3,52}{\electronvolt}. \]
Fréquence seuil du sodium : $\nu_0 = \dfrac{W_0}{h} = \dfrac{\SI{2,28}{\electronvolt}\times\SI{1,60e-19}{\joule\per\electronvolt}}{\SI{6,63e-34}{\joule\second}} = \SI{5,50e14}{\hertz}$.
Comme $\nu = \SI{8,50e14}{\hertz} > \nu_0$, \textbf{l'effet photoélectrique a lieu} et :
\[ E_{c,\max} = h\nu - W_0 = \SI{3,52}{\electronvolt} - \SI{2,28}{\electronvolt} = \SI{1,24}{\electronvolt}. \]
\item Potentiel d'arrêt : $U_0 = \dfrac{E_{c,\max}}{e} = \SI{1,24}{\volt}$.
\item Allure des caractéristiques : pour chaque intensité, $I = 0$ tant que $U < -U_0$ (tous les électrons sont repoussés) ; $I$ croît pour $-U_0 < U < 0$ puis atteint un palier $I_{\text{sat}}$ dès que $U$ est suffisamment positif (tous les électrons émis sont collectés). \textbf{Points particuliers :} les deux courbes coupent l'axe des tensions au \emph{même} point $U = -U_0 = \SI{-1,24}{\volt}$ (le potentiel d'arrêt ne dépend que de $\nu$) ; en revanche le palier est d'autant plus haut que l'intensité est grande : $I_{\text{sat},2} > I_{\text{sat},1}$ car $E_2 > E_1$ (le courant de saturation est proportionnel au flux de photons).
\begin{popfigure}
\begin{tikzpicture}
\begin{axis}[
width=\linewidth, height=5.5cm,
axis lines=middle,
xlabel=$U$ (V), ylabel=$I$,
xmin=-2, xmax=1.2, ymin=-0.3, ymax=5,
xtick={-1.24,0}, xticklabels={$-U_0$,$0$},
ytick={2,4}, yticklabels={$I_{\text{sat},1}$,$I_{\text{sat},2}$},
]
\addplot[popcurve, thick, domain=-2:-1.24, samples=10] {0};
\addplot[popcurve, thick, domain=-1.24:0, samples=60] {2*(1+x/1.24)^1.5};
\addplot[popcurve, thick, domain=0:1.2, samples=10] {2};
\addplot[poporange, thick, domain=-2:-1.24, samples=10] {0};
\addplot[poporange, thick, domain=-1.24:0, samples=60] {4*(1+x/1.24)^1.5};
\addplot[poporange, thick, domain=0:1.2, samples=10] {4};
\draw[dashed, thin, popmuted] (axis cs:-1.24,0) -- (axis cs:-1.24,4.6);
\end{axis}
\end{tikzpicture}
\legende{Caractéristiques $I = f(U)$ pour deux intensités $E_1 < E_2$ à fréquence constante : même potentiel d'arrêt, courants de saturation différents.}
\end{popfigure}
\end{enumerate}
\textbf{Partie B — Effet Compton.}
\begin{enumerate}
\setcounter{enumi}{3}
\item Pour $\theta = 90\degree$, $\cos\theta = 0$ :
\[ \Delta\lambda = \Lambda_C = \SI{2,43e-12}{\meter} = \SI{0,00243}{\nano\meter}, \qquad \lambda' = \SI{0,0708}{\nano\meter} + \SI{0,00243}{\nano\meter} = \SI{0,0732}{\nano\meter}. \]
Énergies : $E = \dfrac{hc}{\lambda} = \dfrac{\SI{1240}{\electronvolt\nano\meter}}{\SI{0,0708}{\nano\meter}} = \SI{17,5}{\kilo\electronvolt}$ et $E' = \dfrac{hc}{\lambda'} = \dfrac{\SI{1240}{\electronvolt\nano\meter}}{\SI{0,0732}{\nano\meter}} = \SI{16,9}{\kilo\electronvolt}$.
Énergie cinétique de l'électron de recul (conservation de l'énergie) :
\[ E_c = E - E' = \SI{17,5}{\kilo\electronvolt} - \SI{16,9}{\kilo\electronvolt} = \SI{0,57}{\kilo\electronvolt} \approx \SI{570}{\electronvolt}. \]
\end{enumerate}
\textbf{Partie C — Comparaison.}
\begin{enumerate}
\setcounter{enumi}{4}
\item Tableau complété :
\begin{center}
\begin{tabularx}{\linewidth}{|l|X|X|}
\hline
\rowcolor{popblueL} \textbf{Critère} & \textbf{Effet photoélectrique} & \textbf{Effet Compton} \\
\hline Cible du photon & Électron lié dans un métal & Électron quasi libre (faiblement lié) \\
\hline Destin du photon & Photon totalement absorbé (il disparaît) & Photon diffusé : il change de direction et perd de l'énergie ($\lambda' > \lambda$) \\
\hline Grandeurs conservées & Énergie (la quantité de mouvement est absorbée par le réseau cristallin tout entier) & Énergie totale et quantité de mouvement totale (collision élastique à deux corps) \\
\hline Domaine spectral typique & Visible et ultraviolet proche & Rayons X et $\gamma$ \\
\hline
\end{tabularx}
\end{center}
\item \textbf{Conclusion :} l'effet photoélectrique (absorption d'un quantum d'énergie $h\nu$) et l'effet Compton (collision avec conservation de l'énergie et de la quantité de mouvement $p = h/\lambda$) sont tous deux inexplicables par le modèle ondulatoire et imposent une description \textbf{corpusculaire} de la lumière. Pourtant, les phénomènes d'interférences et de diffraction exigent une description \textbf{ondulatoire} : la lumière possède donc une double nature, c'est la \textbf{dualité onde-corpuscule}.
\textbf{Applications :} \emph{effet photoélectrique} — panneaux photovoltaïques pour l'électrification rurale et le pompage solaire de l'eau dans les villages du Nord-Cameroun, cellules de commande automatique de l'éclairage public ; \emph{effet Compton} — imagerie médicale par rayons X dans les hôpitaux de Douala et Yaoundé (et protection associée du personnel), radiothérapie des cancers.
\end{enumerate}
\end{corrige}

\begin{attention}[Barème indicatif et points de vigilance — Exercice 11 (sur 12 pts)]
\textbf{Barème :} Partie A : 1) 2 pts ; 2) 1 pt ; 3) 2 pts. Partie B : 4) 2,5 pts. Partie C : 5) 2 pts ; 6) 2,5 pts.
\textbf{Points de vigilance :} question 1, la vérification $\nu > \nu_0$ (ou $h\nu > W_0$) doit être écrite \emph{avant} le calcul de $E_{c,\max}$ ; question 3, les deux éléments attendus sont « même $U_0$ » et « $I_{\text{sat}}$ croît avec l'intensité » — un schéma sans commentaire ne suffit pas ; question 4, bien écrire la conservation de l'énergie $E = E' + E_c$ ; question 6, la conclusion doit articuler les deux aspects (succès du corpuscule, nécessité persistante de l'onde) et citer des applications concrètes et plausibles.
\end{attention}

\newpage

\coursec{Fiche bilan}

\begin{popfigure}
\centering
\begin{tikzpicture}[mindmap, grow cyclic, every node/.style={concept, text=white, font=\small},
  root concept/.append style={concept color=popdark, font=\small\bfseries, minimum size=3.2cm},
  level 1/.append style={level distance=3.6cm, sibling angle=60},
  level 1 concept/.append style={minimum size=2.5cm, font=\scriptsize\bfseries}]
\node[root concept]{Effet photo\-électrique et effet Compton}
  child[concept color=popblue] { node {Mise en évidence (Hertz, Hallwachs)}
  }
  child[concept color=popgreen] { node {Seuil et Einstein}
  }
  child[concept color=poporange] { node {Cellule photo\-électrique $I=f(U)$}
  }
  child[concept color=popgold] { node {Applications}
  }
  child[concept color=poppurple] { node {Quantité de mouvement du photon}
  }
  child[concept color=poppink] { node {Effet Compton}
  };
\node[anchor=west, font=\scriptsize, text=popblue] at (4.6,2.4) {$\nu \geq \nu_0$ : émission};
\node[anchor=west, font=\scriptsize, text=popblue] at (4.6,2.0) {instantanée d'électrons};
\node[anchor=west, font=\scriptsize, text=popgreen] at (4.6,0.9) {$W_0 = h\nu_0 = \dfrac{hc}{\lambda_0}$};
\node[anchor=west, font=\scriptsize, text=popgreen] at (4.6,0.5) {$E_{c\max} = h\nu - W_0$};
\node[anchor=west, font=\scriptsize, text=poporange] at (4.6,-0.6) {$I_{\text{sat}}$, potentiel d'arrêt};
\node[anchor=west, font=\scriptsize, text=poporange] at (4.6,-1.0) {$eU_0 = E_{c\max}$};
\node[anchor=east, font=\scriptsize, text=popgold] at (-4.6,1.6) {photodiodes, photomètres,};
\node[anchor=east, font=\scriptsize, text=popgold] at (-4.6,1.2) {panneaux solaires};
\node[anchor=east, font=\scriptsize, text=poppurple] at (-4.6,0.2) {$p = \dfrac{h\nu}{c} = \dfrac{h}{\lambda}$};
\node[anchor=east, font=\scriptsize, text=poppink] at (-4.6,-0.8) {$\lambda' - \lambda = \dfrac{h}{m_e c}(1-\cos\theta)$};
\node[anchor=east, font=\scriptsize, text=poppink] at (-4.6,-1.2) {$\Lambda_C = \num{2,43e-12}\ \text{m}$};
\end{tikzpicture}
\legende{Carte mentale de la leçon : les deux preuves de la granularité de la lumière.}
\end{popfigure}

\begin{bilan}
\textbf{Définitions clés}
\begin{itemize}
  \item \cle{Effet photoélectrique} : émission d'électrons par un métal éclairé par une radiation de fréquence suffisante. Mise en évidence : expériences de \cle{Hertz} (1887) et de \cle{Hallwachs} (électroscope chargé négativement se décharge sous UV).
  \item \cle{Seuil photoélectrique} : fréquence minimale $\nu_0$ (longueur d'onde maximale $\lambda_0$) en dessous de laquelle aucun électron n'est extrait, quelle que soit l'intensité lumineuse. $\nu_0$ dépend du métal.
  \item \cle{Travail d'extraction} $W_0$ : énergie minimale pour arracher un électron au métal.
  \item \cle{Potentiel d'arrêt} $U_0$ : tension inverse qu'il faut appliquer pour annuler le courant photoélectrique.
  \item \cle{Effet Compton} : diffusion d'un photon X par un électron quasi libre, avec augmentation de la longueur d'onde du photon diffusé.
\end{itemize}

\textbf{Formules à connaître par cœur}
\begin{center}
\begin{tabularx}{\linewidth}{|l|X|l|}
\hline
\rowcolor{popblueL} \textbf{Grandeur} & \textbf{Relation} & \textbf{Unités / remarques} \\
\hline
Énergie du photon & $E = h\nu = \dfrac{hc}{\lambda}$ & joules (J) ; $h = \num{6,626e-34}\ \text{J.s}$ \\
\hline
Seuil photoélectrique & $W_0 = h\nu_0 = \dfrac{hc}{\lambda_0}$ & $\lambda_0 = \dfrac{c}{\nu_0}$ \\
\hline
Relation d'Einstein & $E_{c\max} = h\nu - W_0 = h(\nu - \nu_0)$ & valable si $\nu \geq \nu_0$ \\
\hline
Potentiel d'arrêt & $eU_0 = E_{c\max}$ & $U_0$ en volts, $e = \num{1,60e-19}\ \text{C}$ \\
\hline
Courant de saturation & $I_{\text{sat}} = n\,e$ & $n$ : nombre d'électrons émis par seconde \\
\hline
Quantité de mouvement du photon & $p = \dfrac{h\nu}{c} = \dfrac{h}{\lambda}$ & en $\text{kg.m.s}^{-1}$ \\
\hline
Effet Compton & $\lambda' - \lambda = \dfrac{h}{m_e c}(1 - \cos\theta)$ & $\Lambda_C = \dfrac{h}{m_e c} = \num{2,43e-12}\ \text{m}$ \\
\hline
\end{tabularx}
\end{center}

\textbf{Méthodes express}
\begin{itemize}
  \item \cle{Exploiter $I = f(U)$} : lire $I_{\text{sat}}$ (palier) ; lire $U_0$ (abscisse où $I = 0$) ; en déduire $E_{c\max} = eU_0$, puis $\nu$ ou $W_0$ via Einstein.
  \item \cle{Conversion rapide} : $1\ \text{eV} = \num{1,60e-19}\ \text{J}$ ; pour $\lambda$ en nm : $E\ (\text{eV}) \approx \dfrac{1240}{\lambda\ (\text{nm})}$.
  \item \cle{Compton} : $\Delta\lambda$ maximal pour $\theta = 180^{\circ}$ (rétrodiffusion) : $\Delta\lambda_{\max} = 2\Lambda_C$ ; nul pour $\theta = 0^{\circ}$.
  \item \cle{Bilan de puissance} : puissance lumineuse $P = N h\nu$ ($N$ photons/s) ; rendement quantique $\eta = \dfrac{\text{électrons émis}}{\text{photons incidents}}$.
\end{itemize}
\end{bilan}

\begin{aretenir}[Checklist avant le Bac]
\begin{itemize}
  \item[$\square$] Je sais décrire l'expérience de Hertz / de l'électroscope et conclure : l'émission exige $\nu \geq \nu_0$.
  \item[$\square$] Je sais expliquer pourquoi le modèle ondulatoire échoue (seuil, instantanéité, $E_{c\max}$ indépendante de l'intensité).
  \item[$\square$] Je connais $E = h\nu = hc/\lambda$ et je convertis J $\leftrightarrow$ eV sans erreur.
  \item[$\square$] J'écris la relation d'Einstein $E_{c\max} = h\nu - W_0$ avec la condition $\nu \geq \nu_0$.
  \item[$\square$] Je calcule $\lambda_0 = hc/W_0$ et je vérifie l'homogénéité.
  \item[$\square$] Je trace et j'exploite $I = f(U)$ : $I_{\text{sat}}$, $U_0$, et $eU_0 = E_{c\max}$.
  \item[$\square$] Je sais que $I_{\text{sat}}$ croît avec l'intensité lumineuse (à $\nu$ fixée) mais pas $E_{c\max}$.
  \item[$\square$] Je connais la quantité de mouvement du photon $p = h/\lambda$.
  \item[$\square$] Je sais appliquer $\lambda' - \lambda = \dfrac{h}{m_e c}(1-\cos\theta)$ avec $\theta$ en degrés et les bonnes unités.
  \item[$\square$] Je sais que l'effet Compton confirme la nature corpusculaire de la lumière (choc photon--électron, conservation de l'énergie et de la quantité de mouvement).
  \item[$\square$] Je cite des applications : cellules photoélectriques, photodiodes, panneaux photovoltaïques (éclairage public solaire à Maroua, électrification rurale).
  \item[$\square$] Je soigne les chiffres significatifs et les unités SI dans chaque résultat numérique.
\end{itemize}
\end{aretenir}

\findecours

\end{document}
